% MDPinn working paper / lab notebook
%
% Personal tracking document, not for distribution yet. Structured like a
% paper draft so it's easy to grow into one, but treat every section as
% provisional - fill in citations, tighten claims, and prune hedges as you
% do your own reading. TODO markers throughout mark places that need more
% than what this session's work alone can responsibly claim.
%
% Compiles standalone with pdflatex (no bibtex/biber pass needed - the
% bibliography is a plain thebibliography block). If this grows large enough
% to want proper .bib management, that's an easy later migration.

\documentclass[11pt]{article}

\usepackage[margin=1in]{geometry}
\usepackage{amsmath, amssymb}
\usepackage{booktabs}
\usepackage{graphicx}
\usepackage{hyperref}
\usepackage{xcolor}
% natbib required for \citet/\citep (already used elsewhere in this file,
% e.g. \citet{cheng2019water}) and for \bibliographystyle{plainnat} itself -
% missing until now, which would have made this fail to compile the first
% time anyone actually ran pdflatex on it (never yet done locally - no LaTeX
% install in this checkout, see the file-level comment above).
\usepackage{natbib}

\newcommand{\todo}[1]{\textcolor{red}{\textbf{[TODO: #1]}}}

\title{Does Per-Molecule Momentum Conservation Improve Neural Network
Potentials? \\ \large Working notes and results log}
\author{Dietrich Lachmann}
\date{\today}

\begin{document}
\maketitle

\begin{abstract}
This is a working document, not a finished paper. It tracks the motivation,
methods, and results of an ongoing investigation into whether adding a
physics-informed momentum-conservation loss term improves neural network
interatomic potentials (NNIPs), and specifically whether that effect depends
on the system being a single isolated molecule versus a periodic
multi-molecule system where momentum conservation is not already guaranteed
by architectural equivariance. On MD17 aspirin (single molecule), the
constraint is confirmed architecturally redundant, motivating a move to a
periodic liquid-water benchmark (Cheng et al. 2019, 64 molecules/box), where
per-molecule momentum conservation is confirmed genuinely non-trivial.
TensorNet's own periodic neighbor search is confirmed correctly
periodicity-aware (Section~\ref{sec:periodicity}). The static-metric
momentum-vs-absolute comparison that looked like a clean win at $n=3$ seeds
does \emph{not} survive to $n=6$: it turns out to depend heavily on
evaluation methodology and on one noisy checkpoint, and the only signal that
holds up under every lens tried is a modest force-quality improvement
(Section~\ref{sec:n6-update}). A first real-dynamics evaluation (NVE
rollouts via ASE, Section~\ref{sec:rollout-results}) finds both conditions
produce unstable dynamics from a realistic starting configuration --- a
failure mode invisible to every static metric above --- with the momentum
condition consistently running hotter, replicated across 10 independent
trials on two different replicate axes. Digging into why (Section~\ref{sec:q3-progress}--\ref{sec:q3-seed5})
found a systematic, statistically robust bias toward higher net per-molecule
force in the specific momentum-trained checkpoint used for these rollouts
--- but that bias was already visible in that checkpoint's own static
evaluation before any dynamics were run, and substantially (though not
completely) weakens at a training seed whose static evaluation favored
momentum instead. Seed-to-seed variation in whether the momentum loss
achieves its intended effect, more than any fixed property of the method
itself, is the best-supported explanation so far. Separately, both
conditions turn out to be dynamically unstable regardless of the momentum
question (Section~\ref{sec:q4}) --- a real short-range collapse mechanism
contributes in some but not all trajectories, and the literature-motivated
fix (an explicit ZBL repulsive prior) was tested directly and found to make
stability substantially \emph{worse}, not better, replicated across two
independent training seeds (Section~\ref{sec:q4-negative-result}) --- a
negative result attributed to the prior's cutoff necessarily overlapping the
O--H covalent bond length in this specific hydrogen-bonded system.
Section~\ref{sec:results} reports everything in full; Section~\ref{sec:next}
lists what's still open: whether the (now much weaker) static momentum
effect is real at all, and whether the residual rollout-stability lean at a
``good'' momentum seed is worth chasing further. \todo{replace this abstract
once both have answers.}
\end{abstract}

\tableofcontents

\section{Motivation}
\label{sec:motivation}

Neural network interatomic potentials (NNIPs) built on equivariant
architectures --- e.g.\ TensorNet \cite{tensornet}, as implemented in
TorchMD-Net \cite{torchmdnet2} --- are trained to reproduce reference
energies and forces from a supervised dataset. A natural question for
physics-informed learning is whether adding an explicit soft constraint for
a conservation law (here: momentum conservation) as an auxiliary loss term
improves the model beyond what supervised training on energies/forces alone
achieves.

The naive version of this question turns out to be uninteresting for a
single isolated molecule: an $O(3)$-equivariant architecture already
guarantees that the \emph{net} force on the whole molecule is (numerically)
zero, for any input, before any training pressure is applied to it at all.
We confirmed this empirically before investing further effort in it (see
Section~\ref{sec:aspirin}) rather than assuming it from the equivariance
argument alone.

The more interesting version of the question is whether momentum
conservation checked \emph{per individual molecule within a periodic,
multi-molecule system} is similarly redundant, or whether it is a genuine,
non-trivial constraint the architecture does not already enforce. Physically,
there is no reason to expect it to be redundant: while the whole periodic
box's net momentum is guaranteed zero by the same equivariance argument, an
individual water molecule being pushed on by its hydrogen-bonded neighbors
has no such guarantee. This is the premise of the water-box study
(Section~\ref{sec:waterbox}).

\section{Background}
\label{sec:background}

\subsection{Architecture and datasets}

All models here are TensorNet \cite{tensornet} potentials trained via
TorchMD-Net 2.0's \texttt{LNNP} training module \cite{torchmdnet2}. Two
datasets are in use:

\begin{itemize}
  \item \textbf{MD17} \cite{md17} --- single-molecule ab initio MD
  trajectories (aspirin, benzene, ethanol used so far). Native frame spacing
  is 0.5\,fs \cite{md17}, not 0.1\,fs as an earlier version of this codebase
  assumed --- see the implementation notes in Appendix~\ref{app:impl}.
  \item \textbf{WaterBox} \cite{cheng2019water} --- 1593 periodic liquid-water
  configurations (64 molecules / 192 atoms each), revPBE0-D3 DFT, from
  \citet{cheng2019water}, distributed via Materials Cloud
  (DOI: 10.24435/materialscloud:2018.0020/v1) and exposed as
  \texttt{torchmdnet.datasets.WaterBox}. \todo{this is the same dataset MACE
  \cite{mace} benchmarks against --- worth reading the rest of that paper's
  water section more carefully, not just the headline numbers pulled in
  Section~\ref{sec:waterbox-results}.}
\end{itemize}

\subsection{Loss formulation}
\label{sec:loss-formulation}

Both studies use the same base supervised (energy + force) loss plus an
optional momentum term; they differ only in which momentum-loss variant is
used. Equations below are pulled directly from this codebase
(\texttt{physics\_losses.py}) and \texttt{torchmdnet.loss}, not re-derived
from memory.

\paragraph{Base supervised loss.} For a training batch of $B$ configurations,
each with $N$ atoms, energy predictions $\hat E_b$ and force predictions
$\hat F_{b,i} \in \mathbb{R}^3$:
\begin{align}
\mathcal{L}_y &= \frac{1}{B}\sum_{b=1}^{B} \left(\hat E_b - E_b\right)^2
&&\text{(energy MSE, \texttt{torchmdnet.loss.mse\_loss})} \\
\mathcal{L}_F &= \frac{1}{3BN}\sum_{b=1}^{B}\sum_{i=1}^{N}\sum_{\alpha \in \{x,y,z\}}
\left(\hat F_{b,i,\alpha} - F_{b,i,\alpha}\right)^2
&&\text{(force MSE, same function, flattened over atoms \& components)} \\
\mathcal{L}_{\mathrm{sup}} &= w_E\, \mathcal{L}_y + w_F\, \mathcal{L}_F
\end{align}
with $(w_E, w_F) = (0.05, 0.95)$ throughout the aspirin study and the
water-box study's original 20-epoch runs (Section~\ref{sec:waterbox-results}).
The water-box study's 50-epoch retrain instead anneals $(w_E, w_F)$ partway
through training (Section~\ref{sec:anneal}).

\paragraph{Whole-molecule momentum loss} (\texttt{momentum\_symmetry\_loss} ---
aspirin study only). For one $N$-atom system with positions $r_i$ and
predicted forces $\hat F_i$, centroid $\bar r = \frac{1}{N}\sum_i r_i$:
\begin{equation}
\mathcal{L}_{\mathrm{mom}} = \frac{1}{N}\left[
\underbrace{\left\lVert \textstyle\sum_{i=1}^N \hat F_i \right\rVert^2}_{\text{net force (linear momentum)}}
\;+\;
\underbrace{\left\lVert \textstyle\sum_{i=1}^N (r_i - \bar r) \times \hat F_i \right\rVert^2}_{\text{net torque about centroid (angular momentum)}}
\right]
\label{eq:mom-whole}
\end{equation}
Both terms are normalized once by the same $N$ (whole-system atom count).

\paragraph{Per-fragment momentum loss} (\texttt{per\_fragment\_momentum\_loss}
--- water-box study only). Partition the $N$ atoms in a configuration into
$M$ molecules, fragment $m$ having atom set $\mathcal{A}_m$ and
$n_m = |\mathcal{A}_m|$ atoms ($n_m = 3$ here: one O, two H). For each
fragment, with its own centroid $\bar r_m = \frac{1}{n_m}\sum_{i\in\mathcal{A}_m} r_i$:
\begin{equation}
\mathcal{L}_{\mathrm{mom}}^{(m)} = \frac{1}{n_m}\left[
\left\lVert \textstyle\sum_{i\in\mathcal{A}_m} \hat F_i\right\rVert^2
+
\left\lVert \textstyle\sum_{i\in\mathcal{A}_m}(r_i - \bar r_m)\times \hat F_i\right\rVert^2
\right],
\qquad
\mathcal{L}_{\mathrm{mom\text{-}frag}} = \frac{1}{M}\sum_{m=1}^{M} \mathcal{L}_{\mathrm{mom}}^{(m)}
\label{eq:mom-frag}
\end{equation}
Structurally this is the \emph{same} linear+angular penalty as
Eq.~\ref{eq:mom-whole}, applied per molecule instead of per whole system,
then \emph{averaged} across the $M=64$ molecules --- not summed, and not
normalized once by total atom count $N=3M$. Each fragment's own term is
normalized by its own $n_m$ before averaging, so one fragment's contribution
does not get diluted just because other fragments in the same configuration
exist; this is what makes it sensitive to a single badly-behaved molecule
rather than only the system-wide average. Computed via \texttt{index\_add\_}
(scatter) over all fragments in a batch at once, not a Python loop per
fragment --- see the vectorization lesson in
Appendix~\ref{app:impl}.

\paragraph{Total training loss.}
\begin{equation}
\mathcal{L}_{\mathrm{total}} = \mathcal{L}_{\mathrm{sup}} + \lambda_{\mathrm{mom}}\, \mathcal{L}_{\mathrm{mom}}^{(\cdot)}
\end{equation}
with $\lambda_{\mathrm{mom}} = \texttt{momentum\_weight}$ and
$\mathcal{L}_{\mathrm{mom}}^{(\cdot)} \in \{\mathcal{L}_{\mathrm{mom}},\,
\mathcal{L}_{\mathrm{mom\text{-}frag}}\}$ depending on the study:
\begin{itemize}
  \item \textbf{Aspirin}: \texttt{absolute} ($\lambda_{\mathrm{mom}}=0$) vs.\
  \texttt{absolute+momentum} ($\lambda_{\mathrm{mom}}=0.01$), using
  $\mathcal{L}_{\mathrm{mom}}$ (Eq.~\ref{eq:mom-whole}).
  \item \textbf{Water-box}: \texttt{water\_absolute} ($\lambda_{\mathrm{mom}}=0$)
  vs.\ \texttt{water\_absolute+momentum} ($\lambda_{\mathrm{mom}}=0.01$), using
  $\mathcal{L}_{\mathrm{mom\text{-}frag}}$ (Eq.~\ref{eq:mom-frag}), computed
  via bond inference plus connected components to group atoms into molecules
  (order-independent --- does not assume atoms are stored in contiguous
  per-molecule blocks).
\end{itemize}
When $\lambda_{\mathrm{mom}}=0$, the momentum term is skipped entirely
(\texttt{WaterLNNP.step}) rather than computed and multiplied by zero, so the
\texttt{absolute}/\texttt{water\_absolute} conditions pay no extra
forward-pass cost --- unlike \texttt{train\_physics.py}'s
\texttt{PhysicsInformedLNNP}, which pays for the extra pass unconditionally
(\todo{worth reconciling why these two paths diverged - probably just
written at different times, check if worth unifying}).

\subsection{Anneal-invariant checkpoint selection}
\label{sec:checkpoint-fix}

\texttt{ModelCheckpoint}/\texttt{EarlyStopping} originally monitored
\texttt{val\_total\_mse\_loss} = $w_E \cdot \texttt{val\_y\_mse} + w_F \cdot
\texttt{val\_neg\_dy\_mse}$, using whichever $(w_E, w_F)$ are active that
epoch. Once the anneal schedule above changes those weights partway through
a run, this metric's own scale shifts at exactly the epoch boundary it is
supposed to help evaluate: pre-anneal totals are dominated by the small,
well-behaved force term (95\% weight); post-anneal totals are dominated by
the large, volatile energy term (75\% weight). A post-anneal epoch is
therefore judged by a structurally harsher yardstick than a pre-anneal one,
regardless of whether the model actually improved - confirmed empirically
before this was understood as a bug: every one of six water-box runs picked
its ``best'' checkpoint from before the anneal switch.

Fixed via \texttt{val\_checkpoint\_score}
(\texttt{WaterLNNP.on\_validation\_epoch\_end}, calling \texttt{super()}
first): the same weighted-sum formula, but always using the run's
\emph{original} (pre-anneal) weights regardless of which phase is currently
active, so ``best'' means the same thing on every epoch of the run.
\texttt{ModelCheckpoint} and \texttt{EarlyStopping} now monitor this instead.
For a non-annealed run this is identical to the old metric (a no-op change).
Also added: a second, unconditional \texttt{ModelCheckpoint}
(\texttt{monitor=None}, \texttt{save\_top\_k=0}, \texttt{save\_last=True}) to
get a genuine always-current-epoch checkpoint - the original single
monitored callback's \texttt{save\_last=True} only updates ``whenever a
checkpoint file gets saved'' (Lightning's own docs), which under
\texttt{save\_top\_k=1} only happens on improvement, so it silently froze at
the same epoch as ``best'' whenever nothing later ever beat it - not
actually the final epoch, discovered only after a wasted diagnostic
round-trip assuming otherwise.

\section{Methods}
\label{sec:methods}

\subsection{Aspirin ablation}
\label{sec:aspirin}

Two conditions (\texttt{absolute}, \texttt{absolute+momentum}) on MD17
aspirin, 3 seeds each, 20 epochs, TensorNet as above. \texttt{delta} /
\texttt{delta+momentum} (residual learning against an analytic CHARMM36
baseline) were scoped out after the analytic baseline's unvectorized
$O(N^2)$ implementation made a single epoch take $\sim$50 hours --- not
revisited without vectorizing that code path first.

\subsection{Water-box study}
\label{sec:waterbox}

Two conditions (\texttt{water\_absolute}, \texttt{water\_absolute+momentum},
momentum weight 0.01), 3 seeds each, 20 epochs, same architecture and
optimizer settings as the aspirin study \emph{except} batch size and
gradient-accumulation, which were re-derived specifically for this system
(see Appendix~\ref{app:impl} --- the aspirin batch size does not fit a
192-atom periodic system with TensorNet's per-edge tensor representation).

The evaluation metrics are: energy MAE and force MAE against DFT ground
truth (standard NNIP accuracy metrics), and mean/max per-molecule momentum
violation on held-out configurations (the metric this study actually cares
about --- see Section~\ref{sec:motivation}).

\subsection{Energy/force loss-weight annealing}
\label{sec:anneal}

Motivated by a pattern in the water-box study's own 20-epoch results
(Section~\ref{sec:waterbox-results}, Section~\ref{sec:caveats}): energy
oscillated by orders of magnitude across every seed of both conditions,
while force stayed comparatively well-behaved, under a \emph{fixed}
0.05/0.95 (energy/force) weighting held constant for the entire run. A
plausible contributor is that force receives 19$\times$ the gradient signal
energy does, for the whole run, with nothing ever correcting for it.

Current practice in two of the most directly comparable architectures
addresses exactly this by \emph{annealing} the energy/force weighting during
training rather than holding it fixed: force-dominant early (to fix the
local force/gradient shape), then switched to energy-dominant later
specifically to fix up absolute energy calibration.
\begin{itemize}
  \item MACE \cite{mace}: one reported schedule uses energy:force $=1{:}100$
  for roughly the first 75\% of training, then $300{:}100$ ($3{:}1$
  favoring energy) for the remainder; a related MACE recipe reports
  force-dominant weighting for $\sim$60\% of training before switching to
  energy-dominant for the rest.
  \item NequIP \cite{nequip}: energy:force $=40{:}1000$ early, flipped at a
  fixed epoch.
\end{itemize}
Not universal practice --- Allegro reportedly uses equal, unannealed
energy/force weighting throughout --- and more generally, fixed
loss-weight coefficients (annealed or not) are a known weak point: optimal
weights are dataset- and architecture-dependent, motivating adaptive
schemes that adjust weights automatically during training rather than
following a hand-picked fixed schedule \cite{adaptive-weighting}. This
project uses a hand-picked fixed schedule, not an adaptive one --- worth
revisiting \citet{adaptive-weighting}'s approach directly if the fixed
schedule below doesn't resolve the energy-oscillation problem.

\textbf{Schedule used here} (added 2026-08-04, for the 50-epoch retrain):
fixed 0.05/0.95 (energy/force) for epochs 0--29 (60\% of 50 epochs, matching
the MACE timing convention above), then switched to 0.75/0.25 ($3{:}1$
favoring energy, matching the $\sim3{:}1$ ratio of MACE's reported
final-phase weights) for epochs 30--49. A hard switch at a fixed epoch, not
a smooth interpolation, matching what the cited schedules actually report
rather than inventing an untested variant. Implemented as a Lightning hook
(\texttt{WaterLNNP.on\_train\_epoch\_start}) that mutates
\texttt{self.hparams.y\_weight}/\texttt{neg\_dy\_weight} in place --
\texttt{LNNP.step()} reads these fresh every call rather than caching them
at construction, so this takes effect with no changes needed to
\texttt{torchmdnet} itself.

\textbf{Result} (Section~\ref{sec:q1} has the full story, including a
checkpoint-selection bug this schedule exposed and then helped diagnose):
once checkpoint selection was made anneal-invariant
(\texttt{val\_checkpoint\_score}, Section~\ref{sec:checkpoint-fix}), 5 of the
6 water-box cells selected a genuinely post-anneal epoch as best, versus 0 of
6 under the old monitor - direct, mechanism-level confirmation that this
schedule is doing real work, not just correlating with an artifact.

\subsection{Periodic neighbor search verification}
\label{sec:periodicity}

The bond-inference/molecule-grouping step used by the per-fragment momentum
loss (Section~\ref{sec:loss-formulation}) was fixed and confirmed
periodicity-aware early in this study (Appendix~\ref{app:impl}, item 3) --
but that is a separate, simpler computation this project owns
(\texttt{structural\_metrics.py}). TensorNet's own distance/neighbor module,
inside the installed \texttt{torchmdnet} package, had never been directly
checked the same way: \texttt{box} vectors are passed into every model
forward call (\texttt{LNNP.step()} does this automatically, confirmed by
reading \texttt{torchmdnet}'s own source), but accepting a box tensor and
correctly applying minimum-image distance inside the neighbor search are two
different claims, and only the first was previously verified.

\paragraph{Test.} A naive first instinct -- translate every atom in a
configuration by the same lattice vector and wrap the whole system back into
the primary cell -- turns out to test nothing: for atoms already stored
inside $[0, L)$, adding a full lattice vector and wrapping reproduces the
exact same absolute coordinates you started with (up to float32 noise), so
the model sees an unchanged input. The actual non-trivial invariance a
correctly periodic model must satisfy is: translate a \emph{single} molecule
(identified via the already-verified \texttt{infer\_molecule\_groups}) by one
full lattice vector, \emph{without} wrapping it back into the cell and
without touching any other atom, and compare predicted energy/force against
the unshifted baseline. Under correct minimum-image handling this must be
exactly invariant, since the molecule's new (now out-of-cell) coordinate
describes the same physical position under periodic boundary conditions. A
half-lattice-vector shift is run alongside it as a negative control -- a
genuinely different configuration, which should \emph{not} be invariant, to
confirm the test is sensitive rather than trivially passing. Implemented in
\texttt{src/verify\_periodicity.py}; tests an architectural property of the
forward pass, not training quality, so any trained checkpoint works equally
well as the test subject.

\paragraph{Result.} Across 5 held-out configurations, one trained
\texttt{water\_absolute} checkpoint (seed 0): the five real lattice-vector
shifts ($+a, +b, +c, +a{+}b{+}c, -a$) changed predicted energy by at most
0.0059\,eV and force by at most $9{\times}10^{-5}$\,eV/\AA{} (max over all
192 atoms), while the half-lattice-vector negative control changed energy by
2.9--5.25\,eV and force by up to 11.7\,eV/\AA{} -- roughly 500--1500$\times$
larger. \textbf{TensorNet's neighbor search is confirmed to correctly apply
minimum-image convention in this pipeline}; open item 2 (CLAUDE.md) is
closed.

\paragraph{A precision false alarm, resolved.} The residual diffs on the
five real shifts initially looked suspiciously quantized -- every value an
exact multiple of $2^{-9}$\,eV ($\approx 1.953$\,meV) -- raising a concern
that some part of the energy computation was running at reduced precision
(TF32/fp16) despite \texttt{precision=32} being set explicitly throughout
this codebase. Resolved by checking the actual magnitude involved rather
than assuming: with \texttt{remove\_ref\_energy=False}, the model predicts
\emph{raw absolute} total system energy, confirmed $\approx -30{,}000$\,eV
for this 192-atom system -- not the $O(1$--$10)$\,eV scale the energy-MAE
numbers in Table~\ref{tab:waterbox-summary} might suggest (those are errors
\emph{against} ground truth, not the prediction's own magnitude). Plain
float32's relative machine epsilon ($2^{-23}$) at a $-30{,}000$\,eV magnitude
predicts a rounding floor of $\approx 0.0036$\,eV -- matching the observed
diffs (0--0.0059\,eV) closely. This is ordinary float32 behavior at a large
absolute energy scale, not a precision bug; no code or configuration change
was needed. \todo{worth remembering next time a metric looks suspiciously
quantized: check the absolute magnitude of the underlying number and compute
the expected float32 ULP before chasing exotic explanations.}

\subsection{Real-dynamics evaluation: NVE rollouts via ASE}
\label{sec:rollout-methods}

Static energy/force/momentum-violation metrics on held-out configurations
don't establish that a potential produces stable, physically realistic MD
trajectories. Rather than extending \texttt{rollout\_nve.py} (velocity-Verlet,
no periodic-boundary handling, built for aspirin) to the periodic case,
dynamics here are driven by ASE \cite{ase} instead: \texttt{TensorNetCalculator}
(\texttt{src/waterbox\_ase.py}) wraps a trained checkpoint as an ASE
\texttt{Calculator} (reusing \texttt{evaluate\_waterbox.py}'s checkpoint
loader and forward-pass convention exactly, not a new, separately-trusted
code path), and ASE's own \texttt{VelocityVerlet}/\texttt{NeighborList}/
\texttt{get\_rdf} handle the periodic integration and analysis --- already
tested code, rather than a second hand-rolled minimum-image implementation
on top of the one just verified in Section~\ref{sec:periodicity}.

\textbf{Protocol} (\texttt{src/rollout\_waterbox\_ase.py}): start from a real
held-out test configuration's positions (same split/seed as every other
evaluation script), assign Maxwell-Boltzmann velocities at 300\,K, zero net
momentum, integrate plain NVE --- no thermostat, since a thermostat would
correct away exactly the instability this is meant to detect --- with
\texttt{VelocityVerlet} at $\Delta t = 0.5$\,fs. Two diagnostics: total-energy
drift over the trajectory, and O--O/O--H/H--H RDF compared against a
reference computed identically from real DFT configurations (not the
rollout itself). \texttt{rdf\_rmax} is capped below half the smallest box
edge among whichever frames are actually in play (box size varies
configuration to configuration in this dataset, $\sim$9--14\,\AA{} observed)
rather than a fixed guess, since \texttt{get\_rdf}'s minimum-image
convention requires $r_{\max} < L/2$ --- confirmed necessary in practice: a
fixed 6.0\,\AA{} default hit \texttt{CellTooSmall} on the training box.

\textbf{Two independent random-seed roles}: \texttt{seed} controls the
train/val/test split (which physical configuration a given index refers to)
and must stay fixed within any batch of replicates meant to share a starting
geometry; \texttt{velocity\_seed} controls only the initial
Maxwell-Boltzmann draw. An early version conflated the two, which silently
varies the starting geometry between what were meant to be velocity-only
replicates --- caught by a 335\,K vs.\ 292\,K starting-temperature mismatch
between two runs that were supposed to be identical (Section~\ref{sec:rollout-results}).

\subsection{Force-decomposition investigation (Q3 mechanism)}
\label{sec:force-decomposition-methods}

Motivated by the rollout instability (Section~\ref{sec:rollout-results}):
why does per-fragment momentum-conservation training produce \emph{hotter}
NVE rollouts than plain supervised training, the opposite of the naive
expectation? Three tools, in the order they were used:

\textbf{Force decomposition} (\texttt{src/analyze\_force\_decomposition.py}).
For molecule $m$ with atoms $\mathcal{A}_m$: $F_{\mathrm{net},m} =
\sum_{i\in\mathcal{A}_m} \hat F_i$ (exactly the quantity
\texttt{per\_fragment\_momentum\_loss}'s linear term penalizes) and
$F_{\mathrm{internal},i} = \hat F_i - F_{\mathrm{net},m}/|\mathcal{A}_m|$ for
each atom $i\in\mathcal{A}_m$ --- a pure algebraic split (sums to zero across
each molecule's own atoms by construction, not a model property) isolating
the ``shape-distorting'' force component the momentum loss cannot see: a
model can drive $F_{\mathrm{net}}\to 0$ while $F_{\mathrm{internal}}$ stays
large. \texttt{--mode equilibrium} runs this on real held-out DFT
configurations; \texttt{--mode trajectory} runs it on real frames read back
from a saved rollout (\texttt{ase.io.read}), \emph{cross-evaluating} both
models on frames from both conditions' own trajectories rather than each
model only on its own --- separating ``does a model's force response change
on hot geometries'' from ``are the hot geometries one condition's rollout
reaches different from the other's,'' which a same-model-on-own-trajectory
comparison would conflate.

\textbf{Replication across existing rollouts}
(\texttt{src/run\_force\_decomposition\_study.py}). Reuses the 10 trajectory
pairs already collected by \texttt{run\_rollout\_study.py} (5 velocity-draw
+ 5 config replicates, Section~\ref{sec:rollout-results}) --- no new
rollouts needed. For each replicate, cross-evaluates both models at frames
0/20/80/180 (steps 0/200/800/1800) and computes the momentum-minus-absolute
net-force gap at matched elapsed time, averaged over which model did the
evaluating (justified once the two models were shown to agree closely
point-wise --- see Section~\ref{sec:q3-progress}).

\section{Results so far}
\label{sec:results}

\subsection{Aspirin: whole-molecule momentum is redundant}

Confirmed empirically: \texttt{train\_loss\_momentum} sits at
$\sim 10^{-9}$--$10^{-10}$ (floating-point noise) from epoch 0 onward across
all aspirin seeds, regardless of training. Force MAE converges similarly
($\sim$0.8--1.1\,eV/\AA) across conditions; energy MAE is highly
seed-dependent within a single condition (1.6 to 184.5\,eV across seeds).
\texttt{absolute+momentum} costs $\sim$48\% more wall-clock time for
comparable accuracy. Median rollout-drift score nominally favors
\texttt{absolute+momentum} ($\sim$125$\times$ lower) but this is
\emph{not} distinguishable from seed-to-seed calibration luck at $n=3$.
\todo{this ambiguity is exactly what motivates being careful about the
water-box comparison below - don't repeat the same n=3-and-declare-victory
mistake here.}

\subsection{Water-box: per-molecule momentum conservation is genuinely
violated}
\label{sec:waterbox-results}

The core premise holds up clearly: mean per-molecule momentum violation is
$81.41 \pm 64$ (\texttt{water\_absolute}) and $43.09 \pm 4$
(\texttt{water\_absolute+momentum}) --- many orders of magnitude above
aspirin's $\sim 10^{-9}$ floor, confirming this \emph{is} a non-trivial,
architecturally-unenforced constraint in the periodic multi-molecule
setting, unlike the single-molecule case.

\begin{table}[h]
\centering
\caption{Water-box study, $n=3$ seeds per condition (mean $\pm$ std). These
exact numbers (to reported precision) have now held across three separate
runs: the original 20-epoch run, a 50-epoch annealed retrain still using a
since-fixed, biased checkpoint-selection metric, and the same 50-epoch
annealed retrain with that metric corrected (Section~\ref{sec:checkpoint-fix})
--- see Section~\ref{sec:q1} for why that stability is informative rather
than suspicious. Energy/force MAE are raw per-configuration totals (192
atoms / 64 molecules), \emph{not} yet normalized per-atom or per-molecule
--- see the note below the table.}
\label{tab:waterbox-summary}
\begin{tabular}{lcccccc}
\toprule
Condition & Energy MAE (eV) & Force MAE (eV/\AA) & Mean $\lVert p \rVert_{\mathrm{mol}}$ & Max $\lVert p \rVert_{\mathrm{mol}}$ & Wall (s) \\
\midrule
\texttt{water\_absolute}          & $8.92 \pm 6.6$  & $1.256 \pm 1.1$   & $81.4 \pm 64$ & $1859 \pm 1.0{\times}10^3$ & $5132 \pm 8.1$ \\
\texttt{water\_absolute+momentum} & $6.18 \pm 1.3$  & $0.631 \pm 0.039$ & $43.1 \pm 4$  & $1526 \pm 3.7{\times}10^2$ & $1.04{\times}10^4 \pm 42$ \\
\bottomrule
\end{tabular}
\end{table}

Every one of the four accuracy/physicality metrics moves in the same
direction (momentum condition better), and with much tighter spread across
seeds. Taken at face value against the pre-registered ``clears roughly one
std'' heuristic, none of the four individually clears that bar --- e.g.\ the
energy MAE gap (2.7\,eV) is smaller than \texttt{water\_absolute}'s own std
(6.6\,eV). \textbf{This is not yet a confirmed effect}, for a specific,
diagnosed reason (Section~\ref{sec:caveats}), not just ``small $n$.''

The momentum condition's wall-clock cost is $\sim2\times$
\texttt{water\_absolute}'s ($5132\,\mathrm{s} \to 1.04{\times}10^4\,
\mathrm{s}$), consistent across all three seeds to within 0.4\% --- expected
and mechanically confirmed: the momentum loss requires a second full
forward+backward pass per training step.

\subsubsection{External calibration: comparison to MACE on the same dataset}

The WaterBox dataset used here (\citet{cheng2019water}, 1593 configs) is the
\emph{same} benchmark MACE \cite{mace} reports on. MACE achieves energy RMSE
of 1.9\,meV/H$_2$O and force RMSE of 36.2--37.1\,meV/\AA{} on this dataset.
Converting Table~\ref{tab:waterbox-summary} to the same per-molecule basis
(dividing by 64): energy MAE here is $\sim$75--258\,meV/molecule
(\texttt{water\_absolute}) and $\sim$73--112\,meV/molecule
(\texttt{water\_absolute+momentum}) across seeds --- roughly 40$\times$ to
over 100$\times$ MACE's fully-converged number. Force MAE
(576--2556\,meV/\AA{} and 605--675\,meV/\AA{} respectively) is 15$\times$ to
70$\times$ MACE's force RMSE, well beyond what the MAE-vs-RMSE convention
difference alone would explain.

This is not evidence of a bug --- MACE's numbers reflect a mature,
extensively-tuned architecture trained for presumably far longer than 20
epochs with hyperparameters chosen for this exact system, whereas this
study's water-box runs inherited hyperparameters (batch size aside) from a
completely different single-molecule dataset and were never re-tuned for
absolute accuracy (that was never the goal --- the goal is the
\emph{momentum-violation comparison}, not a competitive water potential).
But it is a useful, concrete sense of how much accuracy headroom remains,
and a reminder to normalize per-atom/per-molecule in any future reporting
so numbers are comparable to published benchmarks
\cite{mlip-error-metrics, mlip-validation}.

\subsection{Why the current comparison isn't trustworthy yet}
\label{sec:caveats}

Digging into the per-epoch training histories (not just the summary table)
revealed that \emph{both} conditions' validation loss oscillates by 1--3
orders of magnitude between adjacent epochs --- it does not converge
smoothly. Checkpoint selection (single epoch with minimum validation loss)
is therefore effectively sampling from noise: e.g.\ one
\texttt{water\_absolute} seed's ``best'' checkpoint came from epoch 3 of 18
and was never matched again for the rest of training, while another seed's
best checkpoint came from epoch 16. The large std in
Table~\ref{tab:waterbox-summary} is substantially an artifact of this
checkpoint-selection process, not a clean measurement of seed-to-seed
variability in what the models actually learned. This means:
\begin{itemize}
  \item Running more seeds would not fix this --- it would just give a more
  precise estimate of a still-not-meaningful quantity (the expected value of
  the minimum of $\sim$16 noisy per-epoch samples).
  \item The right next step is evaluating a \emph{non-cherry-picked}
  checkpoint (e.g.\ the last-epoch checkpoint, already saved via
  \texttt{save\_last=True} at zero additional training cost) across all six
  existing runs, to see whether the comparison holds up, or sharpens, once
  checkpoint-selection noise is removed from the picture.
\end{itemize}

\subsection{Update: the $n=3$ momentum comparison does not survive to $n=6$}
\label{sec:n6-update}

Seeds were extended from 3 to 6 per condition once both prerequisites for
trusting more seeds were satisfied: anneal-invariant checkpoint selection
(Section~\ref{sec:checkpoint-fix}) and confirmed periodic neighbor search
(Section~\ref{sec:periodicity}).

\textbf{Raw best-checkpoint comparison, $n=6$}: energy MAE
$5.94\pm5.27$ (\texttt{water\_absolute}) vs.\ $4.79\pm1.99$
(\texttt{+momentum}); force MAE $0.964\pm0.782$ vs.\ $0.635\pm0.037$; mean
momentum violation $60.9\pm46.4$ vs.\ $43.8\pm5.3$; \textbf{max momentum
violation $1222.9\pm966.6$ vs.\ $1394.1\pm870.8$ --- reversed from the
$n=3$ table above}, absolute now lower. The ``wins on all four metrics''
claim (Table~\ref{tab:waterbox-summary}) does not survive to $n=6$.

\textbf{One seed is doing most of the work}: \texttt{water\_absolute} seed 2
(energy MAE 16.5\,eV, force MAE 2.56\,eV/\AA) is 3--4$\times$ every other
seed in either condition, despite an unremarkable \texttt{best\_model\_score}
(1.83, squarely inside the 1.48--3.15 range of all 11 other seeds) --- the
same checkpoint-selection noise documented in Section~\ref{sec:caveats}, now
visible at the single-seed level instead of smoothed into an aggregate std.
Excluding it: energy MAE reverses to favor absolute ($3.83\pm1.25$ vs.\
$4.79\pm1.99$); force MAE becomes a wash ($0.645\pm0.056$ vs.\
$0.635\pm0.037$); mean violation becomes a wash ($42.1\pm6.4$ vs.\
$43.8\pm5.3$); max violation more clearly favors absolute
($861\pm431$ vs.\ $1394\pm871$). Quietly dropping an inconvenient seed is a
methodological trap, not a fix --- shown here specifically to demonstrate how
much of the original result rode on one data point, not to argue for
dropping it.

\textbf{A more principled, noise-robust check}: the windowed
validation-set comparison already used earlier in this study
(Section~\ref{sec:q1}, item 2 --- average the last 5 epochs' logged
\texttt{val\_y\_l1\_loss}/\texttt{val\_neg\_dy\_l1\_loss} across all 12
cells, immune to single-epoch checkpoint-selection noise since it never
picks one ``best'' epoch at all). Energy: $37.2\pm7.2$ vs.\ $33.3\pm5.5$
(momentum lower, doesn't clear significance either way). Force:
$0.754\pm0.160$ vs.\ $0.622\pm0.073$ (momentum lower and tighter --- the one
signal that holds up under every lens tried).

\textbf{Verdict}: the energy conclusion flips depending on evaluation
methodology (favors momentum in the raw and windowed views, favors absolute
once the noisy outlier is excluded from the raw view) and is not trustworthy
at this $n$. Force quality is the one metric that consistently, modestly
favors momentum across all three lenses. Max per-molecule violation does not
favor momentum under any lens tried. This replaces the original
``directionally consistent on all four metrics'' claim.

\subsection{Real dynamics: both conditions produce unstable NVE rollouts;
momentum consistently runs hotter}
\label{sec:rollout-results}

\textbf{Both conditions fail to hold a real $\sim$300\,K liquid-water
configuration near its starting temperature under plain NVE.} Confirmed
across every rollout in this study ($n=22$ total, including pilots): starting
from a real $\sim$287--335\,K DFT snapshot, both \texttt{water\_absolute} and
\texttt{water\_absolute+momentum} heat to a plateau somewhere in the
$\sim$1000--2600\,K range within well under 1\,ps of integration, then hold
there rather than diverging further, with total-energy drift concentrated
mostly (though not always exclusively) in the initial few hundred fs. This is
exactly the failure mode static accuracy metrics cannot see
\cite{mlip-validation} --- both checkpoints look reasonable on held-out
static configurations (Section~\ref{sec:n6-update}) while being dynamically
unstable once the system evolves under its own predicted forces.

\textbf{Severity is configuration-dependent}: the single starting
configuration used for the first replicate batch below plateaus far hotter
($\sim$1900--2600\,K) than the five other configurations tested in the
second batch ($\sim$1050--1650\,K) --- ``these checkpoints heat a real
snapshot to $\sim$2000\,K'' overstates the general case;
``\ldots to somewhere in 1000--2600\,K, depending on the starting
configuration'' is accurate.

\begin{table}[h]
\centering
\caption{Rollout stability, two independent replicate batches ($n=5$ each,
1\,ps NVE, mean $\pm$ std). \emph{Velocity batch}: one fixed starting
configuration, 5 different Maxwell-Boltzmann velocity draws. \emph{Config
batch}: one fixed velocity draw, 5 different starting configurations (a true
paired comparison, config by config, since the velocity draw is
bit-identical --- Section~\ref{sec:rollout-methods}). Plateau temperature:
mean $\pm$ std over the last 30\% of logged steps.}
\label{tab:rollout-summary}
\begin{tabular}{llcc}
\toprule
Batch & Condition & Drift (meV/atom/ps) & Plateau $T$ (K) \\
\midrule
Velocity ($n=5$) & \texttt{water\_absolute}          & $28.1 \pm 28$  & $1949 \pm 84$ \\
Velocity ($n=5$) & \texttt{water\_absolute+momentum} & $64.9 \pm 27$  & $2317 \pm 104$ \\
Config ($n=5$)   & \texttt{water\_absolute}          & $3.98 \pm 2.7$ & $1285 \pm 180$ \\
Config ($n=5$)   & \texttt{water\_absolute+momentum} & $7.56 \pm 2.5$ & $1459 \pm 140$ \\
\bottomrule
\end{tabular}
\end{table}

\textbf{Momentum consistently runs hotter}: in the config batch's paired
comparison (identical geometry and velocity draw per pair, only the trained
weights differ), momentum's plateau temperature exceeds absolute's in
\textbf{5 of 5} pairs ($+92$ to $+204$\,K); combined with the velocity
batch's full separation (momentum's minimum, 2252\,K, exceeds absolute's
maximum, 2051\,K), that is \textbf{10 of 10} trials across two independent
replicate axes. Drift agrees in 4 of 5 config-batch pairs (reversed for one
configuration) plus the velocity batch's full separation --- 9 of 10. This
does not depend on any single data point the way the $n=3$ static-metric
comparison did, but the underlying \emph{mechanism} --- why momentum
conservation training makes NVE rollouts run hotter, not cooler, than plain
supervised training --- is not yet understood; see Section~\ref{sec:q3}.

\textbf{Confounds caught and fixed along the way}: an initial single-trial
comparison (momentum $3.5\times$ the drift of absolute) turned out to be
confounded by mismatched starting velocities (335\,K vs.\ 292\,K, since
\texttt{MaxwellBoltzmannDistribution} without an explicit \texttt{rng} draws
from whatever state numpy's global RNG happens to be in) --- fixing this
\emph{reversed} which condition showed less drift, before the $n=5$ batches
above established the more robust picture. The same \texttt{seed} used for
the data split was initially reused for the velocity draw too, which would
have silently varied the starting geometry between what were meant to be
velocity-only replicates; split into two independent parameters before
running the batches above (Section~\ref{sec:rollout-methods}).

\subsection{Q3 progress: momentum's rollout systematically drifts toward
higher-net-force geometries}
\label{sec:q3-progress}

\textbf{Point-wise, the two models agree closely --- neither candidate 1 nor
3 holds up.} On real held-out equilibrium configurations
(\texttt{--mode equilibrium}, Section~\ref{sec:force-decomposition-methods}),
momentum's net-force magnitude was not smaller than absolute's (if anything
slightly larger: $1.25$ vs.\ $1.18$\,eV/\AA{} averaged over 6 configs), and
the internal-force share was nearly identical between conditions
($\sim$97\% for both) --- no support for candidate 3 (momentum
``satisfied'' by redistributing force within a molecule rather than reducing
it). Cross-evaluating both models on real rollout trajectory frames
(\texttt{--mode trajectory}) found the same: at every one of 8
(trajectory, frame) combinations tested, the two models' predicted forces on
the identical input geometry differ by only 1--5\%, with no growth from
frame 0 to frame 180 --- no support for candidate 1 either (a
force-accuracy difference that only emerges away from the training
distribution would show a widening gap with frame index; it doesn't).

\textbf{But the two conditions' own trajectories visit systematically
different geometries.} Fixing \texttt{eval\_model} and comparing across
\emph{trajectory} instead: at matched elapsed time, momentum's own rollout
consistently sits at a geometry with higher per-molecule net force than
absolute's own rollout, regardless of which model evaluates it. Replicating
this across all 10 already-collected trajectory pairs
(\texttt{run\_force\_decomposition\_study.py}, both velocity- and
config-axis batches pooled) confirms it is not an artifact of the one pair
first examined, and reveals something sharper: the gap grows monotonically
from an exact zero.

\begin{table}[h]
\centering
\caption{Momentum-minus-absolute per-molecule net-force gap at matched
elapsed time, pooled across 10 replicates (5 velocity-draw + 5 config,
Section~\ref{sec:rollout-results}). Frame 0 is the identical starting
geometry within each replicate by construction.}
\label{tab:force-gap-by-frame}
\begin{tabular}{lccc}
\toprule
Frame (step) & Gap (mean $\pm$ std, eV/\AA) & $n_{\mathrm{positive}}$ \\
\midrule
0 (0)      & $4.5{\times}10^{-8}$        & 4/10 \\
20 (200)   & $0.050 \pm 0.153$           & 6/10 \\
80 (800)   & $0.085 \pm 0.163$           & 8/10 \\
180 (1800) & $0.351 \pm 0.261$           & \textbf{10/10} \\
\bottomrule
\end{tabular}
\end{table}

At frame 0 the gap is machine-precision zero and $n_{\mathrm{positive}}$
(4/10) is indistinguishable from a coin flip, exactly as expected since both
conditions start from the literal same geometry within each replicate. By
frame 180 it is unanimous across all 10 independent replicates. Under a true
null of no real effect, 10/10 agreement by chance has probability
$\approx 2\times(0.5)^{10} \approx 0.002$ --- this is a real, replicated
pattern, not an artifact.

\textbf{Interpretation: a systematic bias, not undirected chaos.} If the two
force fields differed only by ordinary numerical noise with no consistent
bias --- plain sensitive dependence in a nonlinear many-body system, which
would amplify whatever small difference exists in whatever direction it
happens to point --- the \emph{direction} of divergence should vary across
independent replicates: sometimes momentum ends up at the higher-net-force
geometry, sometimes absolute does, roughly at random. Getting the same
direction in all 10 independent replicates (different starting
configurations \emph{and} different velocity draws) rules that out. There
must be a small, systematic, repeatable bias between the two models'
predicted forces --- not just noise --- that consistently pushes momentum's
trajectory toward higher-net-force states, and it is that systematic
component, compounding over $\sim$900\,fs of nonlinear dynamics, that
eventually produces the large, obvious rollout-temperature effect. The
point-wise comparisons above (``the two models agree closely'') were
eyeballing means over a handful of configs --- exactly the kind of small
systematic bias implied here could be sitting inside that noise band without
showing up as an obvious mean difference.

\textbf{The systematic bias, found: net force specifically, not torque.} A
rigorous paired test (\texttt{--mode paired-bias}) --- the momentum-minus-
absolute difference at every individual atom/molecule across 100 held-out
configs, not an eyeballed mean --- finds a large, highly significant,
consistently-signed bias in exactly the quantity the trajectory-divergence
result implicated: net per-molecule force is $0.073$\,eV/\AA{} higher for
momentum on average ($z=+23.8$, $p\approx4{\times}10^{-125}$ across 6587
molecule-pairs), and 67\% of individual molecule-pairs show momentum higher,
not just barely over half. Raw and internal force magnitude show the
opposite (tiny, $<1\%$) sign, so this is specific to the net component, not
a general ``momentum's forces are bigger'' effect. Net torque --- checked as
a candidate reconciliation, since the aggregate \texttt{per\_fragment\_
momentum\_loss} combines net force$^2$ and net torque$^2$ --- shows no
comparable compensating bias (mean not significantly different from zero,
$p=0.85$; the sign test is significant only because $n$ is large, and leans
the same direction as net force, not opposite).

\textbf{No compensating torque effect was needed, because there was no real
contradiction to reconcile.} This seed pair's own static evaluation
(Table~\ref{tab:waterbox-summary}'s underlying per-seed data,
\texttt{results/waterbox\_study/raw\_results.csv}) already showed
\texttt{water\_absolute+momentum} seed 0 with \emph{higher} (worse) mean
per-molecule momentum violation than \texttt{water\_absolute} seed 0
($45.16$ vs.\ $41.93$) --- the apparent tension came from comparing this
seed-0-specific force-decomposition result against the $n{=}6$ \emph{seed
average} instead of seed 0's own number. That average
(Section~\ref{sec:n6-update}) is itself dominated by \texttt{water\_
absolute} seed 2's outlier; the six seeds individually split roughly
3-and-3 on which condition has the lower per-fragment violation, and seed 0
--- the one pair used for every rollout and every force-decomposition check
in this investigation --- happens to be one of the three where momentum is
worse on this metric. The rollout instability isn't decoupled from the
static per-fragment evaluation; it's the dynamical consequence of something
that evaluation already showed for this specific seed pair, once the right
number is checked rather than the aggregate.

\subsection{Confirming check: does the effect track which seed's static
metric favors momentum?}
\label{sec:q3-seed5}

If seed 0's rollout instability is the dynamical consequence of that
specific training run's static per-fragment momentum violation being worse
for \texttt{water\_absolute+momentum} (Section~\ref{sec:q3-progress}), a
seed where the static metric instead favors momentum should show a smaller
or reversed rollout-stability gap. Seed 5 is such a seed
(Section~\ref{sec:n6-update}'s per-seed data: mean violation $36.28$
vs.\ absolute's $41.30$, momentum $\sim$12\% lower). \texttt{run\_rollout\_
study.py} gained a \texttt{--train-seed} option (defaulting to 0, so every
existing result is unaffected) to repeat both replicate batches on seed 5's
checkpoints instead.

\begin{table}[h]
\centering
\caption{Rollout stability, seed 0 vs.\ seed 5, both replicate axes
($n=5$ each, mean $\pm$ std).}
\label{tab:seed5-comparison}
\begin{tabular}{lcccc}
\toprule
& Drift, absolute & Drift, momentum & Plateau $T$, absolute & Plateau $T$, momentum \\
& (meV/atom/ps) & (meV/atom/ps) & (K) & (K) \\
\midrule
Seed 0, velocity & $28.1 \pm 28$ & $64.9 \pm 27$ & $1949 \pm 84$ & $2317 \pm 104$ \\
Seed 0, config    & $3.98 \pm 2.7$ & $7.56 \pm 2.5$ & $1285 \pm 180$ & $1459 \pm 140$ \\
Seed 5, velocity  & $7.65 \pm 19$ & $\mathbf{4.04 \pm 6.5}$ & $967.7 \pm 86$ & $1036 \pm 20$ \\
Seed 5, config    & $-0.74 \pm 0.35$ & $0.70 \pm 3.4$ & $738.6 \pm 130$ & $895.6 \pm 200$ \\
\bottomrule
\end{tabular}
\end{table}

\textbf{The effect weakens substantially at seed 5, but doesn't cleanly
reverse.} At seed 0, the momentum-minus-absolute plateau-temperature gap is
large ($+368$\,K, $+174$\,K) with little to no overlap between the two
conditions' distributions. At seed 5 the gap is much smaller ($+68$\,K,
$+157$\,K) and now sits well inside the conditions' own spread (absolute's
std alone is 86--130\,K) --- the clean, unambiguous separation seen at
seed 0 is gone. Temperature still leans toward momentum running slightly
hotter in both of seed 5's axes, not reversed, but no longer decisively.

\textbf{Drift specifically stops being a reliable signal at seed 5.} At
seed 0, momentum showed higher drift in both replicate axes, consistently.
At seed 5, momentum shows \emph{lower} drift in the velocity-draw axis but
\emph{higher} drift in the config axis --- the sign flips depending on which
axis is checked, using the identical two checkpoints both times. That
internal inconsistency is itself informative: it indicates this comparison
has become noise-dominated at seed 5, in sharp contrast to seed 0's
unanimous 10/10 direction.

\textbf{Verdict}: the seed-0 finding was real but seed-specific, not a
general property of the momentum loss term --- moving to a seed where the
static metric already favored momentum substantially shrank the
rollout-stability gap, exactly as the mechanism in
Section~\ref{sec:q3-progress} predicted. What remains at seed 5 is a
residual, weak, and only partially consistent lean toward momentum being
slightly less stable even here, not a full disappearance or reversal of the
effect. This is weaker than a full confirmation but is not a null result
either: seed-to-seed variation in whether the momentum loss achieves its
intended effect appears to be the dominant factor in the rollout-stability
question, more than any fixed property of ``momentum-conservation training''
as a method.

\section{Open questions / next steps}
\label{sec:next}

Four threads now. \textbf{Q4 (top priority}, Section~\ref{sec:q4}):
\emph{both} conditions produce dynamically unstable rollouts regardless of
the momentum question --- a literature-supported hypothesis (missing
short-range repulsion, Ziegler-Biersack-Littmark) and a concrete, already
partly-tooled next step exist but haven't been run yet. Q1 (static momentum
effect) is weaker than it looked but not fully closed out. Q3 (why rollouts
run hotter under momentum training specifically) has a well-supported answer
for seed 0 (a systematic net-force bias present in that seed's own static
evaluation) that substantially weakens, but doesn't fully vanish, at a seed
where the static metric favors momentum --- seed-to-seed training variance,
not a fixed property of the momentum loss, is the best-supported explanation
overall. Q2 (real-dynamics evaluation generally) continues alongside Q4
rather than blocking it.

\subsection{Q1 (pressing): does the per-molecule momentum loss actually help?}
\label{sec:q1}

Current status: directionally consistent improvement across all four
metrics, tighter variance, but not yet distinguishable from checkpoint-
selection noise (Section~\ref{sec:caveats}). Concrete next steps, roughly in
order:
\begin{enumerate}
  \item \emph{Done, free, zero training cost}: re-evaluated all six existing
  checkpoints at what was intended to be the deterministic last-epoch
  checkpoint (\texttt{last.ckpt}). Result: nearly identical to the original
  best-checkpoint numbers for every one of the six runs --- not because of a
  script bug, but because Lightning's \texttt{save\_last=True} only updates
  ``whenever a checkpoint file gets saved,'' which under
  \texttt{save\_top\_k=1}/\texttt{mode="min"} only happens on improvement.
  Since at least one seed's global-minimum validation loss occurred at
  epoch 3 of 18 and was never beaten again, \texttt{last.ckpt} silently
  froze at the same epoch-3 weights as the ``best'' checkpoint --- it was
  never actually the final epoch. Fixed for the retrain below via a second,
  unconditional \texttt{ModelCheckpoint} (\texttt{monitor=None},
  \texttt{save\_top\_k=0}, \texttt{save\_last=True}), which has no
  improvement gate and so genuinely overwrites every epoch.
  \item \emph{Done, free, zero training cost}: windowed sanity check using
  only the already-logged per-epoch validation histories (last 5 epochs'
  \texttt{val\_y\_l1\_loss}/\texttt{val\_neg\_dy\_l1\_loss}, averaged, instead
  of a single cherry-picked epoch). Split result: the \emph{energy}
  comparison held up and even strengthened (momentum's windowed average
  $42.4 \pm 7.0$ vs.\ absolute's $83.3 \pm 26.0$ --- gap now exceeds
  absolute's own std, unlike the best-checkpoint comparison). The
  \emph{force} comparison did not survive ($0.678 \pm 0.051$ vs.\
  $0.647 \pm 0.064$ --- statistically indistinguishable, momentum nominally
  slightly worse) --- suggesting the original ``momentum clearly wins on
  force'' result was mostly a checkpoint-selection artifact, while the
  energy result looks more likely to be real. Could not address the
  momentum-violation metric this way --- it is never logged during
  training, only computed post-hoc from real model weights.
  \item \emph{Done}: full retrain, all 6 cells, 50 epochs with the anneal
  schedule active (Section~\ref{sec:anneal}) and \texttt{early\_stop\_patience=40}.
  All six cells reached \texttt{final\_epoch}=48 (of 50 requested) - the
  windowed pre/post-anneal comparison (item 2 above, computed for real this
  time via \texttt{src/analyze\_anneal\_effect.py} rather than by hand) held
  up: 5 of 6 cells' median post-anneal energy error improved over pre-anneal,
  all 6 cells' median force error improved despite force's loss weight
  dropping from 0.95 to 0.25. But this retrain still used the \emph{old},
  biased \texttt{val\_total\_mse\_loss} monitor for actual checkpoint
  selection, so its reported best-checkpoint summary numbers were not yet
  trustworthy on their own terms.
  \item \emph{Done}: fixed the checkpoint-selection bug itself
  (Section~\ref{sec:checkpoint-fix}) and re-ran the full 6-cell matrix a
  second time. Result: \textbf{5 of 6 cells now select a genuinely
  post-anneal epoch as best}, versus 0 of 6 before the fix - direct
  confirmation the anneal schedule is doing real work, not an artifact of a
  biased selection metric. The one holdout
  (\texttt{water\_absolute+momentum} seed 2, best epoch 13) isn't evidence of
  a remaining bug - it's simply a seed where no post-anneal epoch happened to
  beat what pre-anneal already found, which is a legitimate outcome once the
  comparison is fair. Reported summary numbers (Table~\ref{tab:waterbox-summary})
  are, notably, \emph{unchanged to reported precision} from every earlier
  version of this study despite the underlying checkpoints genuinely
  differing (different epochs, confirmed via wall-clock and the new
  \texttt{val\_checkpoint\_score} column) - the most likely reading is that
  this model reliably settles into a fairly consistent quality range once
  past its initial noisy phase, regardless of which specific well-converged
  epoch (pre- or post-anneal) ends up selected.
  \item \textbf{Net conclusion so far, superseded}: at $n=3$, this item
  concluded the headline comparison was unchanged from the first honest
  read and standing on solid methodological footing, with significance the
  only open question. \textbf{That does not hold at $n=6$}
  (Section~\ref{sec:n6-update}): the ``wins on all four metrics'' claim is
  gone, the energy conclusion flips depending on evaluation methodology, and
  only the force-quality improvement survives every lens tried. The
  checkpoint-selection footing itself is still solid --- what changed is
  that more seeds revealed how much of the original comparison rode on one
  noisy seed, exactly the outcome running more seeds was supposed to check
  for.
  \item \textbf{Superseded again, by real dynamics}: the static-metric
  comparison above is no longer the most informative evidence available.
  Section~\ref{sec:q1-rollout-update} reports the real-dynamics comparison
  under the corrected (ZBL-bonded-exclusion + anneal-fixed) methodology ---
  a materially stronger and more consistent signal than anything in this
  subsection.
\end{enumerate}

\subsection{Q2: evaluating trained potentials in real dynamics}
\label{sec:q2}

\textbf{Status: initial results in hand} (Section~\ref{sec:rollout-results}
has the full story) --- both conditions produce unstable NVE rollouts from a
realistic starting configuration, momentum consistently running hotter
across 10 of 10 replicate trials. Methodology in place
(Section~\ref{sec:rollout-methods}): ASE-driven NVE rollouts, energy-drift
tracking, and O--O/O--H/H--H RDF against a real-configuration reference.
The RDF side of the diagnostic hasn't been analyzed in the same depth as the
energy/temperature numbers yet --- worth revisiting windowed by time
(early/transient vs.\ plateau region) rather than over a whole trajectory
that spends most of its length in an already-overheated, non-representative
state. \todo{diffusion coefficient / other dynamical observables - read more
before committing to a list here.}

\subsection{Q4 (new, now the top priority): why are BOTH conditions
dynamically unstable at all?}
\label{sec:q4}

Independent of the momentum-vs-absolute question (Q3), every rollout in this
study --- either condition, any seed, any starting configuration tried so
far --- heats a real $\sim$300\,K liquid-water snapshot to somewhere in
1000--2600\,K within under a picosecond of plain NVE
(Section~\ref{sec:rollout-results}). This is the more fundamental problem:
a potential that cannot hold a realistic starting configuration near its own
temperature isn't a usable water potential regardless of how the physics-loss
comparison resolves.

\textbf{Literature-supported hypothesis: missing short-range repulsion.}
\citet{fu2022forces} is the paper that essentially predicts this failure
mode --- force/energy accuracy on static held-out data does not guarantee
stable MD, and they propose rollout stability itself as the metric that
actually matters, exactly the gap this project's Q2 was built to probe. The
literature's standard diagnosis for \emph{why} static accuracy fails to
prevent this: real AIMD training data rarely samples atoms at anomalously
short distances (high-energy, low-probability configurations), so a purely
learned potential has no reliable training signal there. Once a rollout
becomes energetic enough to wander into a slightly-too-close configuration
--- exactly what an already-heating trajectory does --- the network can
predict an arbitrarily wrong force (even attractive instead of repulsive)
instead of the strong repulsion that should be there, releasing energy and
compounding the instability.

The standard fix in the literature is an explicit, non-learned repulsive
term at short range, almost always the Ziegler--Biersack--Littmark (ZBL)
screened-nuclear-repulsion potential \citep{zbl1985}, added directly to the
predicted energy so the network is never solely responsible for short-range
repulsion. Both architectures already cited in this project use it: NequIP
ships a dedicated \texttt{ZBL} prior class specifically to ``mitigate
molecular dynamics failure modes associated with atoms getting too close'';
MACE includes a ZBL-based pairwise repulsion term in its foundation models by
default, for the same reason \cite{nequip, mace}. \citet{shortrangeprior2025}
reports the same benefit directly: adding a short-range empirical (ZBL-style)
term measurably improves both simulation robustness and training efficiency.

\textbf{This project's own tooling already supports it.} \texttt{torchmdnet}
--- the framework used throughout this study --- ships a ready-to-use
\texttt{ZBL} prior (\texttt{torchmdnet/priors/zbl.py}), selectable via the
same \texttt{prior\_model} argument currently set to \texttt{None} in
\texttt{train\_waterbox.py}. Its constructor needs \texttt{cutoff\_distance},
\texttt{max\_num\_neighbors}, and unit-conversion factors
(\texttt{distance\_scale}, \texttt{energy\_scale}) to bridge ZBL's native
unit convention to this dataset's \AA/eV convention --- exactly the kind of
unit-bridging step this project has gotten wrong before (the Bohr/Hartree
mixup in \texttt{waterbox\_data.py}, Appendix~\ref{app:impl} item 2), so it
needs the same empirical verification (check against physical priors) rather
than trusting the defaults.

\textbf{First real run was an artifact, caught before being trusted.}
\texttt{src/diagnose\_short\_range\_collapse.py} scans every trajectory
already on disk (\texttt{results/waterbox\_rollout*/**/rollout.xyz}) for
anomalously short \emph{cross-molecule} O--O/O--H/H--H distances
(minimum-image aware, reusing
\texttt{baseline\_potential.\_infer\_bonds\_from\_positions}'s own
displacement-minus-nearest-box-multiple convention rather than a fresh one),
and checks whether the first such event precedes, coincides with, or follows
each trajectory's heating-transient onset (temperature crossing halfway
between its starting and plateau value). Run once against all 42 saved
trajectories, it reported a summary that looked like a clean confirmation
(24/42 ``precedes'', 18/42 ``coincides'', 0 ``follows'', 0 ``no
violation'') --- but every single trajectory's \emph{first} sub-floor event
was at frame 0/$t{=}0$, the real, unheated DFT starting snapshot itself,
before the model had taken a single integration step. That is not possible
for a genuinely model-caused collapse signal, and inspecting the raw
per-frame CSVs (not just the summary table) confirmed it directly: e.g. the
single-replicate \texttt{waterbox\_rollout} trajectory's frame 0
($T{=}287$\,K) already contained an O--H pair at 1.06\,\AA{} (floor
1.63\,\AA) and an H--H pair at 1.37\,\AA{} (floor 1.44\,\AA), identically
reproduced in \texttt{waterbox\_rollout\_momentum}'s frame 0 (same starting
geometry, as expected) and in a second, independent starting configuration.

Two calibration bugs, both now fixed: (1) the H--H check excluded only
literal covalent bonds (via \texttt{structural\_metrics.infer\_bonds}), not
same-molecule pairs --- a water molecule's own two hydrogens sit
$\sim$1.5\,\AA{} apart purely from the H--O--H bond angle, with no bond
between them, so every molecule's own H\ldots H distance polluted the H--H
statistic from frame 0 onward. Fixed by excluding same-\emph{molecule} pairs
(\texttt{structural\_metrics.infer\_molecule\_groups}) instead of only
literal bonds --- covers the O--H covalent bonds and the intramolecular
H\ldots H geometry in one consistent step. (2) the ``anomalous'' floor was a
percentile computed by pooling every individual pairwise distance across
$\sim$200 reference configs (thousands of pairs each), then compared against
each rollout frame's own \emph{minimum} --- an order statistic checked
against a population-wide percentile, which sits below that threshold for
essentially any config purely from having thousands of pairs to draw a
minimum from, healthy water included. Fixed by computing each reference
config's own minimum \emph{first}, then taking a percentile over that much
smaller pool of per-config minima (the same statistic on both sides of the
comparison) --- default floor percentile raised from 0.1 to 1.0 accordingly,
since the pool shrank from thousands of samples to $\sim$200. Molecule
membership (frame 0 only, reused for later frames, same rationale as the
original bond-freezing design --- distance-based inference on an
already-collapsed frame is unreliable) replaces the earlier bonds-only
exclusion throughout. Both fixes are synthetically unit-tested locally,
including a regression test built directly from the intramolecular-H--H bug
found in the real run.

\textbf{Corrected re-run: a genuinely mixed result, not a uniform
confirmation.} Across all 42 trajectories: 15 PRECEDE the heating onset, 11
COINCIDE (26/42, 62\%, at or before onset), \textbf{16 FOLLOW} (38\% ---
short-range collapse detected only after heating is already underway), 0
show no cross-molecule violation at all. This splits cleanly by training
seed and replicate axis, not randomly:

\begin{itemize}
\item \textbf{Seed 0, velocity-draw batch (10/10): unanimous PRECEDES.}
Spot-checked directly (\texttt{waterbox\_rollout}, \texttt{water\_absolute}):
frame 0 is clean (0 violations everywhere, confirming the fix), but by the
very next logged frame (5\,fs) temperature has already risen $287\to548$\,K
\emph{and} a cross-molecule O--H pair has dropped to 0.80\,\AA{} (floor
0.97\,\AA) with an H--H pair at 0.82\,\AA{} (floor 0.91\,\AA) --- before the
temperature-defined onset threshold (55\,fs) is even crossed.
\item \textbf{Seed 0, config batch (10): mixed} --- 3 precede, 1 coincide, 6
follow, with \texttt{cfg3}/\texttt{cfg4}/\texttt{cfg5} following for
\emph{both} conditions. Consistent with the already-known configuration
dependence of rollout severity (Section~\ref{sec:rollout-results}): these
are among the milder, cooler-plateau configurations.
\item \textbf{Seed 5, config batch (10/10): unanimous FOLLOWS.}
Spot-checked (\texttt{cfg1}, \texttt{water\_absolute}): temperature climbs
steadily $304\to760$\,K over the first 115\,fs with \emph{zero}
cross-molecule violations anywhere in that window; the first violation
doesn't appear until 265\,fs, well after the heating is already underway.
\item \textbf{Seed 5, velocity-draw batch (10/10): COINCIDES, with a
caveat.} This batch's heating is explosively fast --- temperature jumps
$\sim$300$\to$700\,K within the very first logged interval (5\,fs at the
standard \texttt{energy\_log\_stride=10}/\texttt{dt=0.5}) --- so
precedes-vs-coincides isn't meaningfully resolvable at this logging
resolution for this batch; the whole transition happens inside one interval.
\end{itemize}

\textbf{Conclusion: missing short-range repulsion is a real contributing
mechanism for a meaningful fraction of trajectories --- cleanly and
unanimously so for seed 0's velocity-draw replicates --- but not a universal
explanation for Q4.} Seed 5's config batch heats substantially with no
detectable close contact anywhere near the transition, implying at least one
other mechanism also drives the instability there. Next step: enable the ZBL
prior in \texttt{train\_waterbox\_model} and retrain, then re-run the
rollout-stability comparison (\texttt{run\_rollout\_study.py}) to check
whether the shared instability resolves --- but going in expecting a
\emph{partial} fix given the above, not full resolution; if instability
persists (especially in seed-5-config-like trajectories), that would point to
a second, distinct mechanism (e.g.\ an integrator/energy-conservation issue,
or a systematic force bias unrelated to short-range contacts) rather than a
failure of the ZBL fix itself.

\textbf{ZBL wiring is now in place, retrain not yet run.}
\texttt{train\_waterbox\_model} gained an opt-in \texttt{use\_zbl\_prior} flag
(default \texttt{False}, so every existing checkpoint/result stays
reproducible unchanged) that sets \texttt{model\_args["prior\_model"]="ZBL"}
with \texttt{prior\_args} supplying \texttt{cutoff\_distance},
\texttt{max\_num\_neighbors}, \texttt{atomic\_number} (the identity map
\texttt{list(range(100))} --- WaterBox's \texttt{z} field already stores
literal atomic numbers, not a compacted type index), and
\texttt{distance\_scale}/\texttt{energy\_scale} bridging this project's
\AA/eV convention to ZBL's native meters/Joules formula (exact constructor
signature and per-pair math pulled directly from torchmd-net's public
source, since this checkout has no installed copy to introspect). The
\texttt{cutoff\_distance=2.0}\,\AA{} default was chosen empirically via
\texttt{src/verify\_zbl\_units.py} --- a standalone reimplementation of
ZBL's exact screening-function/cosine-cutoff formula that needs no
torchmdnet install, letting it be checked directly rather than guessed (the
same discipline this project already applies to unit-bridging steps, per
the Bohr/Hartree lesson). Comparing candidate cutoffs
(2.0/2.5/3.0/4.0/5.0\,\AA) against both the empirical anomaly floors above
and typical equilibrium separations: 2.0\,\AA{} gives a large repulsive
correction right at the anomaly floors (0.6--3.8\,eV) while staying exactly
zero at real O--O/H--H equilibrium distances and only 0.003\,eV at the O--H
hydrogen-bond distance; cutoffs of 2.5\,\AA{} or more start leaking a
non-negligible correction (up to 0.30\,eV at 5.0\,\AA) into the
hydrogen-bond distance itself, which would fight the bulk-water physics the
network already learned correctly rather than only fixing short-range
collapse. A nonzero ZBL contribution at the covalent O--H bond length itself
($\sim$2.6\,eV at the 2.0\,\AA{} cutoff) is expected, not a bug --- ZBL has no
concept of ``this pair is a real bond'' (identical to NequIP/MACE's own
usage); the network is expected to learn to compensate for it during
retraining. \texttt{run\_waterbox\_study.py --use-zbl-prior} retrains both
conditions into a separate \texttt{checkpoints/waterbox\_study\_zbl/} /
\texttt{results/waterbox\_study\_zbl/} root, mirroring
\texttt{run\_rollout\_study.py}'s \texttt{--train-seed} convention of never
overwriting the existing reference checkpoints/results.

\textbf{Retrain has run; it surfaced a ZBL-specific anomaly that needs
resolving before trusting the static comparison.}
\texttt{results/waterbox\_study\_zbl/raw\_results.csv}: \texttt{water\_absolute}
seed 0 alone shows energy MAE 16404\,eV, force MAE 6.87\,eV/\AA{} ---
an order of magnitude worse than every other cell (all at 2.4--5.3\,eV). This
is not the pre-existing checkpoint-selection noise already documented
elsewhere in this project: the identical seed trained completely normally
\emph{without} ZBL (energy MAE 4.78\,eV, matching the number already cited
above), and every other water\_absolute seed and every
water\_absolute+momentum seed --- including its own seed 0 --- trained fine
under ZBL. The training history
(\texttt{best\_model\_history.csv}) shows \texttt{val\_checkpoint\_score}'s
true minimum landing at the last epoch (48) with an unremarkable
\texttt{val\_y\_mse\_loss}\,$\approx$\,18.9 --- validation looked fine; test
did not. \textbf{Working hypothesis, not yet confirmed}: ZBL's
$1/\text{distance}$ singularity means a single anomalously-close held-out
test configuration (Section~\ref{sec:q4} already found real DFT
configurations occasionally have surprisingly close non-bonded pairs, e.g.\
0.95\,\AA{} O--H) could dominate a $\sim$160-config mean on its own if this
particular checkpoint hasn't learned to compensate well at that exact
distance --- consistent with only the unregularized condition hitting it,
and with validation (a different random subset) not. Added
\texttt{evaluate\_waterbox.py --per-sample} (per-test-config energy/force/
momentum error plus minimum inter-molecular distance, sorted worst-first) to
check this directly.

\textbf{Singularity hypothesis refuted; root cause bounded to an isolated
per-checkpoint fragility, not fully explained.} \texttt{--per-sample} on the
seed-0 checkpoint showed the error is nearly UNIFORM across all 160 test
configs (16480--16530\,eV, a $<$1\% spread) with no correlation to
\texttt{min\_intermolecular\_distance} --- a rare close-contact event would
show one large outlier with a strong distance dependence, not this. Checking
\texttt{config.json}'s \texttt{post\_test.test\_y\_l1\_loss} (Lightning's own
\texttt{trainer.test()}, reloading the identical checkpoint and evaluating
the identical test split) across every ZBL and non-ZBL cell: 11 of 12 show
close agreement between Lightning's number and
\texttt{evaluate\_waterbox\_checkpoint}'s own hand-rolled number, ruling out
a general tooling bug --- and, critically, seed 0's own
\texttt{post\_test.test\_y\_l1\_loss}\,=\,3.0\,eV: Lightning's trusted path
says this exact checkpoint is fine. Also ruled out a box-tensor dtype
mismatch between the two evaluation paths by checking torchmd-net's own
\texttt{WaterBox} dataset source directly --- \texttt{box} is explicitly cast
to \texttt{torch.float} at load time, identical either way.

What remains, bounded but not fully explained: Lightning evaluates
\texttt{water\_absolute} in batches of 2
(\texttt{FIXED\_HPARAMS["batch\_size"]}); \texttt{evaluate\_waterbox\_
checkpoint}'s hand-rolled loop --- and, critically,
\texttt{waterbox\_ase.py}'s \texttt{TensorNetCalculator}, used by every
rollout step --- evaluates one example at a time. This one checkpoint's
learned weights are apparently unusually sensitive to that difference; every
other seed tolerates it fine. \textbf{Practical implication: seed 0's
\texttt{water\_absolute}+ZBL checkpoint is unsafe for any batch-of-1
evaluation, which is exactly what a rollout does at every step} --- using it
as originally planned (the default \texttt{--train-seed 0}) would have
silently produced garbage rollout data attributable to this artifact, not to
ZBL's real effect on stability. Seed 1 is verified safe for both conditions
under both evaluation paths (\texttt{post\_test.test\_y\_l1\_loss}: 3.42\,eV
/ 4.46\,eV for absolute/momentum; matches \texttt{evaluate\_waterbox\_
checkpoint}'s own batch-of-1 numbers closely) --- use it instead.

\textbf{Rollout study (seed 1, dt=0.5, the original default) has run - and
both batches diverged numerically, not just ``were unstable.''}
\texttt{results/waterbox\_rollout\_study\_zbl\_seed1/summary\_table.md}
(velocity axis): plateau temperatures in the millions to hundreds of
millions of K (\texttt{water\_absolute} up to 17.3M\,K,
\texttt{water\_absolute+momentum} up to 264.7M\,K) --- not physical, a
numerical explosion, not the $\sim$1000--2600\,K instability established
above. \texttt{summary\_table\_by\_config.md} (config axis): same pattern.
\textbf{Root cause: the existing 0.5\,fs integration timestep, tuned for the
smooth learned-only potential, is too coarse for ZBL's steep short-range
repulsive force} --- a well-documented risk of adding hard repulsive priors
to NVE MD, not a wiring/calibration bug. Confirmed directly: re-ran the
single worst replicate (\texttt{water\_absolute} seed 1, vseed1, which hit
17.3M\,K at dt=0.5) at dt=0.1 (5x smaller) ---
\texttt{results/zbl\_dt\_test/water\_absolute\_vseed1\_dt0.1/
energy\_history.csv} shows near-perfect total-energy conservation
($E_\text{tot}$ drifts $\sim$0.005\,eV over 400\,fs, essentially float32
noise at this $\sim{-}$30{,}009\,eV absolute scale) while temperature climbs
smoothly and physically (300\,K $\to$ $\sim$4023\,K, not yet plateaued at
400\,fs) --- no explosion. This rules out a ZBL units/calibration bug for
this checkpoint; the earlier catastrophic numbers are integrator artifacts,
not evidence about ZBL's real effect on stability.

\subsection{Q4 resolved (negative result): ZBL, as configured, makes rollout
stability substantially worse, replicated across two independent seeds}
\label{sec:q4-negative-result}

Both replicate batches were redone at the corrected timestep
(dt=0.1, 10000 steps = 1000\,fs, matching the original comparison window) at
seed 1, and a full ZBL-vs-no-ZBL, same-seed, same-timestep ablation was run
at a second independent seed (5) to check whether the result generalizes.

\begin{table}[h]
\centering
\caption{Plateau temperature (K), mean $\pm$ std over $n=5$ replicates, at
the corrected dt=0.1 timestep. Every ZBL cell is compared against the
matched no-ZBL cell at the identical seed, replicate axis, and timestep.}
\label{tab:zbl-negative-result}
\begin{tabular}{llcc}
\toprule
Seed & Axis / condition & No ZBL & ZBL \\
\midrule
1 & velocity, absolute          & $1342 \pm 16$  & $4615 \pm 76$ \\
1 & velocity, momentum          & $1227 \pm 39$  & $1759 \pm 212$ \\
1 & config, absolute            & $843 \pm 143$  & $4598 \pm 190$ \\
1 & config, momentum            & $790 \pm 153$  & $2511 \pm 650$ \\
5 & velocity, absolute          & $933 \pm 17$   & $5052 \pm 110$ \\
5 & velocity, momentum          & $1061 \pm 14$  & $5285 \pm 81$ \\
5 & config, absolute            & $742 \pm 140$  & $5066 \pm 160$ \\
5 & config, momentum            & $838 \pm 170$  & $5375 \pm 180$ \\
\bottomrule
\end{tabular}
\end{table}

\textbf{Every one of 8 matched cells (40 ZBL replicates vs.\ 40 no-ZBL
replicates total) shows ZBL substantially hotter, by a factor of roughly
3--7$\times$} --- checked at the level of individual replicates, not just
aggregates, and every individual ZBL replicate exceeds every individual
no-ZBL replicate in its matched cell (no overlap anywhere). Drift stays
negligible in every case (well under 1\,meV/atom/ps both with and without
ZBL), so this is not a residual integrator artifact --- both regimes are
genuinely energy-conserving, settled dynamical outcomes; ZBL's dynamics
simply settle at a much hotter plateau. \textbf{Verdict: this specific ZBL
configuration does not fix the shared Q4 instability --- it makes it
substantially worse, replicated across two independent training seeds.}

\textbf{Static accuracy failed to predict this too, compounding the exact
problem this line of inquiry started from.} \citet{fu2022forces}'s central
claim --- that static energy/force accuracy on held-out configurations does
not predict MD stability --- motivated testing real dynamics in the first
place (Section~\ref{sec:rollout-results}). Here it goes one step further:
Lightning's own \texttt{trainer.test()} on these exact ZBL checkpoints
reported energy MAE of 3--5\,eV, indistinguishable from the non-ZBL
checkpoints (Section~\ref{sec:q4}) --- static accuracy did not just fail to
predict the \emph{original} instability, it also failed to predict that this
specific, literature-motivated \emph{fix} would backfire.

\textbf{Leading hypothesis for the mechanism, grounded in how ZBL is
actually implemented elsewhere.} torchmd-net's \texttt{ZBL} prior exposes a
freely-chosen \texttt{cutoff\_distance} with no built-in guidance on how to
pick it, and applies uniformly to every pair within that cutoff --- covalent
bonds included, since the class has no concept of molecular topology. The
\texttt{verify\_zbl\_units.py} check in Section~\ref{sec:q4} already showed
the chosen 2.0\,\AA{} cutoff injects a large ($\sim$2.6\,eV) contribution
directly onto
every O--H covalent bond, continuously, for the entire simulation --- not
just at the rare close non-bonded encounters it was meant to correct.
NequIP's own ZBL implementation \citep{nequip-code} handles this
differently: it exposes no independent cutoff at all, instead reusing the
main model's own graph edges and relying on the ZBL screening function's own
exponential decay to vanish before reaching equilibrium separations. That
decay length ($a = 0.8854\,a_0/(Z_1^{0.23}+Z_2^{0.23})$, the same formula
used throughout this implementation) works out to just 0.18\,\AA{} for O--H,
0.15\,\AA{} for O--O, and 0.23\,\AA{} for H--H in this system --- an order of
magnitude tighter than the 2.0\,\AA{} cutoff chosen here. The common
convention of setting a ZBL-style cutoff at the sum of covalent radii would
give 0.97/1.32/0.62\,\AA{} for O--H/O--O/H--H respectively --- but every one
of those falls \emph{below} the empirical anomaly floors this project's own
diagnostic established (0.79--1.8\,\AA, Section~\ref{sec:q4}), meaning that
convention would have been too conservative to ever reach the actual
collapse zone in this system. \textbf{This looks like a genuine tension
specific to a hydrogen-bonded system like water}: the anomalous non-bonded
approach distance and the covalent bond length sit closer together here than
in the systems ZBL priors are typically deployed for, so a cutoff wide enough
to reach the former necessarily also overlaps the latter substantially.
Classical force fields (and this project's own
\texttt{baseline\_potential.py}, for its CHARMM baseline) handle exactly
this tension with explicit 1-2 (bonded) exclusion lists; torchmd-net's stock
\texttt{ZBL} prior has none. \citet{shortrangeprior2025}'s own positive
result used LLZO, a solid electrolyte with no hydrogen-bond network --- a
system where this specific overlap may simply not arise, which would be
consistent with why their implementation saw a clean benefit while this one
did not.

\subsection{Bonded-exclusion ZBL: implemented, not yet trained or tested}
\label{sec:q4-bonded-exclusion}

A dedicated literature review (2026-08-18,
\texttt{paper/literature\_review\_candidates.md} section 0) checked
NequIP's and MACE's actual \texttt{ZBL} source directly rather than relying
on their papers alone: neither implements a bonded-pair exclusion. Their
``cutoff normalized by the sum of covalent radii'' is a per-element-pair
\emph{adaptive cutoff}, not a topological exclusion --- it applies
identically whether a given pair happens to be bonded or not, so it would
not have resolved the O--H-bond-vs-anomaly-floor overlap diagnosed above
either. MACE-OFF, MACE's own flagship model for covalent organic chemistry,
simply does not enable ZBL at all. The only real precedent found is
classical MD: LAMMPS/CHARMM/AMBER all zero out nonbonded interactions for
1-2 (bonded) and 1-3 (angle) pairs via a topology-derived exclusion list,
specifically to avoid double-counting energy already described by
bonded/angle terms --- the direct, 60-year-old analog of what's needed here.
No paper was found reporting our specific failure mode (ZBL backfiring on a
bonded system) --- this appears to be a genuinely underexplored gap, not a
previously-solved or previously-warned-about problem.

\texttt{src/molecular\_zbl.py}'s \texttt{MolecularZBL} (subclasses
\texttt{torchmdnet.priors.zbl.ZBL}) implements this by excluding
same-\emph{molecule} atom pairs rather than literally-bonded pairs ---
exactly equivalent to the CHARMM/AMBER 1-2+1-3 convention for this specific
system, since a water molecule has only 3 atoms (O, H, H): every possible
pair within one molecule is either the O--H bond (1-2) or the H\ldots H
angle pair (1-3), so ``same molecule'' and ``1-2-or-1-3 bonded'' are exactly
the same set of pairs here --- no generic bond-graph distance-2 traversal
needed. The exclusion reuses the identical \texttt{local\_molecule\_ids}
tensor already trusted for the per-fragment momentum loss, not a second,
separately-computed grouping. Because torchmd-net's \texttt{create\_prior\_
models} resolves a prior by name (\texttt{getattr(priors, name)}), a custom
class needs registering into \texttt{torchmdnet.priors}'s own namespace at
\emph{every} site that might reconstruct the model from a checkpoint, not
just at initial training-time construction --- handled here by an
idempotent \texttt{register\_molecular\_zbl\_prior()} call, both in
\texttt{train\_waterbox.py} and inside
\texttt{evaluate\_waterbox.py}'s \texttt{load\_waterbox\_checkpoint} (the
one shared reload path also used by \texttt{waterbox\_ase.py}'s rollout
calculator and \texttt{analyze\_force\_decomposition.py}).

Wired through \texttt{train\_waterbox.py --use-zbl-prior --zbl-bonded-
exclusion}, \texttt{run\_waterbox\_study.py} (same flags, writing to a
further-separate \texttt{checkpoints/waterbox\_study\_zbl\_bonded/} root so
the stock-ZBL negative-result checkpoints are never touched), and
\texttt{run\_rollout\_study.py} (same flags). The masking logic (mapping a
batched edge index to its local, within-one-system atom index via modulo,
then checking molecule-id equality) was synthetically unit-tested locally;
the class itself subclasses torchmd-net's \texttt{ZBL} directly, so it
cannot even be imported on this Windows checkout.
\textbf{Smoke test has run; looked alarming, wasn't.} The single smoke-test
cell (\texttt{water\_absolute} seed 0, 2 epochs) showed
\texttt{energy\_mae=16060\,eV} --- superficially close to the stock-ZBL
seed-0 anomaly's 16404\,eV. Checked against that same stock-ZBL seed-0
checkpoint's own training history: its epoch-0 (before any real training)
\texttt{val\_y\_l1\_loss} was 25732\,eV, the same order of magnitude ---
confirming a barely-trained checkpoint (raw $\sim{-}$30{,}000\,eV total
energy, not a small residual, per \texttt{remove\_ref\_energy=False}) is
expected to show energy error in the thousands of eV regardless of ZBL
variant, and this is not a repeat of the batch-of-1 fragility (which only
appeared after full 48-epoch training). \texttt{force\_mae=3.0\,eV/\AA} is
consistent with this project's own established finding that force converges
faster than energy. The smoke test's only real job --- confirm the pipeline
runs without an exception --- is satisfied.

\textbf{Full retrain has run --- and it reproduced the exact same
batch-of-1 anomaly, in the exact same cell, confirming it is not a
cutoff-overlap artifact after all.} \texttt{results/waterbox\_study\_zbl\_
bonded/raw\_results.csv}: all 12 cells trained for the full 48 epochs
(\texttt{--force-retrain} correctly overwrote the smoke test's 2-epoch
checkpoint), and 11 of 12 land in the same healthy range as the no-ZBL and
stock-ZBL baselines (energy MAE 2.4--6.0\,eV). The sole exception is again
\texttt{water\_absolute} seed 0: \texttt{energy\_mae=16059.9\,eV},
essentially the same magnitude as stock ZBL's seed-0 anomaly
(16404\,eV, Section~\ref{sec:q4}). Checked every cell's \texttt{config.json}
the same way as before: that checkpoint's own
\texttt{post\_test.test\_y\_l1\_loss} (Lightning's \texttt{trainer.test()},
batch size 2) is 3.85\,eV --- indistinguishable from every other seed's
healthy value (3.60--6.00\,eV across all other cells, each closely matching
its own \texttt{raw\_results.csv} entry). Only this one cell diverges
between the two eval paths.

This settles the question left open in Section~\ref{sec:q4}: the batch-of-1
evaluation fragility is a property of \texttt{water\_absolute} seed 0's
particular learned weights, triggered by the presence of \emph{any}
ZBL-family prior term, not specifically by the covalent-bond-overlap problem
bonded exclusion was built to fix. Bonded exclusion neither improved nor
worsened it. Seed 1 remains verified safe under both eval paths with this
variant too (\texttt{post\_test.test\_y\_l1\_loss}=3.60\,eV vs.\
\texttt{raw\_results.csv} energy\_mae=3.598\,eV, a near-exact match) ---
used for the rollout comparison below, consistent with the stock-ZBL choice
in Section~\ref{sec:q4-negative-result}.

Excluding seed 0's outlier, bonded-exclusion \texttt{water\_absolute}
averages energy/force MAE of roughly 3.4\,eV / 0.67\,eV/\AA\ across seeds
1--5 --- statistically indistinguishable from both the no-ZBL baseline
(5.941$\pm$5.3 / 0.964$\pm$0.78) and stock ZBL, reconfirming that static
accuracy does not discriminate between prior variants here. One real
difference did show up: \texttt{water\_absolute}'s mean per-molecule
momentum violation fell from stock ZBL's 132.1 to 64.1 (close to the
no-ZBL baseline's 60.9) --- consistent with bonded exclusion removing the
spurious constant force stock ZBL injected onto every O--H bond.

\textbf{Rollout comparison at seed 1 has run (dt=0.1, both axes) - bonded
exclusion is a large, real improvement over stock ZBL, but a substantial
gap to the true $\sim$300\,K starting temperature remains.}
\texttt{results/waterbox\_rollout\_study\_zbl\_bonded\_seed1}: plateau
temperatures of 821$\pm$15\,K (\texttt{water\_absolute}) / 902$\pm$20\,K
(\texttt{water\_absolute+momentum}) on the velocity axis, and
697$\pm$150\,K / 735$\pm$160\,K on the config axis --- an order of
magnitude better than stock ZBL's millions-of-K numerical explosion at the
same dt, and also better than the no-ZBL baseline (1227--1342\,K velocity
axis, seed 1). Bonded exclusion is confirmed as a genuine fix for the
double-counting problem, not just a less-catastrophic failure mode.

\subsection{A second, independent confound in this same comparison: checkpoint selection}
\label{sec:q4-anneal-confound}

A second issue, unrelated to ZBL, surfaced while tracing why the momentum
condition ran hotter than absolute above --- the opposite of the no-ZBL
baseline's direction. Checking which epoch \texttt{val\_checkpoint\_score}
actually selected as ``best'' for every one of the 12
\texttt{waterbox\_study\_zbl\_bonded} cells (6 seeds $\times$ 2 conditions)
found it ranges from epoch 17 to 47 out of a 50-epoch budget, and is
essentially uncorrelated with anything but which epoch's noisy post-anneal
validation score happened to dip lowest within the epochs remaining. The
seed-1 pair used for the rollout above is a clean example of the problem:
\texttt{water\_absolute} selected epoch 46 (16 epochs of post-anneal
fine-tuning), while \texttt{water\_absolute+momentum} selected epoch 29 ---
\emph{one epoch before \texttt{anneal\_epoch=30} even fires}, i.e. this
``momentum'' checkpoint never received any energy-focused fine-tuning at
all. Comparing the two is comparing a converged model against an
essentially pre-anneal one, not a fair test of the momentum loss term.

Fixed via \texttt{eligible\_epoch\_start} (\texttt{train\_waterbox.py}'s
\texttt{WaterLNNP}): every epoch before this index is logged with
\texttt{val\_checkpoint\_score}=$+\infty$, so \texttt{ModelCheckpoint}'s
plain running-minimum tracker can mathematically never select one of them;
the true (unmasked) score is logged separately as
\texttt{val\_checkpoint\_score\_raw} so \texttt{EarlyStopping} keeps
watching genuine convergence rather than being fooled by the $+\infty$
floor into stopping before the eligible window ever opens.
\texttt{run\_waterbox\_study.py --extended-anneal} pairs this with
\texttt{EXTENDED\_NUM\_EPOCHS=70}/\texttt{EXTENDED\_ELIGIBLE\_EPOCH\_START=50}
--- unchanged \texttt{anneal\_epoch=30}, but a mandatory 20-epoch post-anneal
burn-in (epochs 30--49) before any epoch becomes selectable, leaving a real
20-epoch eligible window (50--69) instead of the original schedule's
effectively-0-to-20-epoch one. Writes to a further-separate
\texttt{checkpoints/waterbox\_study\_..\_ext70/} root. Verified working as
designed on all 4 retrained cells (seeds 0/1 $\times$ both conditions):
every selected epoch (68, 64, 68, 66) falls inside the 50--69 window, and
every cell's \texttt{post\_test.test\_y\_l1\_loss} (Lightning's own
\texttt{trainer.test()}) again closely matches its own
\texttt{evaluate\_waterbox\_checkpoint} number --- no repeat of the
seed-0 batch-of-1 anomaly (Section~\ref{sec:q4}) at any of these new
epochs, suggesting that fragility tracks specific learned weights at
whichever epoch got selected rather than being a fixed property of
seed 0 + \texttt{water\_absolute}.

\subsection{The static-accuracy/dynamics gap has a specific, measured mechanism here}
\label{sec:q4-epot-gap}

Tracing the still-elevated plateau temperature (both conditions, both ZBL
variants, before or after the anneal fix) through
\texttt{energy\_history.csv} directly: \texttt{etot\_ev} is essentially
perfectly conserved throughout every rollout (drift $\sim$0.001--0.02\,eV
out of a $\sim$30{,}000\,eV system) --- ruling out an integrator or
non-conservative-force bug. What actually happens: \texttt{epot\_ev} drops
sharply in the first $\sim$100\,fs while \texttt{ekin\_ev} rises by the same
amount, i.e. the model assigns the real DFT starting configuration
measurably more potential energy than wherever its own nearby learned
minimum sits, and energy-conserving dynamics converts that one-time gap
into heat. The gap is a \emph{static, per-configuration} property, not
dynamical noise: across 5 velocity-seed replicates sharing one fixed
starting geometry it is nearly identical every time (seed-1 ext70, velocity
axis: 11.6$\pm$0.37\,eV absolute, 9.96$\pm$0.24\,eV momentum) --- initial
velocity doesn't matter, the model has already ``decided'' this geometry
has too much potential energy before any dynamics run. Across different
starting geometries the gap varies substantially and predicts the outcome:
the config-axis batch (seed 1, ext70) shows a clean monotonic relationship
between per-geometry gap and final temperature, independently for both
conditions (cfg2: 2.9/2.7\,eV $\to$ 439/425\,K; cfg1: 7.9/7.0\,eV $\to$
637/632\,K; cfg5: 9.0/8.3\,eV $\to$ 635/688\,K; cfg3: 9.1/8.3\,eV $\to$
671/625\,K; cfg4: 13.3/11.9\,eV $\to$ 842/763\,K) --- the smallest-gap
geometry barely heats at all, the largest-gap geometry heats the most, with
every intermediate config falling in between in the same order.

This is a specific, previously-characterized MLIP failure mode, not a
symptom unique to this pipeline. \citet{fu2022forces} (already this
project's central citation) established the core disconnect: static
force/energy test-set accuracy is an average over many held-out
configurations and does not predict simulation stability, because NVE
dynamics has no mechanism to average anything out --- it deterministically
converts whatever residual error exists at the \emph{one} configuration
a rollout launches from into heat, regardless of the model's accuracy
elsewhere. \citet{stocker2022robust} directly quantify this: a GNN
potential with test force MAE of only 0.035\,eV/\AA\ still accumulated an
absolute-energy error exceeding 10\,eV during a 300\,K rollout --- the same
order of magnitude measured here --- which they describe as the PES having
``holes'': localized regions the trajectory reaches where the true error is
far above the reported test-set average, invisible to static evaluation
and exposed only by running real dynamics. Most directly relevant:
\citet{raja2025stable} study almost exactly this project's system (a
64-water-molecule box, the same underlying benchmark \citet{fu2022forces}
introduced) and hit the same disconnect --- a model with good validation
accuracy that nonetheless produces localized instabilities invisible to
standard metrics (median stable time only 312\,ps at baseline). Their fix,
StABlE training, is not another static-data or architecture change but
differentiable fine-tuning directly against the model's own simulated
trajectories, supervised by a target observable (they used mean O--H bond
length); it took their baseline to 1\,ns stable and recovered correct
diffusivity and radial distribution functions. The standard alternative in
the wider MLIP literature is DP-GEN-style concurrent/active learning
\citep{zhang2020dpgen}: retraining with configurations sampled from the
model's own rollouts, specifically because a model trained only on
pre-existing AIMD snapshots has never seen the configurations its own
dynamics will actually explore, even ones close to equilibrium.

\textbf{Decisive check has run: static ENERGY error does not predict rollout
heating, but static FORCE error does.} \texttt{evaluate\_waterbox.py
--per-sample} on both ext70 seed-1 checkpoints, restricted to the 6 configs
actually used in the rollouts above (\texttt{test\_config\_index} 0--5).
Spearman correlation between each config's static per-sample error and its
rollout-derived \texttt{epot} gap: \texttt{energy\_error} gives
$\rho=0.20$ (absolute) / $\rho=-0.14$ (momentum) --- essentially no
relationship, slightly negative for momentum. \texttt{force\_error} gives
$\rho=0.83$ (absolute) / $\rho=0.71$ (momentum) --- a real relationship,
and the two extremes match exactly and independently in both conditions:
config index 2 has the lowest \texttt{force\_error} of all 6
(0.389/0.360\,eV/\AA) and the coldest rollout in both (439/425\,K, smallest
gap); config index 4 has the highest \texttt{force\_error}
(0.728/0.551\,eV/\AA) and the hottest rollout in both (842/763\,K, largest
gap) --- despite being near-\emph{best} on static energy accuracy
(0.73/0.21\,eV, among the lowest of the 6), ruling out ``the model gets the
energy value wrong here'' as the mechanism for this specific configuration.

This sharpens, rather than overturns, the mechanism proposed above: NVE
dynamics is driven by forces (the gradient of the PES), not energy values
--- a model can predict the right total energy at a configuration while
still having the wrong gradient there, which is exactly what would slide
the system toward the model's own nearby basin and convert that mismatch
into heat, independent of whether the energy value at the starting point
happened to be accurate. This is a direct, project-specific confirmation of
what \citet{fu2022forces}'s own title argues in general (\emph{forces} are
not enough) --- here, force error specifically, not energy error, is the
static quantity that actually predicts rollout stability.

\subsection{Path forward: stability-aware fine-tuning via StABlE}
\label{sec:q4-stable-path}

\textbf{Where this leaves Q4.} The two bugs fixed so far this session (ZBL
bonded-pair double-counting, Section~\ref{sec:q4-negative-result}--\ref{sec:q4-bonded-exclusion};
checkpoint-selection landing pre-anneal,
Section~\ref{sec:q4-anneal-confound}) were both real and both produced
real, measured improvement --- but neither was ever a fix for the
mechanism identified in Section~\ref{sec:q4-epot-gap}, and neither closes
it: both conditions still plateau at 600--900\,K against a $\sim$300\,K
start at seed 1, the best (least confounded) data currently available. The
mechanism is now specific and quantitatively confirmed: per-configuration
\emph{force} error (not energy error) at the rollout's starting geometry
predicts how much a given rollout heats ($\rho=0.83$/$0.71$,
Section~\ref{sec:q4-epot-gap}). Two candidate fixes were considered
(Section~\ref{sec:q2}): DP-GEN-style active learning on the model's own
rollout trajectories, and StABlE-style stability-aware fine-tuning
\citep{raja2025stable}. Active learning was set aside for a concrete,
practical reason rather than a scientific one: it requires a DFT oracle to
label newly-explored configurations at the reference level of theory
(revPBE0-D3), and this pipeline has none --- \texttt{waterbox\_data.py}
only ever consumes Cheng et al.'s pre-computed static dataset, with no
CP2K/VASP labeling step anywhere in the codebase. StABlE does not have
this dependency, which is the deciding factor for pursuing it first.

\textbf{How StABlE actually works (verified against the paper directly, not
just its abstract).} The method alternates two phases, and --- critically
for what infrastructure this needs --- never backpropagates through the MD
trajectory itself. In the \emph{simulation phase}, $R$ replicas run
thermostatted (Langevin) dynamics from training-set initial conditions;
every $t$ steps each replica is checked against a stability criterion
(e.g.\ bond-length deviation), and replicas that fail are rewound. In the
\emph{learning phase}, replicas that stayed stable are simulated for $t$
further steps, subsampled every $S$-th frame for decorrelation, and used to
estimate $\partial_\theta \mathbb{E}_{P_\theta}[g(\Gamma)]$ for a target
observable $g$ (they use the RDF and mean O--H bond length) via what they
call the differentiable Boltzmann estimator:
\[
\mathcal{E} = \frac{N}{k_BT(N-1)}\Big[\hat{\mathbb{E}}[g(\Gamma)]\,
\hat{\mathbb{E}}[\nabla_\theta U_\theta(\Gamma)]^\top -
\hat{\mathbb{E}}\big[g(\Gamma)\,\nabla_\theta U_\theta(\Gamma)^\top\big]\Big],
\]
a covariance-based, unbiased estimator built on thermodynamic perturbation
theory \citep{zwanzig1954} --- the same reweighting idea underlying
\citet{thaler2021difftre}'s DiffTRe, which StABlE builds on directly and
which demonstrated the same strategy on a coarse-grained water model. The
key practical consequence: $\mathbb{E}_{P_\theta}[g(\Gamma)]$'s dependence
on $\theta$ only enters through the Boltzmann distribution itself, not
through which integrator produced the samples --- so the samples $\Gamma$
can come from \emph{any} sampler, including a non-differentiable one, and
the only per-sample quantity requiring autograd is
$\nabla_\theta U_\theta(\Gamma)$, an ordinary single-configuration gradient
of the kind \texttt{evaluate\_waterbox.py} already computes for every test
config. A differentiable, PyTorch-native integrator --- which this project
does not have, since \texttt{waterbox\_ase.py} wraps the model as a
numpy-based ASE \texttt{Calculator} --- is \emph{not} required. The total
loss is $\mathcal{L}_{\text{StABlE}} = \mathcal{L}_{\text{obs}} +
\lambda\,\mathcal{L}_{\text{QM}}$, where $\mathcal{L}_{\text{QM}}$ is the
ordinary supervised energy/force loss over the original fixed training set
(unchanged, purely a regularizer against drifting away from DFT accuracy)
and $\mathcal{L}_{\text{obs}}$ is the Boltzmann-estimator loss above ---
notably, $\mathcal{L}_{\text{QM}}$ is never computed on newly-explored
structures, only the original data, so this method (unlike active
learning) needs no new labels at any point.

\textbf{Concrete plan, mapped onto infrastructure already in this repo:}
\begin{itemize}
  \item \textbf{Target observable}: the O--O/O--H/H--H radial distribution
  function, since \texttt{rollout\_waterbox\_ase.py} already computes
  exactly this against a real-DFT-configuration reference for every
  rollout --- the target and the computation both already exist, reused
  rather than rebuilt. Mean O--H bond length (via
  \texttt{structural\_metrics.py}'s existing bond-length utilities) is a
  cheaper secondary/sanity observable, matching one of the two used in
  \citet{raja2025stable} directly.
  \item \textbf{Instability/rewind criterion}: reuse
  \texttt{diagnose\_short\_range\_collapse.py}'s already-calibrated
  \texttt{heating\_onset\_frame}/\texttt{element\_pair\_intermolecular\_
  distances} utilities (Section~\ref{sec:q4}) as the simulation-phase
  stability check, rather than inventing a new criterion from scratch.
  \item \textbf{Sampling engine}: since gradients need not flow through the
  trajectory (see above), the existing ASE-based rollout infrastructure
  (\texttt{waterbox\_ase.py}'s \texttt{TensorNetCalculator},
  \texttt{rollout\_waterbox\_ase.py}) can generate the simulation-phase
  samples directly, with a Langevin thermostat added for the training-time
  sampling dynamics specifically (kept distinct from the NVE protocol used
  for stability \emph{evaluation} --- the two serve different purposes and
  should not be conflated; the existing NVE rollout study remains the
  correct tool for measuring whether a fine-tuned checkpoint actually
  improved).
  \item \textbf{Which checkpoints}: fine-tune both ext70 seed-1 checkpoints
  (\texttt{water\_absolute} and \texttt{water\_absolute+momentum}), since
  the momentum-vs-absolute question (Section~\ref{sec:q1}) is still open
  and orthogonal to this fix --- doing both and re-running the existing
  rollout comparison answers ``did this close the gap'' and ``does momentum
  still look better once both are stability-fine-tuned'' simultaneously,
  rather than confounding them.
\end{itemize}

\textbf{Honest scope of the remaining work.} This is a real implementation
effort, not a flag flip: a parallel-replica simulation-phase driver with
rewind logic, the covariance-form estimator itself, and a Langevin
integrator all need building (or sourcing) --- none of this project's
existing training code performs any of them. Before committing that effort,
the RDF/bond-length target values should be verified against literature
(Cheng et al.\ 2019's own reported RDF, or \citet{thaler2021difftre}'s
water numbers) the same way \texttt{verify\_zbl\_units.py} empirically
checked ZBL's cutoff rather than trusting a default --- this project has
already been burned once by an unverified hyperparameter choice
(Section~\ref{sec:q4-negative-result}) and should not repeat it here with
an unverified target observable. The force-error correlation motivating
this whole path (Section~\ref{sec:q4-epot-gap}) was measured on only 6
configurations; that does not block starting this work (StABlE targets the
general mechanism via a global structural observable, not that specific
correlation), but a larger-sample confirmation of the correlation remains
worth doing in parallel, as a independent check that the diagnosis this
plan rests on is solid.

\subsubsection{Implementation plan}
\label{sec:q4-stable-plan}

New modules, reusing existing infrastructure wherever possible rather than
rebuilding it:

\begin{description}
  \item[\texttt{src/boltzmann\_estimator.py}] the core, highest-risk piece.
  Implements $\mathcal{L}_{\text{obs}}$ (Section~\ref{sec:q4-stable-path})
  as a constructed pseudo-loss rather than a hand-coded backward pass: given
  a batch of already-sampled, detached configurations, re-run the TensorNet
  forward pass on each with gradients enabled to get $U_\theta(\Gamma_i)$,
  compute the target observable $g(\Gamma_i)$ from the frozen positions (no
  gradient needed there --- only $U_\theta$ does), and assemble
  $\mathcal{L}_{\text{obs}}=\frac{2}{(N-1)k_BT}\big[\overline{g}\cdot
  \overline{U_\theta}-\overline{g\cdot U_\theta}\big]$ as a literal torch
  expression; calling \texttt{.backward()} on it reproduces the estimator
  automatically via ordinary autograd, avoiding a hand-implemented backward
  pass for the covariance formula.
  \item[\texttt{src/waterbox\_langevin.py}] the sampler. Reuses
  \texttt{waterbox\_ase.py}'s \texttt{TensorNetCalculator} unchanged, swaps
  \texttt{VelocityVerlet} (used for NVE evaluation) for ASE's
  \texttt{ase.md.langevin.Langevin} (already an ASE dependency, nothing new)
  for thermostatted training-time sampling. Adds a snapshot/rewind
  mechanism: periodically stash each replica's positions and velocities,
  run \texttt{diagnose\_short\_range\_collapse.py}'s already-calibrated
  distance checks as an attached callback (the same
  \texttt{.attach(callback, interval=N)} pattern
  \texttt{rollout\_waterbox\_ase.py} already uses), and on a stability trip
  reload the last good snapshot instead of continuing.
  \item[\texttt{src/train\_waterbox\_stable.py}] the fine-tuning
  entrypoint. Loads a checkpoint via \texttt{evaluate\_waterbox.py}'s
  existing \texttt{load\_waterbox\_checkpoint} (so ZBL/bonded-exclusion
  prior registration is not duplicated), runs the alternating
  simulate$\to$rewind$\to$learn loop, logs via a
  \texttt{training\_history.py}-style callback, saves to a new checkpoint
  root.
  \item[\texttt{src/run\_stable\_finetune.py}] orchestration, mirroring
  \texttt{run\_waterbox\_study.py}'s CLI conventions
  (\texttt{--train-seed}, \texttt{--use-zbl-prior --zbl-bonded-exclusion
  --extended-anneal} to locate the source ext70 checkpoint) and its
  separate-root discipline: writes to
  \texttt{checkpoints/waterbox\_study\_zbl\_bonded\_ext70\_stable/}, never
  touching the checkpoints being fine-tuned from.
\end{description}

No new evaluation code: \texttt{run\_rollout\_study.py} gets pointed at the
new checkpoint root by extending \texttt{\_roots()} with one more suffix
flag --- the existing NVE rollout comparison is already the right
instrument, and reusing it keeps the before/after numbers directly
comparable to every result already in this paper.

\textbf{Build order, staged so the highest-risk piece is proven in
isolation before it touches a real model:}
\begin{enumerate}
  \item \texttt{boltzmann\_estimator.py} alone, unit-tested against
  synthetic data with a hand-computable expected gradient --- this is the
  piece most likely to carry a silent sign or normalization error that
  would otherwise run without raising an exception and produce
  plausible-looking but wrong gradients, the same failure mode a bug would
  take here as \texttt{verify\_zbl\_units.py} was built to catch for ZBL.
  Fully testable on this project's Windows checkout, no \texttt{torchmdnet}
  required.
  \item Sampler smoke test on the training box: a few replicas, few steps,
  confirm no crash, and --- critically --- artificially lower the
  stability threshold to force an early trip and confirm the rewind
  mechanism actually engages. A rewind that silently never triggers would
  let the rest of the pipeline appear to work while doing nothing.
  \item Wire the full alternating loop; smoke-test with 2--3 outer
  iterations on one seed/condition, confirming checkpoints save and the
  loss doesn't diverge.
  \item Real run on seed 1 (both conditions --- the cleanest existing
  before/after data, Section~\ref{sec:q1-rollout-update}), evaluated with
  the unmodified \texttt{run\_rollout\_study.py} against the 772/706\,K
  ext70 baseline.
  \item Seed 0, once seed 1 looks sane --- the more informative test given
  its much larger starting gap (1682/1072\,K), and where the
  margin-scales-with-severity hypothesis from
  Section~\ref{sec:q1-rollout-update} actually gets checked.
\end{enumerate}

\textbf{Hyperparameters requiring empirical verification, not a default}:
the RDF target checked against Cheng et al.'s own reported values (or
\citet{thaler2021difftre}'s) before trusting the fine-tuning is being
pulled toward a correct target; $\lambda$ started conservative and checked
via \texttt{evaluate\_waterbox.py --per-sample} after fine-tuning to
confirm static accuracy hasn't quietly regressed; the fine-tuning learning
rate set well below the original $10^{-4}$ (this is fine-tuning a
converged model, not training from scratch) and chosen empirically rather
than inherited by default.

\subsubsection{Step 1 implemented and verified: the Boltzmann-estimator pseudo-loss}
\label{sec:q4-stable-step1}

\texttt{src/boltzmann\_estimator.py} implements $\mathcal{L}_{\text{obs}}$,
the highest-risk piece of the plan
(Section~\ref{sec:q4-stable-plan}). Documented here in full since a wrong
sign or normalization in this one piece would silently corrupt every
downstream fine-tuning run without raising an error.

\textbf{What $\mathcal{L}_{\text{obs}}$ actually is, derived from scratch.}
For a Boltzmann distribution $P_\theta(\Gamma) \propto
e^{-U_\theta(\Gamma)/kT}$, differentiating $\mathbb{E}_\theta[g(\Gamma)] =
\int g(\Gamma) P_\theta(\Gamma)\,d\Gamma$ with respect to $\theta$ and using
$\partial_\theta \log Z_\theta = -\mathbb{E}_\theta[\nabla_\theta
U_\theta]/kT$ (the standard partition-function identity) gives, after
substitution,
\[
\frac{\partial}{\partial\theta}\mathbb{E}_\theta[g(\Gamma)] =
-\frac{1}{kT}\,\mathrm{Cov}_\theta\big(g(\Gamma),\,\nabla_\theta U_\theta(\Gamma)\big).
\]
Replacing the exact (population) covariance with its unbiased finite-sample
estimator ($N/(N{-}1)$ Bessel correction) reproduces
\citet{raja2025stable}'s reported estimator exactly --- an independent
confirmation of their formula rather than a copied one. The full
observable-matching loss then follows from one more application of the
chain rule on $\mathcal{L}_{\text{obs}}=\tfrac12(\mathbb{E}_\theta[g]-g_{\text{target}})^2$:
\[
\frac{\partial\mathcal{L}_{\text{obs}}}{\partial\theta} =
\big(\hat{\mathbb{E}}[g]-g_{\text{target}}\big)\cdot
\frac{N}{kT(N-1)}\Big[\hat{\mathbb{E}}[g]\,\hat{\mathbb{E}}[\nabla_\theta U_\theta]
-\hat{\mathbb{E}}[g\,\nabla_\theta U_\theta]\Big].
\]
The module implements this not by hand-computing the right-hand side, but
by constructing a scalar \emph{pseudo-loss} --- a literal torch expression
built from per-sample energies $U_\theta(\Gamma_i)$ evaluated with
autograd enabled (an ordinary forward pass at each already-\emph{sampled,
frozen} configuration, ``already-sampled'' meaning the observable $g$ is
always \texttt{.detach()}-ed and no gradient is ever required through
however $\Gamma$ was generated) --- such that calling \texttt{.backward()}
on it reproduces the formula above exactly, via ordinary autograd rather
than a hand-coded backward pass.

\textbf{A literature check surfaced a real implementation choice, not just
a formula to copy.} Fetching \citet{thaler2021difftre}'s DiffTRe paper
directly (the reweighting method \citeauthor{raja2025stable}'s estimator is
built on) before implementing this found that DiffTRe's production
implementation does \emph{not} resample fresh configurations every gradient
step --- it reuses one reference trajectory across many steps via
importance-sampling weights $w_i \propto e^{-(U_\theta(\Gamma_i)-U_{\hat\theta}(\Gamma_i))/kT}$
for computational efficiency, tracking an effective sample size
$N_{\text{eff}}\approx e^{-\sum_i w_i \ln w_i}$ and re-simulating only once
$N_{\text{eff}}$ drops too low (importance-weight degeneracy). StABlE's own
description (Section~\ref{sec:q4-stable-path}) --- simulate a short window,
take one gradient step, reset replicas, repeat --- is instead genuinely
\emph{on-policy}: every gradient step samples fresh from the
\emph{current} $\theta$'s own distribution, which is exactly what the
derivation above assumes and needs no importance weights at all. This
module deliberately implements the on-policy form, not DiffTRe's
importance-reweighted variant --- copying the latter in without its
accompanying $N_{\text{eff}}$ safeguard would have introduced a real,
documented failure mode this design avoids by construction. Also checked
against DiffTRe: per-bin loss weighting for a vector-valued (histogram)
observable like the RDF. DiffTRe uses a plain, unweighted mean-squared
error across bins for a single observable type (their inter-observable
weights $\gamma$ balance RDF against ADF against pressure, not bins within
one RDF) --- this module matches that convention rather than inventing
per-bin variance weighting that has no precedent in the literature it is
built on.

\textbf{Verification (\texttt{src/verify\_boltzmann\_estimator.py}), all 6
checks passing:} a synthetic discrete ``toy Boltzmann system'' of 300
states with known linear-in-$\theta$ energies gives an \emph{exact}
(sampling-noise-free) ground truth for $\mathbb{E}_\theta[g]$ and its
finite-difference derivative, independent of anything in
\texttt{boltzmann\_estimator.py}. States are then Monte Carlo sampled
on-policy from the true $P_{\theta_0}$ distribution and fed through the
real module. Results: a single $N{=}20{,}000$ draw already matches the
ground truth to 2.6\% relative error; averaged over 400 independent trials
at $N{=}500$, the mean recovered gradient matches the exact-expectation
target to 0.056\% --- strong evidence against a subtle bias a single lucky
draw could otherwise hide. The construction was also checked to generalize
correctly to a 3-parameter $\theta$ (matching finite-difference partial
derivatives per component, real TensorNet fine-tuning will use thousands of
parameters simultaneously) and to a genuinely multi-bin observable
(cross-checked two ways: a $B{=}1$ vector path reproduces the scalar path
bit-for-bit, and a real $B{=}3$ case matches its own finite-difference
ground truth to 0.86\%). Edge cases (single-sample batches, a detached $U$
that would otherwise silently zero out the gradient, mismatched sample
counts) all raise clear errors rather than failing silently.

\textbf{Status}: this is step 1 of the build order in
Section~\ref{sec:q4-stable-plan} --- the pseudo-loss construction is
implemented and independently verified. \texttt{train\_waterbox\_stable.py}
(the alternating fine-tuning loop tying everything together) is not yet
built.

\subsubsection{Step 2 (partially verified): the sampler and rewind mechanism}
\label{sec:q4-stable-step2}

\texttt{src/waterbox\_langevin.py} implements the simulation/rewind half of
the plan, split the same way \texttt{diagnose\_short\_range\_collapse.py}
already is (that module's own docstring): a pure-numpy state machine
(\texttt{Snapshot}, \texttt{ReplicaState}) testable on this checkout without
\texttt{ase}/\texttt{torchmdnet}, and an ASE-driven half
(\texttt{run\_stable\_sampling\_phase}) that is not.

\textbf{Reused, not duplicated}: the live per-step stability check calls two
new small functions added to \texttt{diagnose\_short\_range\_collapse.py}
(\texttt{frame\_min\_distances}, \texttt{frame\_violates\_floors}), factored
out of that module's already-validated \texttt{analyze\_trajectory} so a
running simulation and the existing post-hoc scan
(Section~\ref{sec:q4}) share one distance/exclusion implementation rather
than a second, separately-trusted copy of ``is this frame anomalous.''

\textbf{Two things checked against ASE's actual documented behavior before
writing this, rather than assumed}:
\begin{itemize}
  \item \texttt{ase.md.langevin.Langevin}'s \texttt{friction} parameter is
  documented to be supplied already divided by \texttt{ase.units.fs}
  (their own example: \texttt{0.01 / units.fs} for a physical
  0.01\,fs$^{-1}$ rate) --- the \emph{opposite} direction from
  \texttt{timestep=dt*units.fs}, already used elsewhere in this project
  (\texttt{rollout\_waterbox\_ase.py}). Getting this backwards would not
  crash --- it would silently make the thermostat couple roughly 10$\times$
  too weakly or too strongly (\texttt{units.fs}$\approx0.0982$) --- exactly
  the failure signature of the Bohr/Hartree unit mixup already caught once
  in this project (\texttt{waterbox\_data.py}). The module's own
  \texttt{friction\_fs\_inv} parameter takes the physical value and does
  the conversion once, internally, at the single call site.
  \item ASE's documentation does not state whether a \texttt{Langevin}
  object tolerates an externally-mutated \texttt{Atoms} state (a hard
  \texttt{set\_positions}/\texttt{set\_velocities} reset, exactly what a
  rewind does) between \texttt{.run()} calls without carrying stale
  internal bookkeeping. Rather than trust that undocumented behavior, a
  fresh \texttt{Langevin} object is constructed after every rewind ---
  cheap relative to the simulation itself, and removes the question
  entirely instead of resolving it optimistically.
\end{itemize}
Also reused rather than reintroduced: explicit \texttt{rng=np.random.
RandomState(seed)} for both Langevin's own random-force draws and the
post-rewind velocity redraw --- the identical discipline this project
already learned the hard way for
\texttt{MaxwellBoltzmannDistribution} (Lessons, \texttt{CLAUDE.md}: it
silently draws from whatever state numpy's \emph{global} RNG happens to be
in unless told otherwise).

\textbf{Verified (\texttt{src/verify\_waterbox\_langevin.py}), 7/7 checks
passing}: the \texttt{ReplicaState} state machine --- recording stable
frames, rewinding to the last known-good snapshot (not the original one,
and correctly discarding the invalidated window), stride-based
subsampling with correct edge-case behavior (empty buffer, stride larger
than the collected window), and a full realistic
simulate$\to$rewind$\to$collect$\to$subsample sequence end to end.

\textbf{Not yet verified}: \texttt{run\_stable\_sampling\_phase} itself
(needs a real \texttt{ase}+\texttt{torchmdnet} install) has not been run
anywhere. Before trusting it, run the thermostat-only sanity check this
module ships specifically to catch a friction-unit mistake directly rather
than downstream in an opaque fine-tuning run:
\begin{quote}
\texttt{python src/waterbox\_langevin.py --ckpt
checkpoints/waterbox\_study\_zbl\_bonded\_ext70/water\_absolute/seed1/
best\_model.ckpt --smoke-test}
\end{quote}
This runs one replica's Langevin dynamics with no stability check or
rewind logic involved, and logs temperature over the run. A correctly
configured thermostat should show the temperature equilibrate toward
\texttt{--temperature-k} (default 300\,K) at a physically reasonable rate;
barely any coupling (temperature drifting the way plain NVE already does,
Section~\ref{sec:q4-epot-gap}) or wild oscillation/collapse would both
directly indicate the friction convention above needs revisiting before
anything is built on top of it.

\textbf{Thermostat smoke test: run, and passes.} 2000 steps (1\,ps) at
\texttt{friction\_fs\_inv=0.01} on \texttt{water\_absolute} seed 1's ext70
checkpoint: tail mean $320.1\pm13.1$\,K against a 300\,K target. Checked
against equipartition rather than eyeballed: for this 192-atom system
($N_{\text{dof}}=3N{-}3=573$), the expected relative fluctuation in
instantaneous temperature is $\sqrt{2/N_{\text{dof}}}\approx5.9\%$, i.e.\
$\approx17.7$\,K at 300\,K --- close to the observed 13.1\,K (a bit smaller,
consistent with logged samples not being fully decorrelated). A
wrong-direction or $\sim\!10\times$-off friction would look qualitatively
different: either near-total decoupling (unbounded drift, matching plain
NVE's own heating tendency) or fluctuations far smaller than equipartition
predicts (an overdamped, artificially rigid thermostat) --- neither
happened. The $+20$\,K mean offset is plausible rather than alarming: this
exact checkpoint is already known to have a real NVE heating tendency
(772\,K plateau, Section~\ref{sec:q1-rollout-update}), so a modest residual
pull above 300\,K over just 1\,ps of thermostatting, while friction is
still working against it, is consistent with the very problem StABlE
exists to fix, not a units bug.

\textbf{A second smoke test, \texttt{--smoke-test-rewind}, targets the
rewind mechanism itself} --- the \texttt{ReplicaState} logic is unit-tested
(Section~\ref{sec:q4-stable-step2} above), but that never exercises a real
\texttt{ase.md.langevin.Langevin} object or a real
\texttt{set\_positions}/\texttt{set\_velocities} reload. Forces every
stability check to trip by passing deliberately impossible floors (1000\,\AA,
far beyond any real distance in a $\sim$12\,\AA\ box) --- matching the
build order's own prescription (Section~\ref{sec:q4-stable-plan}:
``artificially lower the stability threshold to force an early trip'')
--- and checks three concrete signals: \texttt{n\_rewinds} matches the
expected check count exactly (every single check trips, by construction);
zero samples collected (since \texttt{last\_good} can never advance past
the initial snapshot when every check fails); and, most directly, the
replica's final positions land back on exactly the original starting
positions (max drift $<10^{-6}$\,\AA) --- confirming the reload actually
happens, not just that a counter increments.

\textbf{Run, and passes exactly}: $n_{\text{rewinds}}=5$ (matching the
predicted $100/20$ exactly, every check tripped as forced), 0 samples
collected, and max position drift $=0.000000$\,\AA\ --- the reload is
confirmed bit-exact, not merely approximately correct. One harmless-but-
worth-fixing side effect surfaced while running this:
\texttt{molecule\_group\_ids} (\texttt{diagnose\_short\_range\_collapse.py})
warned on a non-writable array, since \texttt{atoms.get\_cell().diagonal()}
returns a read-only numpy view and \texttt{torch.as\_tensor} warns rather
than copying automatically. Fixed at the source with an explicit
\texttt{.copy()} (checked empirically that \texttt{np.ascontiguousarray}
alone would not have been a general fix here, since it only forces a copy
incidentally when the input happens to be non-contiguous) --- one line,
zero behavior change, and this function is about to be called constantly
by the real fine-tuning loop, so worth silencing at the source rather than
downstream at every call site.

\textbf{Status}: both halves of step 2 are now verified --- the pure
\texttt{ReplicaState} logic (Section~\ref{sec:q4-stable-step2}) and the
real ASE-driven thermostat and rewind mechanism (this subsection).

\subsubsection{Step 3 (written, not yet run): the fine-tuning loop}
\label{sec:q4-stable-step3}

\texttt{src/train\_waterbox\_stable.py} ties \texttt{boltzmann\_estimator.py}
and \texttt{waterbox\_langevin.py} together: each outer iteration runs a
sampling phase across $R$ replicas, pools every replica's collected
Snapshots into one batch, computes each pooled sample's RDF ($g$, pure
geometry) and the model's predicted energy there ($U$, a fresh
grad-enabled forward pass), takes one $\mathcal{L}_{\text{obs}} +
\lambda\mathcal{L}_{\text{QM}}$ gradient step, and repeats. Explicitly a
\emph{fine-tuning} loop, not a from-scratch training loop --- it loads an
existing ext70 checkpoint and continues optimizing the same weights, never
reinitializing them.

Deliberately not batched into one multi-graph forward pass per gradient
step: each pooled sample's energy is computed with its own single-graph
forward call, the same pattern \texttt{evaluate\_waterbox.py} and
\texttt{waterbox\_ase.py} already use everywhere else in this project,
rather than a new, unverified multi-graph batching path whose correctness
can't be checked without a real \texttt{torchmdnet} install. The sample
count per gradient step here is small (tens, not thousands) --- a real but
modest cost, deferred as a future optimization once this simpler version
is confirmed correct.

\textbf{Three real issues caught by review before this ever reached the
training box} (this project's own established practice --- see
\texttt{verify\_zbl\_units.py}, the checkpoint-selection confound
Section~\ref{sec:q4-anneal-confound} --- of not trusting a first draft
without checking it):
\begin{itemize}
  \item \textbf{The shared-model bug.} The first draft constructed each
  replica's \texttt{TensorNetCalculator} directly from the checkpoint path
  (\texttt{TensorNetCalculator(ckpt, ...)}), which loads its own
  independent copy of the model on disk --- a completely different object
  from the one the optimizer updates. Every replica would have silently
  sampled forever from the original, un-fine-tuned weights, never
  reflecting any training progress: exactly the DiffTRe-without-
  reweighting mistake \texttt{boltzmann\_estimator.py}'s own docstring
  warns against reintroducing by accident (Section~\ref{sec:q4-stable-step1}).
  Fixed with a small, backward-compatible extension to
  \texttt{waterbox\_ase.py}'s \texttt{TensorNetCalculator}: an optional
  \texttt{model=} constructor argument that uses an already-loaded model
  object directly instead of loading a fresh one, so every replica and the
  optimizer share the exact same live parameter tensors (PyTorch's
  \texttt{optimizer.step()} mutates parameters in place, never replacing
  the tensor objects, so this requires no further plumbing once the
  reference is actually shared). Existing callers
  (\texttt{rollout\_waterbox\_ase.py}, \texttt{waterbox\_langevin.py}) are
  unaffected --- \texttt{checkpoint\_path} remains the first positional
  argument, behaving exactly as before.
  \item \textbf{Unsafe RDF \texttt{rmax}.} The first draft computed
  \texttt{rmax} only from the 200 sampled reference DFT configurations, not
  from the replicas' own starting boxes --- reintroducing the exact bug
  \texttt{rollout\_waterbox\_ase.py} already learned the hard way (its own
  comment: box size varies configuration to configuration in this dataset,
  and ``the reference sample is what actually triggered
  \texttt{CellTooSmall} on the training box''). Fixed by computing
  \texttt{rmax} from the reference frames \emph{and} every replica's
  starting atoms object together, matching that established precedent
  rather than re-deriving a narrower one.
  \item A stale code comment, written before the shared-model bug above was
  found and fixed, that had come to assert the opposite of what the
  corrected code actually does --- corrected for anyone reading the code
  later.
\end{itemize}

\textbf{Not yet run.} Before a real fine-tuning run, the build order's own
prescription (Section~\ref{sec:q4-stable-plan}): a short smoke test
confirming the whole loop runs end to end without crashing or diverging.
\begin{quote}
\texttt{python src/train\_waterbox\_stable.py --ckpt
checkpoints/waterbox\_study\_zbl\_bonded\_ext70/water\_absolute/seed1/
best\_model.ckpt --out checkpoints/waterbox\_study\_zbl\_bonded\_ext70\_
stable/water\_absolute/seed1\_smoke --smoke-test}
\end{quote}
2 replicas, 3 outer iterations, short windows --- confirms checkpoints
save and \texttt{L\_obs}/\texttt{L\_QM} don't crash or produce
NaN/Inf, nothing more. \texttt{lr=1e{-}6} and \texttt{lambda\_qm=1.0} are
explicitly starting points, not tuned values (Section~\ref{sec:q4-stable-plan}
already flagged both as needing empirical verification) --- 3 iterations
cannot show whether they are well-chosen, only whether the loop is wired
correctly.

\textbf{Run: no crash, checkpoint saved, but every one of the 3 iterations
skipped its gradient step (only 1 pooled sample, below the estimator's
$N{\ge}2$ floor).} Not random noise --- a clean, diagnosable pattern.
Replica 0: 0 rewinds throughout. Replica 1: rewinds climbing 2, 4, 6 ---
every single stability check tripped, all 6 checks so far. Tracing the
mechanics: only the very first check of the very first iteration ran under
the (then-buggy) zero-velocity start below; every check after that,
including replica 1's own second check onward, already runs from a
properly redrawn 300\,K Maxwell-Boltzmann velocity (the rewind branch
inside \texttt{run\_stable\_sampling\_phase} always redraws regardless of
how a replica started). Replica 1 failing even with correctly thermalized
velocities argues against the cold start being the (sole) cause --- more
consistent with this replica's randomly-drawn training-split starting
\emph{geometry} itself being one of the ``bad'' configurations this
project has already established exist (the config2-vs-config4 epot-gap
spread, Section~\ref{sec:q1-rollout-update}, ranged 2.7--13\,eV). Read this
way, the rewind mechanism catching it immediately and repeatedly is it
doing exactly its job, not a bug.

The actual defect is that the smoke test had no headroom to survive that.
With only 2 replicas and \texttt{subsample\_stride=2} against a 2-check
window, even a perfectly well-behaved replica caps out at 1 collected
sample (Python's \texttt{frames[::2]} on a 2-element list keeps only index
0) --- so the moment either replica has a bad starting draw, the whole
phase is mathematically stuck below the sample floor, independent of
anything about $\lambda_{\text{QM}}$/$lr$/the estimator itself. The real
(non-smoke) defaults do not share this problem: 4 replicas $\times$ 5
checks/phase means a single well-behaved replica alone already clears
$N{\ge}2$ (3 samples), so one chronically-unstable replica among 4 would
not starve a real run the way it starved this two-replica smoke test.

Two fixes applied. \textbf{(1)} Replicas were being initialized at exactly
zero velocity in \texttt{fine\_tune\_stable}, inconsistent with the
already-\emph{verified} Maxwell-Boltzmann convention used everywhere else
in this project (Section~\ref{sec:q4-stable-step2}'s thermostat smoke
test, \texttt{rollout\_waterbox\_ase.py}). Per the mechanics above this
was not the cause of this specific outcome, but it was still an
unjustified departure from the one regime already confirmed to equilibrate
sanely, with no reason to introduce a second, unverified one for the real
fine-tuning loop --- fixed to draw a proper \texttt{MaxwellBoltzmannDistribution}
per replica (explicit \texttt{RandomState}, this project's own established
discipline). \textbf{(2)} The smoke test's own sizing was widened (2
$\to$ 3 replicas, 20 $\to$ 40 learn-window steps $\to$ 4 checks/phase) so
that, matching the real-default headroom argument above, a single
well-behaved replica alone clears $N{\ge}2$ even if every other replica
happens to be persistently unstable --- a real, expected outcome for some
starting configurations, not something a smoke test should be sized to be
fragile against.

\textbf{Re-run with both fixes: real gradient steps at every iteration.}
All 3 smoke-test iterations executed (\texttt{n\_samples=4} each,
comfortably above the $N{\ge}2$ floor), checkpoints saved at every step.
$\mathcal{L}_{\text{obs}}$(real): 0.403, 0.463, 0.481;
$\mathcal{L}_{\text{QM}}$: 17.6, 40.3, 12.8 --- both bouncing rather than
monotonically improving, expected and not concerning at $n{=}3$ iterations
with a batch size of 2 (explicitly never claimed to demonstrate anything
about convergence, only that the loop runs). The rewind pattern repeated
in a mildly interesting way: a \emph{different} replica (index 0 this
time, versus index 1 in the earlier under-provisioned run) tripped every
single check (4, 8, 12 rewinds), while the other two stayed clean --- a
different random 3-replica draw naturally pulls different starting
configs from the training split, so this isn't the same replica behaving
inconsistently, but it is now two independent draws out of five total
replicas that included a chronically-unstable one, worth keeping in mind
when scaling up.

\textbf{One thing surfaced in this run's output that needed resolving
before trusting anything built on it, not waved past on circumstantial
grounds alone}: a \texttt{UserWarning} that
\texttt{torchmdnet\_extensions::get\_neighbor\_pairs\_bkwd} has no
registered autograd kernel for backprop through it, ``may lead to silently
incorrect behavior.'' This targets \texttt{\_qm\_regularizer\_loss}'s force
term specifically: \texttt{force\_pred} comes from the model's own
internal \texttt{autograd.grad(energy, pos, create\_graph=True)} call, so
differentiating force-MSE with respect to the model's own parameters (what
\texttt{total\_loss.backward()} does) is a double-backward through that
same internal call --- exactly what the warning is about. Circumstantial
reasoning for why this is likely fine: ordinary supervised training
(\texttt{train\_waterbox.py}) needs this identical double-backward for any
run with nonzero force weight, i.e.\ every water-box run in this project;
if it were silently wrong, force accuracy across every checkpoint this
project already trusts would likely look far worse than it does. Not
proof, though --- an explicit ``may be silently wrong'' warning doesn't
get waved past on an analogy alone, this project's own established
discipline (\texttt{verify\_zbl\_units.py} and every other
\texttt{verify\_*.py} script) is to check directly.

\texttt{src/verify\_qm\_gradient.py} settles it directly: compares
\texttt{torch.autograd}'s own gradient of \texttt{\_qm\_regularizer\_loss}
against a manual central-finite-difference estimate (perturb a real model
parameter by $\pm\epsilon$ and re-evaluate the loss with no autograd
involved on that side at all) for the largest-magnitude-gradient parameter
in each of several tensors. If they agree, the double-backward path is
confirmed numerically correct for this exact model/PyTorch/torchmdnet
combination --- not just plausible by analogy.

\textbf{Run: 4/5 parameters agree within 5\%, one at $\approx$8.4\%
(\texttt{layers.5.linears\_tensor.3.weight}) --- a single borderline
mismatch, not the wildly-wrong-or-wrong-signed pattern a genuinely broken
double-backward would be expected to produce.} Rather than treat this as
confirmed evidence either way, the script was extended (not just re-argued
about) with two further checks before drawing a conclusion: (1) whether the
\emph{unperturbed} loss itself is bit-identical across repeated
evaluations with unchanged parameters and an identical batch --- GNN
message-passing architectures commonly rely on CUDA scatter/reduce
operations that are non-deterministic by default (this project's own code
already uses \texttt{scatter} elsewhere, e.g.\ \texttt{molecular\_zbl.py}),
which would inject noise into any finite-difference comparison independent
of whether the gradient itself is correct; (2) an automatic eps-sweep on
any mismatching parameter, since a genuine gradient bug should stay
robustly wrong across a wide range of step sizes while noise-limited
estimation typically does not.

\textbf{Resolved: confirmed noise, not a bug.} The determinism check found
real, nonzero non-determinism in this pipeline (5 repeated evaluations of
the identical unperturbed loss spread by $2.2{\times}10^{-3}$) --- direct
confirmation, not just plausibility, that CUDA scatter/reduce
non-determinism (or similar) is present here independent of any gradient
question. The eps-sweep on the one mismatching parameter settled it
cleanly: $\text{rel\_err}$ fell monotonically as eps grew ($26.4\%
\to 5.9\% \to 1.0\% \to 0.8\%$ at $\text{eps}=10^{-5},10^{-4},10^{-3},10^{-2}$),
tracking the noise-floor-implied uncertainty at each step size almost
exactly ($\sim\!110, 11, 1.1, 0.11$ respectively) --- at $\text{eps}=10^{-5}$
the noise floor exceeds the gradient's own magnitude ($\approx\!247$),
explaining the large apparent error there directly, while at
$\text{eps}=10^{-2}$ (noise floor $\approx\!0.11$, negligible) the
numerical estimate ($-251.0$) matches the analytic one ($-246.8$) to
within 0.8\%. A genuine double-backward bug would stay just as wrong
regardless of step size; noise-limited estimation does not, and this is
exactly that pattern. This also explains why only one of five parameters
flagged at the default eps: the noise floor at $\text{eps}=10^{-4}$
($\approx\!11$) sits at almost the same ratio to this parameter's gradient
magnitude as to \texttt{linear.weight}'s (247 vs.\ 248, the parameter that
passed one line above it) --- a coincidence of which side of the 5\%
threshold a shared noise floor happened to land on for two
near-identically-margined parameters, not a real outlier. \textbf{The
\texttt{get\_neighbor\_pairs\_bkwd} warning is confirmed benign for this
model/PyTorch/torchmdnet combination} --- a deprecation notice about a
legacy compatibility fallback that still computes the correct gradient,
not evidence of a silent error. $\mathcal{L}_{\text{QM}}$'s force term is
safe to use as implemented; no further action needed here.

\subsubsection{First real (non-smoke) pilot run: CUDA OOM, root-caused and fixed}
\label{sec:q4-stable-oom}

The first real pilot (\texttt{water\_absolute} seed 1, \texttt{n\_replicas=4},
30 outer iterations, otherwise real-run defaults) crashed on iteration 0's
$\mathcal{L}_{\text{QM}}$ computation: \texttt{torch.OutOfMemoryError},
46.54\,GiB of 47.5\,GiB \emph{actually allocated} (not fragmentation --- the
reported reserved-but-unallocated figure was only 444\,MiB), inside the
model's own internal \texttt{autograd.grad(energy, pos,
create\_graph=True)} call.

\textbf{Root cause}: \texttt{\_qm\_regularizer\_loss} accumulated a
differentiable running sum across \texttt{qm\_batch\_size=8} samples
(\texttt{total = total + per\_sample\_loss}) before a single
\texttt{.backward()} call. Each sample's forward pass computes forces
internally via \texttt{autograd.grad(\ldots, create\_graph=True)} ---
\texttt{create\_graph=True} is expensive specifically because it retains
the entire first-order backward graph so it can be differentiated a
second time. Summing 8 such terms before backing off held all 8 of these
graphs alive simultaneously. This project's own established
\texttt{batch\_size} for ordinary supervised training on this exact model
is 1--2, not 8 --- \texttt{qm\_batch\_size=8} was chosen without weighing
what \texttt{create\_graph=True} actually costs per sample, and this run
is the direct, no-longer-hypothetical confirmation of that cost.
\texttt{boltzmann\_estimator\_pseudo\_loss}'s own samples (\texttt{pooled\_U})
cannot avoid this same accumulation --- the covariance estimator
genuinely needs every sample's energy jointly available before its one
pseudo-loss backward call --- but the crash traceback shows the failure
happened specifically \emph{inside} $\mathcal{L}_{\text{QM}}$'s loop, after
\texttt{pooled\_U}'s own construction had already completed successfully:
direct evidence that \texttt{pooled\_U}'s footprint alone was tolerable,
and $\mathcal{L}_{\text{QM}}$'s additional accumulation on top of it is
what pushed memory over the edge --- not a case where every part of the
pipeline needed rework.

\textbf{Fix}: $\mathcal{L}_{\text{QM}}$ (unlike the covariance estimator)
is just a sum of independent per-sample terms, so it can legitimately call
\texttt{.backward()} immediately after each sample rather than
accumulating first. \texttt{\_qm\_regularizer\_loss\_and\_backward} does
exactly this --- gradients from repeated \texttt{.backward()} calls
accumulate into \texttt{.grad} by PyTorch's normal default behavior
(reset only by the caller's own \texttt{optimizer.zero\_grad()}), so this
is mathematically identical to one \texttt{.backward()} on the summed
loss, just without ever holding more than one sample's graph in memory at
once. The outer loop was restructured to match: \texttt{optimizer.
zero\_grad()} once, then \emph{two separate} backward calls
($\mathcal{L}_{\text{obs}}$'s own, then $\mathcal{L}_{\text{QM}}$'s
per-sample ones) rather than summing both into one loss tensor first ---
which has the further benefit that $\mathcal{L}_{\text{obs}}$'s own
(potentially large) graph is freed by its backward call before
$\mathcal{L}_{\text{QM}}$'s loop ever starts, rather than the two
competing for GPU memory simultaneously. The original tensor-returning
version was kept, not deleted, under its original name
(\texttt{\_qm\_regularizer\_loss}) specifically for
\texttt{verify\_qm\_gradient.py}, which needs to call \texttt{.backward()}
itself and re-evaluate the loss at perturbed parameter values with no side
effects --- safe there because that script only ever uses a batch size of
2, where accumulating a couple of \texttt{create\_graph=True} graphs was
never the problem. Both versions share the identical per-sample
computation via a new \texttt{\_qm\_per\_sample\_loss} helper, so this is
a change in \emph{how} gradients get accumulated, not a change to what is
being computed --- the gradient-correctness verification above
(Section~\ref{sec:q4-stable-step3}) still fully applies to both.

\textbf{Re-run: fix confirmed, full 30 iterations completed, no crash.}
\texttt{pooled\_U}'s own footprint was indeed tolerable on its own, exactly
as the earlier traceback evidence suggested --- no further memory lever
needed.

\subsubsection{Pilot run result: inconclusive on the loss trend, but a striking cross-validated finding on replica behavior}
\label{sec:q4-stable-pilot1}

$\mathcal{L}_{\text{obs}}$(real) is not yet interpretable as
succeeding or failing: after an initial low reading (0.36--0.48,
iterations 0--1), it settles into a noisy $0.57$--$0.80$ band for the rest
of the run with no clear trend either way --- unsurprising given each
iteration's value is a Monte Carlo estimate over only 2--5 pooled samples,
and 30 iterations is short. $\mathcal{L}_{\text{QM}}$ bounces between
$0.97$ and $68.4$ with no systematic drift --- the expected behavior for a
regularizer meant to hold existing DFT accuracy steady, not actively
improve it further.

\textbf{The striking result is in the rewind counts.} Tracking each of the
4 replicas' cumulative rewinds across all 30 iterations (150 stability
checks each):

\begin{center}
\begin{tabular}{lll}
\toprule
replica & final rewinds & trip rate \\
\midrule
A & 146/150 & 97\% \\
C & 141/150 & 94\% \\
D & 105/150 & 70\% \\
B & 8/150 & 5\% \\
\bottomrule
\end{tabular}
\end{center}

Three of four replicas are near-\emph{permanently} unstable across the
entire run, not merely occasionally --- only replica B is reliably
well-behaved. This directly explains why several iterations produced too
few samples to take a gradient step (iterations 4, 28, 29): the pipeline
is effectively running on one good replica's worth of signal, with the
other three mostly consuming rewind cycles rather than contributing
samples.

This is not a new, isolated problem --- it closely matches the seed-1
ext70 config-axis epot-gap finding from Section~\ref{sec:q1-rollout-update}
on this \emph{same} checkpoint: only cfg2 (gap 2.9\,eV, plateau 439\,K)
stood out as well-behaved among five starting configurations, with
cfg1/cfg3/cfg4/cfg5 all showing substantially larger gaps (7--13\,eV) and
correspondingly hotter rollouts --- also roughly a 1-in-5 good-config
ratio. Two independent diagnostics (post-hoc rollout-derived epot gap
there; live short-range-collapse checks during Langevin dynamics here)
converging on the same ``most starting configurations are bad for this
checkpoint, a clear minority are fine'' picture is meaningfully stronger
evidence than either alone --- not a coincidence of one particular
analysis.

\textbf{Practical implication}: with \texttt{n\_replicas=4}, effective
sample diversity per gradient step is lower than the replica count
suggests --- the pipeline mostly samples replica B's own neighborhood
of configuration space, with thin, intermittent contributions from D, and
essentially none from A/C. Two non-exclusive ways to address this before
scaling up further: increase \texttt{n\_replicas} (more absolute
well-behaved replicas even at roughly the same $\sim$25\% good ratio,
directly fixing the too-few-samples-per-iteration problem at iterations
4/28/29); or simply run longer, since the chronically-unstable replicas
are not wasted --- StABlE's whole design point is sampling near the edge
of stability, exactly what A/C/D are doing, they each just contribute very
little per iteration individually.

\textbf{Decision: both.} The two fixes are not competing --- more replicas
raises the expected number of well-behaved replicas per iteration (at the
same $\sim$25\% good ratio, $4\to8$ replicas takes the expectation from
$\sim$1 to $\sim$2, directly reducing the risk of an iteration like 4/28/29
starving below the $N\ge2$ sample floor), while more iterations is
needed independently to get $\mathcal{L}_{\text{obs}}$'s trend out of
pilot-scale noise regardless of replica count. Scaling only one axis would
leave the other pilot-identified problem unaddressed. Next run (not yet
executed): \texttt{water\_absolute} seed 1, \texttt{n\_replicas=8},
\texttt{n\_outer\_iterations=100} --- $8\times100=800$ replica-iterations
of sampling budget vs.\ the pilot's $4\times30=120$ ($\sim$6.7$\times$),
still short of \texttt{train\_waterbox\_stable.py}'s coded default of 200
iterations so this remains a checkpoint to inspect (rewind-rate pattern,
$\mathcal{L}_{\text{obs}}$ trend) before any longer commitment, not yet the
final production run. All other hyperparameters (learning rate, $\lambda_{QM}$,
window/stride, temperature) unchanged from the pilot. On success, repeat
identically for \texttt{water\_absolute+momentum} seed 1, then evaluate both
fine-tuned checkpoints via \texttt{evaluate\_waterbox.py --per-sample} and
\texttt{run\_rollout\_study.py} against the existing ext70 baseline.

\textbf{First attempt at $n\_replicas=8$ hit a new, real CUDA OOM ---
uncapped pool size scales with $n\_replicas$, not just $n\_outer\_iterations$.}
The crash happened inside the \texttt{pooled\_U} accumulation loop itself
(not $\mathcal{L}_{\text{QM}}$'s already-fixed accumulation), on the very
first outer iteration: 46.53 of 47.5\,GiB actually allocated, not
fragmentation. Root cause: \texttt{run\_stable\_sampling\_phase} records one
stable snapshot per \texttt{check\_interval} window
($\texttt{learn\_window\_steps}/\texttt{check\_interval}=5$ here), and
\texttt{take\_learn\_window(stride)} keeps every \texttt{stride}-th one --- a
best case (every window stable) of $\lceil 5/2\rceil=3$ samples per
replica. At iteration 0 every replica is still freshly-equilibrated from
its Maxwell-Boltzmann draw, so this best case is exactly what happens
before any replica has had a chance to drift into instability, giving a
\emph{theoretical} worst-case pool at $n\_replicas=4$ of 12 simultaneous
\texttt{\_config\_energy} calls (each holding a \texttt{create\_graph=True}
graph, since --- unlike $\mathcal{L}_{\text{QM}}$ --- the covariance-based
observable loss needs the whole pooled batch jointly and can't call
\texttt{.backward()} per sample). \textbf{First fix attempt (wrong):} a
\texttt{max\_pooled\_samples} parameter capping the pool at 12, reasoned
from that worst-case arithmetic and described at the time as matching
``the already-validated $n\_replicas=4$ worst case exactly.'' That
description was false, not merely imprecise: the 4-replica pilot's own
logged history never actually pooled more than 5 samples in any real
iteration (organic instability kept the true count well under the
12-sample theoretical ceiling), so 12 was never actually exercised, let
alone validated. Re-running with this ``fix'' at $n\_replicas=8$ OOM'd
again, at essentially identical memory usage (46.54 of 47.5\,GiB) to the
original uncapped crash --- direct evidence the cap size itself, not just
its absence, was the problem.

\textbf{Corrected fix, calibrated from measurement rather than arithmetic
on paper (the exact discipline CLAUDE.md's own GPU-OOM lesson already
prescribes, not followed carefully enough the first time):
\texttt{max\_pooled\_samples} lowered to 5, the pilot's actually-observed
maximum with zero assumed headroom.} The pooling loop now also logs
\texttt{torch.cuda.memory\_allocated()} after each sample on outer
iteration 0 specifically, so any future increase to this cap is set from a
real measured number rather than a second guess. This still decouples what
$n\_replicas$ was actually meant to buy (more diverse starting points to
sample from) from an unbounded per-iteration memory cost --- excess
collected snapshots beyond the cap are randomly subsampled away (via the
run's own rng) before any energy call is made, cheap since Snapshots are
plain numpy until then --- just at a genuinely conservative ceiling this
time. Not yet re-run.

\subsubsection{$n\_replicas=8$, $n\_outer\_iterations=100$ run: completed cleanly, but the substantive result is mixed-to-concerning, not a clean success}
\label{sec:q4-stable-8x100}

Re-run with \texttt{max\_pooled\_samples=5} completed all 100 iterations
with no crash --- confirms the corrected cap works. The pilot's original
``too few samples'' problem is genuinely fixed: the minimum
\texttt{n\_collected} across the whole run was 3, never below the
$N\ge2$ floor, zero iterations skipped. The cap is, however, binding
almost every iteration (99/100; mean 9.5 collected vs.\ 5 forwarded) ---
never down to the point of starving a gradient step, but discarding real
collected diversity most of the time. No console log with the
per-sample \texttt{[mem]} lines was retained from this run, so there is
still no direct measurement of how much memory headroom exists above 5;
that remains for a future run to log and use before raising the cap
again.

\textbf{$\mathcal{L}_{\text{obs}}$(real) plateaus, does not converge.} It
climbs from a low initial value to a noisy 0.6--0.9 band by iteration
$\sim$15, then stays there for the remaining 85 iterations --- a linear
fit over iterations 15--99 gives a slope of essentially zero
($-0.00017$/iteration, $-0.014$ total over 85 iterations). The pilot's 30
iterations could not distinguish ``not enough data yet'' from ``genuinely
flat''; 100 iterations now can, and the answer is flat. Overall stability
across all 8 replicas got slightly worse over the run, not better: total
trip rate rose from 49.5\% (iterations 0--49) to 59.5\% (iterations
50--99).

\textbf{A new finding: replica 3 underwent an abrupt, apparently permanent
regime change partway through the run.} For iterations 0--43 it behaved
almost as well as the two best replicas (mean 0.25 rewinds/iteration). At
iteration 44 its rate jumped $\sim$15$\times$ to 3.77 rewinds/iteration and
stayed there through iteration 99 --- every other replica's rate stayed
flat across the same boundary (the chronically-unstable ones were already
saturated near ceiling; the good ones stayed good). The coincidence with a
cluster of unusually large $\mathcal{L}_{\text{QM}}$ values right at
iterations 40/43/44 (34--50 vs.\ a typical 2--20) initially looked like a
plausible causal mechanism --- one anomalous training sample in that
iteration's \texttt{qm\_batch\_size=8} draw dominating the gradient step
enough to destabilize replica 3's neighborhood. \textbf{Checked
systematically rather than trusted as a coincidence, and it does not hold
up}: Spearman correlation between $\mathcal{L}_{\text{QM}}$ at iteration
$i$ (the gradient step actually applied at $i$, hence affecting sampling
at $i+1$, not $i$ itself --- the correct causal ordering given the loop
applies the gradient step \emph{after} that iteration's own sampling phase)
and the total rewind jump at iteration $i+1$ is essentially null
($\rho=-0.14$, $n=99$); the 8 largest $\mathcal{L}_{\text{QM}}$ values in
the entire run are followed by entirely unremarkable next-iteration
instability (15--24 total rewinds, right at the overall mean of 21.9). The
more likely explanation is ordinary long-run stochastic drift: a
replica's position is never reset between outer iterations (only within a
phase, on a rewind), so over up to 10{,}000 net MD steps across the full
run it can genuinely random-walk from a stable into an unstable
neighborhood --- the same ``$\sim$1-in-4 configurations are good, most
aren't'' pattern this project already established across different
\emph{starting} configurations (Section~\ref{sec:q4-stable-pilot1}'s
epot-gap cross-validation), now showing up within a single replica's own
long trajectory rather than only across distinct starting points.

\textbf{Not yet decided: whether to repeat this for
\texttt{water\_absolute+momentum} seed 1 as planned, or first investigate
whether \texttt{lr}/\texttt{lambda\_qm} need adjusting given the flat
$\mathcal{L}_{\text{obs}}$ trend, and whether to pull back the saved
\texttt{stable\_iter*.ckpt} checkpoints for a direct
\texttt{evaluate\_waterbox.py --per-sample} + rollout comparison against
the ext70 baseline} --- a flat pseudo-loss does not by itself establish
that the checkpoint's real rollout behavior did or did not improve, since
RDF-matching is a proxy for the actual target (stable NVE dynamics), not
the target itself.

\subsubsection{Both checks now done: static accuracy unchanged, rollout stability shows no detectable effect --- convergent evidence the run's null $\mathcal{L}_{\text{obs}}$ trend reflects a real lack of progress, not just a noisy proxy}
\label{sec:q4-stable-8x100-eval}

\textbf{Static accuracy (\texttt{evaluate\_waterbox.py --per-sample} on
\texttt{stable\_final.ckpt}, compared directly against this exact
checkpoint's own pre-fine-tuning row in
\texttt{results/waterbox\_study\_zbl\_bonded\_ext70/raw\_results.csv}):}
energy MAE 2.451$\to$2.227\,eV, force MAE 0.798$\to$0.739\,eV/\AA, mean
per-molecule momentum violation 62.65$\to$58.25, max violation
2086.0$\to$1840.5 --- every metric moved slightly in the favorable
direction, a clean single-checkpoint comparison not confounded by this
pipeline's known checkpoint-selection noise (no ``best epoch'' cherry-picking
involved, just the literal final checkpoint). The per-sample worst rows
(up to 22.5\,eV) are nowhere near this project's own previously-established
batch-of-1-anomaly scale ($\sim$16{,}000\,eV) and their
\texttt{min\_intermol\_dist\_A} values sit within the range of already-known
real anomalous DFT configurations, not a computational artifact --- the
reported mean is trustworthy as-is.

\textbf{Rollout comparison (\texttt{src/compare\_stable\_rollout.py}, velocity
axis, $n=5$, matching the existing ext70 seed-1 baseline's exact
\texttt{DATA\_SEED}/\texttt{test\_config\_index}/\texttt{dt} for direct
comparability): no detectable effect.} Plateau temperature 772.1$\pm$20\,K
(before) vs.\ 768.2$\pm$19\,K (after) --- and, critically, the per-replicate
paired differences do not even agree in sign: 3 of 5 velocity seeds run
\emph{hotter} after fine-tuning, 2 run \emph{cooler}, spanning $-30.1$ to
$+15.2$\,K (a 45.3\,K spread) against a mean shift of only $-3.9$\,K. This is
the opposite of the unanimous 10/10 directional pattern this project has
previously treated as evidence of a real, systematic effect
(Section~\ref{sec:q3-progress}'s momentum-vs-absolute finding) --- by that
same standard, this is indistinguishable from run-to-run noise, not a small
win. (Sanity check: the \texttt{before\_finetune} arm reproduces the
original baseline almost exactly, 772.1$\pm$20\,K here vs.\ 772.2$\pm$15\,K
on record --- confirming the comparison script is wired correctly and this
is a fair test, not a broken harness silently returning null.)

\textbf{Conclusion: this specific $n\_replicas=8$, $n\_outer\_iterations=100$,
$lr=10^{-6}$, $\lambda_{QM}=1$ configuration produced no measurable change on
any of three independent measures} --- the training loss itself, static
accuracy, and the actual downstream quantity StABlE fine-tuning exists to
improve. Static accuracy's own mild improvement is reassuring (fine-tuning
did not break the model) but is not itself the goal. Two of the three
measures (loss trend, rollout outcome) point the same direction and were
obtained independently, which is stronger evidence than either alone that
the flat $\mathcal{L}_{\text{obs}}$ trend reflects a real lack of learning
progress rather than $\mathcal{L}_{\text{obs}}$ simply being too noisy a
proxy to see real improvement through. The most likely lever, not yet
tested: \texttt{lr=1e-6} was flagged from the start as an untuned starting
point (``two orders of magnitude below the original training run's
$10^{-4}$'') --- with a dead-flat loss curve and zero rollout change now
both on record, that is the natural next thing to try adjusting before
repeating this recipe unchanged on \texttt{water\_absolute+momentum} seed 1.

\subsubsection{Which hyperparameters to change, checked against primary sources rather than guessed}
\label{sec:q4-stable-hparam-lit}

Before changing anything further, checked what StABlE's own authors actually
used for the same kind of system, and what fine-tuning-specific learning-rate
practice looks like elsewhere in the NNIP literature --- per this project's
own established discipline of verifying against primary sources rather than
folklore.

\textbf{What \citet{raja2025stable} actually used, read directly from the
paper's own Appendix A.6 (Supplementary Table 2), not assumed:}

\begin{center}
\begin{tabular}{lllllll}
\toprule
system & estimator & $\alpha$ (lr) & $\lambda$ & $t$ (steps/cycle) & $R$ (replicas) & $B$ (minibatch) \\
\midrule
Aspirin (SchNet) & Global & 0.001 & 10 & 2000 & 128 & 40 \\
Ac-Ala3-NHMe (NequIP) & Global & 0.001 & 10 & 2000 & 128 & 40 \\
Water (GemNet-T) & Local & 0.003 & 0 & 1000 & 8 & 4 \\
\bottomrule
\end{tabular}
\end{center}

The water row is the directly comparable case --- same system, and their
``local'' per-molecule Boltzmann estimator is the direct analogue of our own
per-fragment approach --- and every one of its four numbers differs sharply
from ours:
\begin{itemize}
  \item $\alpha=0.003$ vs.\ our $10^{-6}$ --- roughly 3000$\times$ higher.
  \item $\lambda=0$ vs.\ our $\lambda_{QM}=1.0$ --- they use \emph{zero} QM
  regularization during StABlE training for water specifically, ``so that
  any improvements in stability or accuracy... can be attributed to the
  learning signal from the reference observable, as opposed to the
  regularization'' (Appendix A.6).
  \item $t=1000$ steps at their 1\,fs timestep (their own Table 4) = 1\,ps of
  simulated time per learning window, vs.\ our
  $\texttt{learn\_window\_steps}=100$ at $dt=0.5$\,fs = 0.05\,ps --- a
  20$\times$ shortfall. Their own stated general guideline: ``a frequency of
  1 picosecond should be sufficient for structural observables of small
  molecular systems'' (Appendix A.6) --- not just their specific choice, a
  stated rule of thumb for this exact problem class.
  \item $R=8$ --- matches our own choice exactly, a reassuring confirmation
  that the replica-count decision was already well-aligned with this primary
  source, even though the other three were not.
\end{itemize}

Crucially, water's entire reported result (median stable time 312\,ps
$\to$ stable for 1\,ns across all 5 held-out initial conditions) came from
\textbf{a single StABlE cycle} --- ``We perform a single cycle of simulation
and learning'' (Section 4.4) --- not from repeating a much smaller step many
times, the strategy attempted here (100 outer iterations at a tiny learning
rate and a 20$\times$-undersized window).

\textbf{Optimizer mismatch complicates directly copying $\alpha=0.003$.} The
paper states an SGD optimizer explicitly for aspirin (``We perform four
complete StABlE cycles of simulation and learning with a SGD optimizer,''
Section 4.1) and does not state a different optimizer for water;
\texttt{train\_waterbox\_stable.py} uses \texttt{torch.optim.Adam} instead.
Adam's per-parameter update is normalized by an estimate of the gradient's
second moment, so it does not scale the same way as SGD's raw
gradient-times-lr update --- the two are not directly interchangeable at the
same nominal learning rate. Three independent, Adam-family reference points
avoid this ambiguity entirely:
\begin{itemize}
  \item \citet{kingmaba2015adam}'s own Adam paper states $\alpha=0.001$ as a
  ``good default setting for the tested machine learning problems''
  (Algorithm 1) --- already three orders of magnitude above our $10^{-6}$,
  for an arbitrary problem, not even a fine-tuning-specific recommendation.
  \item This exact model's own original supervised pre-training
  (\texttt{train\_waterbox.py}) uses \texttt{torchmdnet}'s own
  \texttt{LNNP.configure\_optimizers}, confirmed directly from its source
  \citep{torchmdnet-code} to construct an \texttt{AdamW} optimizer at
  $lr=10^{-4}$. \texttt{AdamW} shares Adam's normalization behavior (differs
  only in how weight decay is applied), so this is a same-optimizer-family,
  same-architecture, already-empirically-successful reference point, not a
  cross-method guess.
  \item MACE's official fine-tuning documentation \citep{mace-finetuning-docs}
  (an AdamW-based foundation-model fine-tuning workflow --- the closest
  real-world analogue to ``fine-tune an already-converged NNIP without
  forgetting it'') states verbatim: ``The default learning rate for
  multihead finetuning is 0.0001, which is lower than the default learning
  rate for single head finetuning (0.001). This helps to stabilize the
  training process and prevent overfitting.''
\end{itemize}

All four reference points --- StABlE's own raw water rate, Adam's generic
literature default, this exact model's own proven from-scratch rate, and
MACE's own most-conservative NNIP fine-tuning default --- land in the same
$10^{-4}$--$10^{-3}$ band, one to three orders of magnitude above our
$10^{-6}$. None of them individually justifies matching Raja et al.'s raw
0.003 exactly (the SGD/Adam mismatch), but all four agree the two-orders-of-
magnitude discount used here (\texttt{train\_waterbox\_stable.py}'s own
docstring: ``well below the original training run's $10^{-4}$'') was too
conservative given the now-established null result.

\textbf{Recommended changes for the next run, in order of confidence:}
\begin{enumerate}
  \item $\alpha$: $10^{-6}\to10^{-4}$ --- matches this exact model's own
  proven from-scratch AdamW rate and MACE's own most-conservative NNIP
  fine-tuning default exactly; a 100$\times$ increase, not an attempt to
  match Raja et al.'s un-comparable SGD rate.
  \item $\lambda_{QM}$: $1.0\to0$ --- directly matches Raja et al.'s own
  water-specific choice and stated rationale. Lower risk than it looks: the
  static-accuracy check already completed
  (Section~\ref{sec:q4-stable-8x100-eval}) found accuracy did not regress
  even at $\lambda_{QM}=1$ over 100 iterations, so the ``regularizer
  prevents forgetting'' concern that motivated a nonzero $\lambda_{QM}$ in
  the first place does not appear to be a live risk at the scale tried so
  far.
  \item $dt$: $0.5\to0.1$\,fs --- a newly-identified, previously-unaddressed
  likely confound, not from this literature check: this checkpoint carries
  the bonded-exclusion ZBL prior, and this project already established
  (Section~\ref{sec:q4-negative-result}) that ANY ZBL-variant checkpoint
  needs $dt=0.1$ to avoid a genuine numerical-integration artifact from
  ZBL's steep short-range repulsive force being under-resolved at 0.5\,fs
  --- ``a well-documented risk of adding hard repulsive priors to NVE MD,
  not a wiring/calibration bug.'' That finding was established for NVE
  rollouts specifically; \texttt{train\_waterbox\_stable.py}'s Langevin
  sampling phase was run at the module's own default $dt=0.5$ in both
  completed runs (the pilot and the $n\_replicas=8$ run) without this
  correction applied. A thermostat can dissipate some of the resulting
  energy error rather than let it compound catastrophically as in NVE, so
  this may not be equally severe here --- but the rewind mechanism's own
  stability criterion is exactly the kind of short-range-collapse event this
  timestep mismatch would be expected to trigger spuriously, so an unknown
  fraction of the very high rewind rates observed so far may be a numerical
  artifact rather than a genuine reflection of model quality. Worth
  controlling for before trusting rewind-rate patterns (e.g., replica 3's
  regime change, Section~\ref{sec:q4-stable-8x100}) as purely physical
  findings.
  \item $\texttt{learn\_window\_steps}$: $100\to10{,}000$ --- not copying
  Raja et al.'s raw $t=1000$ directly (their number was computed at their
  own 1\,fs timestep), but matching the underlying quantity that actually
  matters, physical simulated time (their own explicit $\sim$1\,ps choice
  and general guideline): $10{,}000\times0.1\text{ fs}=1\text{ ps}$, exactly
  matching their water-specific choice once expressed in the same units.
  \item $\texttt{check\_interval}$: $20\to100$ --- not from the literature
  (their paper does not describe an equivalent mid-window rewind-check
  parameter), but the author's own derived consistency choice: preserves the
  same $\sim$10\,fs physical-time interval between stability checks already
  validated in this project's own thermostat/rewind smoke tests, scaled for
  the new, smaller $dt$.
  \item $\texttt{n\_outer\_iterations}$: $100\to\sim5$ --- Raja et al.'s
  entire water result came from a single such cycle; with the window itself
  now 100$\times$ longer in raw step count (though only
  $\sim5\times$ more total simulated time across a full run than before,
  once $\texttt{n\_outer\_iterations}$ is reduced accordingly --- $10{,}000
  \times5=50{,}000$ vs.\ the previous $100\times100=10{,}000$ total steps per
  replica), a handful of properly-sized iterations is a closer match to
  their design than continuing the ``many tiny steps'' strategy that just
  produced a null result.
  \item $\texttt{n\_replicas}$: unchanged at 8 --- already matches Raja et
  al.'s own choice.
  \item $\texttt{max\_pooled\_samples}$: unchanged at 5 for now, still an
  unmeasured GPU-memory ceiling rather than a deliberate choice --- but Raja
  et al.'s own explicit guidance (``The minibatch size $B$ should also be
  chosen as large as possible to minimize variance in the Boltzmann
  Estimator, while remaining within GPU memory limits,'' Appendix A.6) is
  now direct literature support for actually obtaining the per-sample
  memory measurement (the \texttt{[mem]} logging already added, not yet used
  to inform a decision) before the next run, rather than continuing to
  guess.
\end{enumerate}

\subsubsection{That recipe has run --- catastrophic, unambiguous divergence, not a subtle result}
\label{sec:q4-stable-lr-divergence}

$\mathcal{L}_{\text{QM}}$ (logged every iteration regardless of
$\lambda_{QM}$ --- a pure monitoring signal of the model's real supervised
accuracy on ordinary training-split samples, decoupled from what actually
drives the gradient when $\lambda_{QM}=0$): iter0 $=20.13$ (baseline, before
any gradient step), iter1 $=1189.72$ (a $59\times$ jump from a
\emph{single} gradient step alone, since $\lambda_{QM}=0$ means that step
used only $\mathcal{L}_{\text{obs}}$), iter2 $=9883.22$, iter3
$=36324.11$, iter4 $=83850.35$ --- a $4166\times$ blowup over 5 iterations,
monotonic, no sign of stabilizing. $\mathcal{L}_{\text{obs}}$(real) itself
stays too noisy to read a trend from at only 5 points ($0.69$, $0.84$,
$0.61$, $0.73$, $0.56$ --- the same ``too few iterations'' problem as the
original 30-iteration pilot, not itself alarming on its own); overall
rewind trip rate stays in the same 47--53\% noisy band as the previous
$n\_replicas=8$ run, no clear direction either way. \textbf{The final
checkpoint from this run should be treated as unusable} --- the
$\mathcal{L}_{\text{QM}}$ trajectory implies severe, real degradation of
ordinary supervised accuracy; not worth spending compute on
\texttt{evaluate\_waterbox.py}/rollout checks to confirm what is already
unambiguous.

\textbf{Root cause, most likely: stacking the $\alpha$ and $\lambda_{QM}$
changes together removed the only thing anchoring supervised accuracy while
also taking a much bigger step, and Adam's per-parameter normalization ---
the SGD/Adam mismatch already flagged as an open risk in
Section~\ref{sec:q4-stable-hparam-lit}, before this run --- likely makes
that step even more aggressive in practice than Raja et al.'s own SGD-based
recipe intended.} Each individual change was well-justified against the
literature in isolation; the interaction between them was not checked
before running both at once in the same experiment.

\textbf{Diagnostic run (3 iterations, $\alpha=10^{-4}$, $\lambda_{QM}$
restored to 1.0): falsified the missing-regularizer hypothesis --- restoring
$\lambda_{QM}$ made the immediate blowup \emph{worse}, not better.}
$\mathcal{L}_{\text{QM}}$: iter0 $=20.13$ (baseline), iter1 $=163{,}192$
(an $8107\times$ jump from a single step, vs.\ $59\times$ at the same
$\alpha$ with $\lambda_{QM}=0$ --- restoring the regularizer made the very
first step's damage $\sim\!137\times$ worse, not better), iter2 $=9734$
(a large partial pullback, not a continued monotonic climb like the
$\lambda_{QM}=0$ run showed). This is the wrong direction for ``the missing
regularizer caused it,'' and the spike-then-pullback shape (rather than
smooth monotonic divergence) is itself informative --- it is the signature
of a step size too large relative to the noise in the gradient estimates,
not of a missing counterweight. At $\alpha=10^{-6}$, even an occasional bad
gradient (from either loss term) only ever produced a tiny, harmless nudge;
at $\alpha=10^{-4}$ ($100\times$ larger), any sufficiently noisy or unlucky
gradient draw --- an outlier sample in the 8-sample QM regularizer batch, or
the inherently noisy 5-sample $\mathcal{L}_{\text{obs}}$ estimate --- can
now produce a devastating parameter jump regardless of which loss term
supplies it. Turning $\lambda_{QM}$ back on did not add a counterweight; it
added a second, independently noisy gradient source riding on top of the
same oversized step.

\textbf{Conclusion: the learning rate itself is the problem, not which loss
term is active.} Two real, bracketing data points are now on record:
$\alpha=10^{-6}$ confirmed too small to move anything measurably
(Section~\ref{sec:q4-stable-8x100-eval}), $\alpha=10^{-4}$ confirmed
catastrophically too large under both $\lambda_{QM}=0$ and $\lambda_{QM}=1$.
\textbf{Next step (not yet run): bisect --- try $\alpha=10^{-5}$}, ten times
below the confirmed failure and ten times above the confirmed-ineffective
rate, with $\lambda_{QM}=1.0$ (the already-safe default) restored.

\subsubsection{Bisection run ($\alpha=10^{-5}$): better than $10^{-4}$, but still not safe}
\label{sec:q4-stable-lr-bisect1}

$\mathcal{L}_{\text{QM}}$ relative to its own iter0 baseline (20.10): iter1
$=81.8\times$, iter2 $=9.5\times$ (a partial recovery), iter3
$=13.5\times$ (climbing again), iter4 $=41.7\times$ (climbing further). The
peak (1645 at iter1) is $\sim\!100\times$ less severe than
$\alpha=10^{-4}$'s peak (163{,}192), but the value never returns near the
healthy baseline and is trending upward again by the run's end rather than
settling --- a recover-then-climb-again shape, the signature of continued
instability rather than convergence to a new stable point.
Rewind/rollout-stability trip rate stayed flat around 47--50\% overall,
consistent with every other run so far --- not yet informative at this
scale (the supervised-accuracy damage clearly present in
$\mathcal{L}_{\text{QM}}$ has not yet had time to manifest there).

\textbf{Three real, bracketing data points are now on record}: $\alpha=
10^{-6}$ confirmed too small (zero measurable effect over 100 iterations,
Section~\ref{sec:q4-stable-8x100-eval}); $\alpha=10^{-5}$ confirmed still
causing real, sustained supervised-accuracy damage (10--80$\times$
baseline, still climbing at the end of a short run); $\alpha=10^{-4}$
confirmed catastrophic ($8000\times$+ baseline,
Section~\ref{sec:q4-stable-lr-divergence}). The safe range is narrower than
$10^{-5}$. \textbf{Next step (not yet run): continue bisecting --- try
$\alpha=3\times10^{-6}$} (the geometric midpoint between the confirmed-safe
$10^{-6}$ and the confirmed-still-too-large $10^{-5}$), $\lambda_{QM}=1.0$
unchanged.

\subsubsection{Bisection run ($\alpha=3\times10^{-6}$): a qualitatively different, bounded result}
\label{sec:q4-stable-lr-bisect2}

$\mathcal{L}_{\text{QM}}$ relative to its own iter0 baseline (20.12, itself
remarkably consistent across every diagnostic run so far --- 20.10--20.13
on the same untouched starting checkpoint under different random batch
draws, a tight, low-noise floor): iter1 $=7.85\times$, iter2 $=2.71\times$,
iter3 $=0.59\times$ (below the baseline reading), iter4 $=2.73\times$. This
is a qualitatively different shape from both higher rates tested --- noisy,
but bounded within roughly $0.6$--$7.9\times$ baseline, with no accelerating
or sustained-climb trend across the 5 iterations.

One replica (index 6) flipped from consistently well-behaved in the
$\alpha=10^{-5}$ run (1--13\% per-iteration trip rate,
Section~\ref{sec:q4-stable-lr-bisect1}) to chronically unstable here
(41--88\% per-iteration, driving overall trip rate from
$\sim$47--53\% up to $\sim$51--61\%) --- same starting configuration both
times (same default seed). Read as the same kind of chaotic single-replica
drift already documented for replica 3 in
Section~\ref{sec:q4-stable-pilot1}'s original run, not evidence that a
lower learning rate is worse for rollout stability: every other replica's
qualitative behavior (which stay good, which are chronically bad) is
consistent across both runs --- one replica's contribution, not a
population-wide shift.

\textbf{Given $\mathcal{L}_{\text{QM}}$ no longer shows runaway or
sustained-climb behavior at this rate, the next step is not to keep
bisecting lower, but to extend this exact setting to more iterations} ---
both to confirm the boundedness holds over a longer horizon (5 iterations
is still not enough to rule out later acceleration) and to finally
accumulate enough data points to read whether $\mathcal{L}_{\text{obs}}$
shows real movement, which is the actual point of this entire
hyperparameter search. \textbf{Not yet run: $\alpha=3\times10^{-6}$,
$\lambda_{QM}=1.0$, extended to $\sim$20 outer iterations} (comparable
statistical power to the original pilot, Section~\ref{sec:q4-stable-pilot1}).

\subsubsection{Extended run (20 iterations): a genuinely stable, promising loss trajectory --- and a still-null rollout result}
\label{sec:q4-stable-lr-final}

$\mathcal{L}_{\text{QM}}$ stayed bounded for the entire 20 iterations
(0.16--7.85$\times$ baseline, its two lowest readings of the whole run at
iterations 18--19), confirming $\alpha=3\times10^{-6}$ is genuinely
non-destructive over a real horizon, not just a lucky first few steps.
$\mathcal{L}_{\text{obs}}$ shows a hint of real movement --- first-10 vs.\
last-10 iteration means, $0.851\to0.724$, Welch $t=-2.42$ (nominally
$p<0.05$) --- though a full linear fit across all 20 points is weaker
($t=-1.28$, $R^2=0.08$), so this is encouraging rather than established,
consistent with this project's own repeated lesson that small-$n$ trends
need replication before being trusted. Replica 6, which flipped bad in the
5-iteration diagnostic, stayed chronically unstable (63--100\% trip rate)
for the full 20 iterations and never recovered --- consistent with
Section~\ref{sec:q4-stable-pilot1}'s finding that a replica which crosses
into a bad neighborhood tends to stay there.

\textbf{Both direct target checks on this run's final checkpoint: a mixed
static-accuracy result, and still no detectable rollout effect.}
\texttt{evaluate\_waterbox.py --per-sample} vs.\ this checkpoint's own
pre-fine-tuning baseline: energy MAE $2.451\to3.821$\,eV ($+56\%$, a real
regression, not noise), force MAE $0.798\to0.753$\,eV/\AA\ ($-5.7\%$), mean
momentum violation $62.65\to60.18$ ($-4.0\%$), max violation
$2086.0\to1878.1$ ($-10.0\%$). The worst per-sample rows are the same
magnitude and the same known-anomalous configurations already
characterized in this project (not a fresh ZBL-singularity artifact) ---
the energy regression is real, and is itself useful evidence: unlike the
$\alpha=10^{-6}$ run (which barely moved the model and correspondingly
barely moved static accuracy), real parameter movement is now happening.
\texttt{compare\_stable\_rollout.py} (velocity axis, $n=5$, identical
settings to every prior comparison): plateau temperature $775.2\pm27$\,K
(before) vs.\ $768.5\pm11$\,K (after) --- a $-6.7$\,K shift, smaller than
the first attempt's already-null $-3.9$\,K, and again mixed in direction
per replicate (2 hotter, 2 cooler, 1 flat), dominated by a single outlier
replicate ($-46.1$\,K; every other replicate within $\pm14$\,K). No
detectable effect, by the same standard applied throughout this project.

\textbf{This combination --- bounded loss, a hint of $\mathcal{L}_{\text{obs}}$
movement, a real (if unwanted) shift in static accuracy, yet zero
measurable rollout effect --- points at something more structural than a
learning-rate problem.} Re-reading \citet{raja2025stable}'s Section 3.1
(Definition 2, their Equation 5) rather than just their hyperparameter
table: ``unphysical configurations can occur within localized regions of
the simulation domain, such as collisions between two molecules in a
large condensed-phase system. Due to spatial averaging, global observables
$g(\Gamma)$ may be insensitive to these localized events, limiting the
ability to identify unphysical states.'' This is exactly this project's
own situation --- \texttt{\_stacked\_rdf} computes one whole-box RDF per
configuration, averaged over all 64 water molecules, while
Section~\ref{sec:q1-rollout-update}'s own decisive check already
established that rollout heating is predicted by per-configuration,
localized force error, not a global/uniform property. A global RDF is, by
the original authors' own stated reasoning, structurally unable to see or
correct exactly the failure mode this project independently characterized.

\textbf{This should have been flagged during the initial literature check
(Section~\ref{sec:q4-stable-hparam-lit}) --- the water row's ``Local''
estimator type was noted in the table but its methodological significance
was not investigated at the time, only the numeric hyperparameters were.}
Raja et al.'s own fix, used specifically for water (not aspirin): extract a
local neighborhood of $n$ atoms (one water molecule) from each sampled
global state, using the fact that an MLFF already parameterizes its global
energy as a sum over atomic contributions --- ``the local energy is easily
obtained by noting that MLFFs parameterize their global energy prediction
$U_\theta$ as a sum over individual atomic energies'' --- so
$U_\theta(\gamma)=\sum_i U_\theta(\gamma^{(i)})$ for just the atoms in the
neighborhood, computed from the model's real, full-box forward pass (not a
re-run on an isolated fragment, which would lose the true local chemical
environment). The localized estimator (Equation 5) is otherwise
identical in form to the global one (Equation 4) --- same covariance
structure, just $U_\theta(\gamma)/g(\gamma)$ in place of
$U_\theta(\Gamma)/g(\Gamma)$. Their chosen water observable was the mean
O--H bond length, a scalar per-molecule quantity directly targeting the
known failure mode (localized bond stretching), not a structural histogram.
They also note extracting \emph{multiple} local neighborhoods per global
state (up to 64 per configuration here, one per molecule) to increase
sample size --- a likely structural explanation for why their water
recipe needed only $R=8$ replicas and a single cycle while this project's
global-observable implementation has been chronically sample-starved
(motivating \texttt{max\_pooled\_samples} in the first place): their
design gets $\sim$64 samples per collected snapshot ``for free''; this
project's gets exactly 1.

Feasibility check (not yet a commitment to implement): TorchMD-Net's own
architecture already computes per-atom energy contributions internally,
aggregated by a separate \texttt{output\_model} step into the final scalar
--- confirmed via its public documentation, not assumed --- so accessing a
per-atom decomposition is a matter of reaching into that intermediate
tensor, the same ``operate on the full-box output, then group/sum by
molecule membership'' pattern this project already uses and trusts for the
per-fragment momentum loss (\texttt{physics\_losses.py}), just applied to
energy instead of force. Not yet verified against an installed
\texttt{torchmdnet} copy (same caveat as every other water-box script in
this repo) --- the exact API needed (whether the per-atom tensor is
reachable without monkey-patching internals) is unconfirmed.

\textbf{Decision point, not yet resolved}: continue tuning within the
current global-RDF implementation (weak justification now, given the
authors' own stated reasoning for why a global observable should not be
expected to fix a localized failure mode), or implement the localized
per-molecule Boltzmann estimator matching Raja et al.'s actual validated
water recipe --- a substantially larger engineering effort, but the first
approach with a specific, primary-source-grounded reason to expect it
would address the mechanism this project has already established as the
real cause of rollout instability.

\subsubsection{Implementation scope for the localized per-molecule Boltzmann estimator}
\label{sec:q4-stable-local-estimator-scope}

Verified directly against two primary sources before committing to any
design choice below --- \citet{raja2025stable}'s own algorithmic
description (Section 3.1, Definitions 1--2, Equations 3--5; Section 4.4)
for what must change conceptually, and torchmd-net's actual source
\citep{torchmdnet-code} (\texttt{torchmdnet/models/model.py}'s
\texttt{TorchMD\_Net.forward}, \texttt{torchmdnet/priors/base.py}'s
\texttt{BasePrior}, \texttt{torchmdnet/priors/zbl.py}) for what is
mechanically possible, not assumed from general documentation.

\textbf{What changes conceptually (grounded in Raja et al., Section 3.1/4.4).}
Definition 2 (Equation 5) is otherwise identical in form to Definition 1
(Equation 4, already implemented and verified in
\texttt{boltzmann\_estimator.py}) --- same covariance-based construction,
$\frac{N}{k_BT(N-1)}\left[\hat{E}[g(\gamma)]\hat{E}[\nabla_\theta
U_\theta(\gamma)]^\top - \hat{E}[g(\gamma)\cdot\nabla_\theta
U_\theta(\gamma)^\top]\right]$, just with the global energy/observable
$U_\theta(\Gamma)/g(\Gamma)$ replaced by their local counterparts
$U_\theta(\gamma)/g(\gamma)$ for a neighborhood $\gamma$ of $n<N$ atoms.
\textbf{This means \texttt{boltzmann\_estimator\_pseudo\_loss} itself needs
NO changes} --- it is already generic over sample count and observable
dimensionality (verified for both scalar and multi-bin vector observables
in its own test suite), so it is reused as-is; only the layer that
constructs $g$ and $U$ upstream of it changes. Concretely, per water
molecule (one local neighborhood $\gamma=\{O,H,H\}$, $n=3$):
\begin{itemize}
  \item $U_\theta(\gamma) = \sum_i U_\theta(\gamma^{(i)})$, the sum of that
  molecule's own 3 atomic energy contributions, computed from the model's
  real, full-box forward pass (per the paper's own statement, quoted
  above) --- not a separate, truncated forward pass on 3 isolated atoms,
  which would lose the true local chemical environment the full box
  provides.
  \item $g(\gamma)$ = the mean O--H bond length of that molecule's own two
  O--H distances (``we use the mean hydrogen-oxygen bond length as our
  training observable,'' Section 4.4) --- a single scalar per molecule,
  not a histogram; no RDF binning/Gaussian smearing needed at all, a real
  simplification over the current \texttt{\_stacked\_rdf} machinery.
  \item Reference $g_{\text{target}}$: the same scalar, averaged across
  every molecule in a large sample of real reference DFT configurations
  (reusing \texttt{\_sample\_reference\_frames}) --- one number, not a
  reference RDF curve.
  \item The rewind/instability criterion in \texttt{waterbox\_langevin.py}
  does \emph{not} change --- Raja et al. themselves use ``a minimum
  intermolecular distance criterion (Section A.5) to detect instabilities
  during training'' for water specifically, which is exactly this
  project's own already-implemented, already-validated criterion
  (\texttt{diagnose\_short\_range\_collapse.frame\_violates\_floors}). Only
  what happens \emph{after} a snapshot is accepted as stable changes.
\end{itemize}

\textbf{How this is implemented mechanically (grounded directly in
torchmd-net's source, read line-by-line, not assumed).}
\texttt{TorchMD\_Net.forward} (lines 587--628 of
\texttt{torchmdnet/models/model.py}) computes per-atom features via
\texttt{self.representation\_model(...)}, converts them to per-atom scalar
energies via \texttt{x = self.output\_model.pre\_reduce(x, v, z, pos,
batch, box=box)}, then gives any active prior models a chance to modify
\emph{this same per-atom tensor} before the final reduction
(\texttt{for prior in self.prior\_model: x = prior.pre\_reduce(x, z, pos,
batch, extra\_args)}) --- before finally reducing to one scalar per
configuration via \texttt{self.output\_model.reduce(x, batch,
num\_systems=...)}. \textbf{Checked directly against
\texttt{torchmdnet/priors/zbl.py}'s actual source rather than assumed:
\texttt{ZBL} (and therefore \texttt{molecular\_zbl.py}'s
\texttt{MolecularZBL}) overrides only \texttt{post\_reduce}, not
\texttt{pre\_reduce}} --- its repulsive correction is computed from its own
independent edge/distance calculation and added directly to the
already-reduced scalar $y$, never touching the per-atom tensor $x$ at all.
This means a capture prior hooking \texttt{pre\_reduce} sees the
\emph{bare} GNN atomic energy decomposition, with no ZBL contribution,
regardless of where in \texttt{prior\_model}'s list it is placed --- not
the ``ZBL-inclusive'' tensor an earlier draft of this scope incorrectly
assumed. This is treated here as the correct behavior, not a gap to work
around: Raja et al.'s own local-energy formula
($U_\theta(\gamma)=\sum_i U_\theta(\gamma^{(i)})$) assumes a clean
per-atom decomposition, which is exactly what the bare GNN output is and
what ZBL's own pairwise, inter-atomic correction structurally is not; and
\texttt{MolecularZBL} specifically excludes same-molecule pairs from its
correction by design (Section~\ref{sec:q4-bonded-exclusion}), so it
contributes exactly zero to any single-molecule local neighborhood's true
energy in the first place --- there is nothing lost for the one-molecule
neighborhood this scope actually uses. \texttt{L\_QM} is unaffected by any
of this and continues to regularize against the model's full forward pass
(GNN + ZBL together), exactly as before --- only the local-observable
capture is GNN-only, a deliberate, documented simplification rather than
an oversight.

Critically, \texttt{torchmdnet/priors/base.py}'s \texttt{BasePrior.pre\_reduce}
default implementation is a pure identity (\texttt{return x}), confirming
that a new, side-effect-only prior class can be added to
\texttt{model.prior\_model} (a mutable list, already iterated exactly this
way) purely to \emph{capture} a reference to \texttt{x} before returning it
unchanged --- no monkey-patching of \texttt{forward()} itself needed, and
zero change to the model's real output (energy, forces) for every other
purpose; list position does not matter for this capture specifically,
since nothing else in this project's own \texttt{prior\_model} touches
\texttt{x} at the \texttt{pre\_reduce} stage. Per-molecule grouping then
reuses \texttt{molecule\_group\_ids}
(already established, from \texttt{diagnose\_short\_range\_collapse.py})
to scatter-sum the captured per-atom tensor by molecule membership --- the
same grouped-sum pattern \texttt{physics\_losses.py}'s
\texttt{per\_fragment\_momentum\_loss} already uses and this project
already trusts, just for energy instead of force, and can reuse
\texttt{torchmdnet.models.utils.scatter} directly (already imported and
used by \texttt{torchmdnet/priors/zbl.py} itself, so it is not a new,
separately-trusted dependency).

\textbf{A derived efficiency argument (this project's own reasoning
combining the two verified sources above, not asserted by either source
directly): this should also structurally fix the chronic
sample-starvation problem that motivated \texttt{max\_pooled\_samples} in
the first place.} The dominant per-snapshot memory cost identified in
Section~\ref{sec:q4-stable-oom} is the model's internal force computation
(\texttt{dy = grad([y], [pos], create\_graph=self.training, ...)}, line 620
of \texttt{model.py}) --- unavoidable because this model instance always
runs with \texttt{derivative=True}, independent of whether forces are
actually needed for a given call. That cost is paid \emph{once per
forward pass}, before any of our own downstream processing of \texttt{x}.
Extracting one summed scalar (the current whole-box approach) versus up to
64 molecule-level sums (the local approach) from the \emph{same},
already-computed, already-paid-for per-atom tensor costs only a cheap
scatter-sum on top --- negligible next to the force computation already
incurred either way. In other words: the same per-snapshot memory budget
that currently yields exactly 1 usable sample (capped further down to
\texttt{max\_pooled\_samples}=5 for memory safety) should, under the local
estimator, yield up to $\sim$64 samples per snapshot at no additional
memory cost --- directly addressing why this project's own recipe has
needed far more replicas/iterations than Raja et al.'s water recipe
(R=8, a single cycle) to reach comparable statistical power.

\textbf{Concrete file-by-file scope:}
\begin{enumerate}
  \item \textbf{New: \texttt{src/local\_molecule\_observable.py}} --- the
  capture-prior class (subclassing \texttt{BasePrior}-equivalent behavior,
  storing a reference to \texttt{x} and returning it unchanged); a function
  computing $g(\gamma)$ (mean O--H bond length) per molecule, vectorized
  across all molecules in a configuration given \texttt{molecule\_group\_ids};
  a function computing $U_\theta(\gamma)$ per molecule via the scatter-sum
  described above; a function sampling the scalar $g_{\text{target}}$ from
  reference DFT configurations.
  \item \textbf{Modify \texttt{train\_waterbox\_stable.py}}: replace the
  \texttt{\_stacked\_rdf}/whole-box-\texttt{\_config\_energy} pooling loop
  with, per collected snapshot: one forward pass (unchanged cost) capturing
  per-atom energies via the new prior, then extraction of all
  molecule-level $(g,U)$ pairs from that single snapshot.
  \texttt{max\_pooled\_samples} is reinterpreted as a cap on
  \emph{snapshots forwarded} (the real memory-relevant quantity), not on
  raw local samples --- a new, separate cap (or none at all, if memory
  measurement shows no need) governs how many of the resulting
  $\sim$64-per-snapshot local samples are pooled into one gradient step.
  \item \textbf{No changes needed}: \texttt{boltzmann\_estimator.py}
  (already generic), \texttt{waterbox\_langevin.py} (rewind criterion
  unchanged), \texttt{evaluate\_waterbox.py}/\texttt{waterbox\_ase.py}
  (the capture prior is attached transiently, in-memory, only inside
  \texttt{fine\_tune\_stable}'s own model-loading step --- it changes no
  model output, so it does not need to be part of the saved checkpoint or
  re-registered at every other reload site the way \texttt{MolecularZBL}
  does).
\end{enumerate}

\textbf{What remains unverified until run on the training box (same
discipline as every other water-box script that touches torchmdnet
internals --- \texttt{molecular\_zbl.py}, \texttt{waterbox\_data.py}'s
Bohr/Hartree check, \texttt{verify\_zbl\_units.py} --- none of this can be
executed on this checkout):}
\begin{itemize}
  \item The exact attribute path from the loaded \texttt{LNNP}/\texttt{WaterLNNP}
  wrapper down to the real \texttt{TorchMD\_Net} instance's
  \texttt{.prior\_model} list, and whether it is a plain list or an
  \texttt{nn.ModuleList} (both support \texttt{.append()}, but the exact
  type affects whether the appended prior's own (nonexistent) parameters
  need registering).
  \item Whether \texttt{x}'s shape/atom-ordering immediately after
  \texttt{output\_model.pre\_reduce} aligns exactly with
  \texttt{z}/\texttt{pos}/\texttt{molecule\_group\_ids}'s own atom
  ordering (expected, since both derive from the same input tensors in the
  same forward call, but not directly confirmed against a running model).
  \item Whether appending a prior after \texttt{MolecularZBL} at
  post-load time (rather than at original construction time, before the
  checkpoint was saved) behaves identically to a prior list built at
  construction --- expected, since \texttt{BasePrior.pre\_reduce}'s
  identity default has no construction-order dependence, but unconfirmed
  against a real checkpoint.
\end{itemize}

\textbf{Effort/risk assessment}: a real, multi-file implementation (new
module, a rework of the pooling loop's inner structure), not a parameter
change --- larger than anything attempted so far in this StABlE line of
work, but every design decision above traces to a specific, quoted primary
source rather than an assumption, and the mechanism proposed
(side-effect-only prior, reusing already-trusted grouping/scatter
utilities) is deliberately the least invasive path found, avoiding any
monkey-patch of \texttt{TorchMD\_Net.forward} itself. A smoke test
(confirm the capture prior fires, confirm per-molecule sums match a
hand-computed reference on one held-out configuration, confirm
$g(\gamma)$ values fall in a physically sensible O--H bond-length range)
should precede any real fine-tuning run, matching this project's own
established practice for every new capability in this codebase.

\textbf{Implementation complete, one correction made while writing it.}
Split across three files, following this project's own established
convention (\texttt{diagnose\_short\_range\_collapse.py},
\texttt{waterbox\_langevin.py}) of separating pure logic from
torchmdnet-dependent code so the former stays locally testable:
\texttt{src/local\_molecule\_geometry.py} (pure numpy/torch --- the
per-molecule energy scatter-sum and mean-O--H-bond-length computation),
\texttt{src/local\_molecule\_observable.py} (\texttt{AtomicEnergyCapture},
\texttt{attach\_atomic\_energy\_capture}, and the reference-sampling
function --- needs torchmdnet, cannot be imported on this checkout, same
caveat as \texttt{molecular\_zbl.py}), and
\texttt{src/verify\_local\_molecule\_geometry.py} (a new synthetic test,
5/5 checks passing on this checkout: known-ground-truth bond lengths,
periodic-boundary correctness, gradient-routing through
\texttt{index\_add\_}, a degenerate-molecule skip case, and a mixed
valid/degenerate case confirming the skip realigns by molecule id rather
than silently index-shifting - see the first-smoke-test correction below
for why the last two exist).
\texttt{train\_waterbox\_stable.py} gained a \texttt{local\_observable}
parameter (\texttt{--local-observable} on the CLI) --- default \texttt{False}
preserves the original whole-box-RDF behavior exactly, unchanged, so every
existing result stays reproducible; \texttt{max\_pooled\_samples} keeps its
existing meaning (a cap on snapshots forwarded through a model call, not on
raw pooled samples) unchanged in either mode.

While writing \texttt{local\_molecule\_observable.py}, a claim already
written into this scope turned out to be wrong and was corrected on the
spot rather than left standing: an earlier pass here asserted that
appending the capture prior \emph{after} \texttt{MolecularZBL} in
\texttt{prior\_model}'s list would capture the ZBL-inclusive atomic energy.
Checked directly against \texttt{torchmdnet/priors/zbl.py}'s actual source
while implementing this: \texttt{ZBL} overrides only \texttt{post\_reduce},
never \texttt{pre\_reduce}, so it never touches the per-atom tensor at all
--- list position was irrelevant, and the captured tensor is unconditionally
GNN-only. Both this section and \texttt{CLAUDE.md} were corrected in place
(not left as a silent discrepancy between the plan and the code) once this
was found; see the corrected paragraph above under ``How this is
implemented mechanically'' for the fixed reasoning.

\textbf{First smoke-test attempt: crashed, on a real, informative edge case
--- not one of the previously-flagged unverified assumptions.} A random
real DFT reference configuration produced a molecule with 0 O and 1 H atoms
during reference-target sampling: \texttt{structural\_metrics.infer\_bonds}
(shared, distance-threshold bond inference, used throughout this project)
occasionally misses a real but thermally-stretched O--H bond, splitting one
water molecule into two degenerate groups.
\texttt{evaluate\_waterbox.py}'s own \texttt{summarize\_molecule\_groups}
already anticipated exactly this possibility (``asserting every water
molecule works out to exactly \{O: 1, H: 2\}, not some other split caused
by a missed/spurious bond'') without it ever having been strictly enforced
before this function existed. The first version of
\texttt{per\_molecule\_mean\_oh\_bond\_length} raised \texttt{ValueError}
on this, too strict given $\sim$12{,}800 molecules get sampled for a
typical 200-config reference average and a rare miss is a negligible
approximation, not a bug to guard against by crashing. \textbf{Fixed}: skip
a degenerate molecule instead of raising, returning
\texttt{(values, valid\_group\_ids)} rather than a plain array so the
skip's effect on molecule identity is explicit, not implicit. This
mattered beyond the reference-sampling call site:
\texttt{local\_molecule\_energies} computes one energy per molecule
regardless of composition (a pure scatter-sum does not care about
elemental identity), so \texttt{train\_waterbox\_stable.py}'s own pooling
loop was updated to index the energy tensor by the SAME
\texttt{valid\_group\_ids} before pairing it with the (now possibly
shorter) observable array --- without this, a skip would have silently
misaligned $g$ and $U$ by molecule identity rather than simply dropping one
degenerate sample cleanly. Two new synthetic checks
(\texttt{check\_malformed\_molecule\_skipped},
\texttt{check\_mixed\_valid\_and\_degenerate\_molecules}) cover this
directly, the latter specifically placing the degenerate molecule between
two valid ones in storage order to catch exactly the index-shift bug this
fix needed to avoid.

\textbf{Re-run smoke test: passed cleanly, and directly confirmed the core
efficiency claim.} 4 forwarded snapshots yielded 254 pooled local samples
--- 63.5/snapshot, essentially the theoretical ceiling of 64 (only 2
molecules skipped as degenerate across all 4 configurations, matching the
expected rare-miss rate). $\mathcal{L}_{\text{obs}}$ values look tiny
($\sim$0.00015) purely because the observable's units changed (a squared
difference between two $\sim$1\,\AA{} bond lengths, vs.\ the old
dimensionless RDF-density mismatch) --- not comparable in magnitude to any
prior run's logged value. One replica tripped every single stability check
across all 3 iterations while the other two stayed perfectly stable
(0 rewinds each) --- consistent with the already-established $\sim$1-in-3
good-replica pattern, though $n=3$ replicas is too small to read further
into.

\textbf{Real per-snapshot memory cost, measured directly from this run's
own \texttt{[mem]} log lines (9.11/13.56/17.98\,GB at snapshots 2/3/4):
$4.44$\,GB/snapshot marginal cost, $\sim$0.24\,GB baseline} --- confirming
almost the entire cost is per-snapshot (the model's internal
\texttt{create\_graph=True} force computation, paid once per forward pass
regardless of how many local samples get extracted afterward), essentially
zero fixed overhead. Raja et al.'s own guidance (``the minibatch size $B$
should... be chosen as large as possible... while remaining within GPU
memory limits,'' Appendix A.6) can now be acted on with real numbers
instead of a guess: $\texttt{max\_pooled\_samples}=8$ extrapolates to
$\sim$35.7\,GB used, $\sim$11.8\,GB headroom (of 47.5\,GB) --- a real,
measured, comfortably-safe doubling of the previous default of 5, not
pushed further toward 9--10 (7.3\,GB/2.9\,GB headroom respectively),
since both would extrapolate beyond the actually-measured 2--4-snapshot
range.

\textbf{Next run (not yet executed): the first real, non-smoke local-observable
fine-tuning attempt}, changing only the observable type and the
now-measurement-justified \texttt{max\_pooled\_samples} relative to the
last validated global-observable configuration
(Section~\ref{sec:q4-stable-lr-bisect2}) --- $\alpha=3\times10^{-6}$,
$\lambda_{QM}=1.0$, $n\_replicas=8$, $n\_outer\_iterations=20$, and the
already-established $dt=0.1$/$\texttt{learn\_window\_steps}=10{,}000$/
$\texttt{check\_interval}=100$ --- held fixed so that any difference in
outcome is attributable to the observable switch, not a confound with
simultaneously re-tuning the learning rate again from scratch.

\subsubsection{First real local-observable run: safe, but a genuinely mixed result}
\label{sec:q4-stable-local-v1}

\texttt{max\_pooled\_samples=8} performed exactly as extrapolated ---
$n\_samples$ sits at 511--512 almost every iteration (occasionally one or
two lower from a skipped degenerate molecule), confirming the measured
per-snapshot cost held at this scale. $\mathcal{L}_{\text{QM}}$ stayed
safely bounded throughout (min 5.8, max 122.0, no runaway trend) ---
$\alpha=3\times10^{-6}$ remains non-destructive under the local observable,
carrying over safely from the global-observable bisection.

\textbf{$\mathcal{L}_{\text{obs}}$ itself shows essentially no resolvable
trend, and this time that is stronger evidence of a real plateau, not just
noise}: Welch $t=-0.18$ between first-10 and last-10 iteration means, a
full linear fit $R^2=0.013$, and $R^2=0.109$ even after excluding
iteration 0's apparent transient --- despite $\sim$100$\times$ more samples
per iteration than any previous run, ruling out ``too few samples per
gradient step'' as the explanation this specific time.

\textbf{The rewind/stability dynamics tell a different, genuinely mixed
story than the loss curve alone would suggest.} Population-level trip rate
across all 8 replicas shows only a small, non-significant decline
(first-10 mean 58.9\% $\to$ last-10 mean 57.1\%, $t=-1.23$) --- no broad
improvement. But replica 3 shows a real, statistically clear stability
improvement: first-10 mean trip rate 91.4\% $\to$ last-10 mean 76.2\%,
$t=-3.22$, with a visibly abrupt regime change partway through (chronically
79--98\% through iteration 13, then 61--76\% for the remainder) ---
the same kind of single-replica regime-change phenomenon already
documented in Section~\ref{sec:q4-stable-pilot1}, here running in the
\emph{stabilizing} direction rather than the destabilizing one. Replica 6
shows no real change either way ($t=-0.03$).

\textbf{Read together: real parameter movement (the bounded but nonzero
$\mathcal{L}_{\text{QM}}$ fluctuation), no resolvable movement in the
specific first-moment observable being matched, and a narrow but real
stability improvement concentrated in one replica rather than the
population.} A plausible explanation, not yet confirmed: matching only the
ensemble \emph{mean} of a per-molecule observable does not directly
constrain which specific configurations drive instability
(Section~\ref{sec:q1-rollout-update}'s own force-error mechanism is a
per-configuration, not population-average, property) --- the model's mean
O--H bond length may already be close to correct on average even while
specific problem molecules/configurations remain undercorrected, so a
mean-matching loss could plateau immediately while still producing real,
if narrow, downstream stability changes for individual replicas. This is
exactly the kind of result the training loss alone cannot resolve; the
direct target checks (\texttt{evaluate\_waterbox.py --per-sample},
\texttt{compare\_stable\_rollout.py} against the ext70 baseline) on this
run's \texttt{stable\_final.ckpt} are the next step.

\textbf{Static accuracy check: a larger energy-accuracy cost than any
prior run at this same nominal learning rate.} Compared against this exact
checkpoint's own pre-fine-tuning baseline: energy MAE
$2.451\to4.928$\,eV ($+101\%$), force MAE $0.798\to0.744$\,eV/\AA\
($-6.8\%$), mean momentum violation $62.65\to59.67$ ($-4.8\%$), max
violation $2086.0\to1883.5$ ($-9.7\%$). Force/momentum move favorably by
almost exactly the same margin every prior run has shown, but the energy
regression is notably worse than the identical $\alpha=3\times10^{-6}$,
20-iteration run under the \emph{global} observable produced ($+56\%$,
Section~\ref{sec:q4-stable-lr-final}) --- a real difference attributable to
the observable switch itself at matched hyperparameters, not noise
(worst-case per-sample rows are the same known-anomalous configurations at
the same magnitude as every prior check, not a fresh ZBL-singularity
artifact, so the aggregate mean is trustworthy as reported). Plausible
explanation, consistent with Section~\ref{sec:q4-stable-local-v1}'s own
finding that this run produced a real (replica 3), not just noisy, stability
change despite an unmoving loss: the local observable's much larger,
lower-variance per-iteration sample pool ($n\approx512$ vs.\ $n\approx5$--9)
likely yields a more consistent net gradient direction per step, which can
mean genuinely \emph{larger} real parameter movement per iteration at the
same nominal $\alpha$, not just a lower-variance estimate of the same
movement --- i.e., ``safe under the noisier global estimator'' may not
transfer cleanly to ``equally safe under the local one.''

\textbf{Rollout comparison: another clean null result, and the concerning
case above is now confirmed, not just possible.} Plateau temperature
$771.5\pm14$\,K (after) vs.\ $768.0\pm26$\,K (before) --- a $+3.6$\,K mean
difference, i.e.\ if anything very slightly \emph{hotter}, not cooler.
Paired per-replicate differences do not agree in sign: 3 of 5 velocity
seeds run hotter after fine-tuning, 2 run cooler, spanning $-46.2$ to
$+43.2$\,K (an 89\,K spread) --- by the same unanimous-direction standard
already applied throughout this project, this is noise, not a real effect
in either direction.

\textbf{Taken together with Section~\ref{sec:q4-stable-lr-final}'s replica-3
finding: an internal training-loop stability proxy (the Langevin rewind
rate under the model currently being fine-tuned) did not predict actual
downstream rollout stability from a real DFT starting configuration} ---
a distinct, new instance of the same general lesson that motivated this
entire StABlE line of work in the first place
(Section~\ref{sec:q1-rollout-update}: static accuracy does not predict
rollout stability), now shown to extend to a \emph{different} proxy metric
than the one originally identified.

\textbf{Where this leaves the local-observable line of work}: two
observable formulations (global whole-box RDF, local per-molecule mean
O--H bond length matching Raja et al.'s own validated water recipe) and
three learning rates ($10^{-6}$ confirmed too small; $10^{-5}$ confirmed
damaging under the noisier global estimator; $3\times10^{-6}$ confirmed
non-destructive under both observables) have now all been tried against
this exact checkpoint (\texttt{water\_absolute} seed 1 ext70), and none has
produced a measurable rollout-stability improvement --- the local
observable's only measured effect relative to the global one was a
\emph{larger} static-accuracy cost ($+101\%$ vs.\ $+56\%$ energy MAE) for
the same null downstream result. This is a real, not yet resolved, decision
point for the project, not a conclusion to paper over with another
hyperparameter attempt.

\subsection{Q3 (resolved for this seed pair): why does momentum training run
hotter under NVE?}
\label{sec:q3}

The mechanistic question the rollout results raised. Momentum-conservation
training is designed to reduce spurious net force/torque on individual
molecules --- naively, that should make dynamics \emph{more} physically
consistent, not less. Original candidate explanations, and how each fared:
\begin{itemize}
  \item \textbf{Ruled out}: the momentum loss trades off against force
  accuracy in a way invisible to the static force MAE metric but that
  matters once forces are integrated forward in time. Cross-evaluating both
  models on real rollout frames found no widening point-wise disagreement
  with frame index --- the two models agree closely (1--5\%) at every frame
  tested, early or late.
  \item \textbf{Weakened, then superseded}: an artifact of \emph{which}
  checkpoint got selected. The 10/10-consistent replicated direction
  (Section~\ref{sec:q3-progress}) already made pure checkpoint-selection
  luck unlikely ($p\approx0.002$); the seed-specific reconciliation below
  makes it unnecessary as an explanation --- the effect is real and present
  in this seed's original static evaluation too, not a training-run fluke.
  \item \textbf{Ruled out as originally stated}: momentum conservation
  satisfied by redistributing force \emph{within} a molecule (large
  internal, small net) rather than by genuinely reducing force error. The
  internal-force share was nearly identical between conditions
  ($\sim$97\% both at equilibrium); a rigorous 100-config paired test
  confirmed no such compensating internal-force bias exists.
\end{itemize}

\textbf{Resolved}: a rigorous paired test found the systematic bias itself
--- net per-molecule force is $0.073$\,eV/\AA{} higher for momentum on
average, highly significant ($p\approx4{\times}10^{-125}$), not a general
``momentum's forces are bigger'' effect (raw/internal force show the
opposite, tiny sign) and not compensated by an opposite torque bias (net
torque: no significant mean shift, $p=0.85$). Checking this seed pair's own
\emph{static} evaluation resolved the apparent tension with
Section~\ref{sec:n6-update}'s aggregate result: seed 0 already showed
momentum with \emph{higher} per-fragment momentum violation than absolute
($45.16$ vs.\ $41.93$) before any rollout was ever run --- the $n{=}6$
aggregate favoring momentum is dominated by \texttt{water\_absolute} seed
2's outlier, and the six seeds individually split roughly 3-and-3. Seed 0,
used for every rollout in this investigation, happens to be one of the three
where momentum is worse. The rollout instability is not a hidden dynamical
effect decoupled from the static metrics --- it is the direct consequence of
something the static per-fragment evaluation already showed for this exact
seed pair, compounded over $\sim$900\,fs of nonlinear dynamics into a large,
obvious temperature difference. See Section~\ref{sec:q3-progress} for the
full derivation.

\textbf{Generalization, checked}: does this track seed-to-seed? Repeating
both replicate batches on seed 5 (where the static metric favors momentum,
unlike seed 0) found the rollout-stability gap shrinks substantially and its
drift comparison stops being directionally consistent, though a weak
temperature-side lean toward momentum remains --- see
Section~\ref{sec:q3-seed5} for the full comparison. Seed-to-seed variation
in whether the momentum loss achieves its intended effect, not a fixed
property of the method, was the best-supported explanation for the
rollout-stability question at that point in the investigation. This
conclusion was reached, however, entirely on checkpoints later found to
carry two independent confounds (stock ZBL's bonded-pair double-counting,
Section~\ref{sec:q4-negative-result}; unfair checkpoint selection,
Section~\ref{sec:q4-anneal-confound}) --- Section~\ref{sec:q1-rollout-update}
revisits this question on checkpoints with both fixed.

\subsection{Q1 revisited: momentum vs.\ absolute under the corrected methodology}
\label{sec:q1-rollout-update}

With both the ZBL bonded-exclusion fix and the anneal/checkpoint-selection
fix in place (Sections~\ref{sec:q4-bonded-exclusion},
\ref{sec:q4-anneal-confound}), the rollout comparison was repeated at two
seeds, both axes:

\begin{center}
\begin{tabular}{llll}
\toprule
seed & axis & \texttt{water\_absolute} & \texttt{water\_absolute+momentum} \\
\midrule
0 & velocity & 1682$\pm$89\,K & 1072$\pm$22\,K \\
0 & config   & 1122$\pm$260\,K & 828$\pm$170\,K \\
1 & velocity & 772$\pm$15\,K & 706$\pm$18\,K \\
1 & config   & 644$\pm$150\,K & 612$\pm$130\,K \\
\bottomrule
\end{tabular}
\end{center}

Momentum wins --- is more stable --- in all four seed$\times$axis
combinations, a materially stronger and more consistent signal than
anything obtained before these confounds were fixed. This directly
reverses the original seed-0, no-ZBL finding (Section~\ref{sec:q3-progress}:
momentum ran \emph{hotter} in 10/10 replicates on that same seed), which is
now understood to have been measured on checkpoints without either fix
applied.

Two further observations from this same data. \textbf{Seed 0 is still far
hotter than seed 1 in absolute terms, for both conditions} ---
roughly $2\times$ across every axis --- confirming the seed-to-seed
variance already established elsewhere in this project is not an artifact
of the bugs just fixed. Checking this against the Section~\ref{sec:q4-epot-gap}
mechanism: seed 0's velocity-axis \texttt{epot} gap is 34.2\,eV
(\texttt{water\_absolute}) / 19.1\,eV (\texttt{water\_absolute+momentum}) ---
roughly 3$\times$ and 2$\times$ seed 1's 11.6/9.96\,eV respectively ---
cleanly explaining why seed 0 runs so much hotter overall, and confirming
this isn't a residual checkpoint-selection artifact: seed 0's selected
epoch (68, identical to seed 1) already passed the same eligibility and
batch-of-1 safety checks. Seed 0 simply converged to a checkpoint with a
genuinely larger force/energy residual near its own training-distribution
configurations than seed 1 did, despite statistically indistinguishable
aggregate test MAE between the two seeds (2.5--3.6\,eV vs.\ 2.2--2.5\,eV) ---
yet another instance of aggregate static accuracy hiding what actually
matters for rollout stability.

\textbf{Momentum's margin over absolute scales with how bad the underlying
problem is}: roughly 610\,K (36\% relative) at seed 0's velocity axis versus
66\,K (9\% relative) at seed 1's, a similar pattern on the config axis. A
plausible, testable hypothesis rather than a conclusion at $n=2$ seeds ---
the per-molecule momentum constraint may act as a mild stabilizing
regularizer that matters more when the underlying force-accuracy problem is
worse --- worth checking once StABlE-fine-tuned checkpoints
(Section~\ref{sec:q4-stable-path}) are available for both conditions at
both seeds, since that would either shrink both conditions' gaps toward
each other (consistent with the hypothesis, momentum's edge being most
useful exactly where the base force error is largest) or leave the current
margin roughly intact (consistent with momentum and the force-accuracy fix
being unrelated, additive effects).

\textbf{Verdict on Q1 as it stands}: still not the $n=6$ standard this
project otherwise holds itself to for the rollout comparison (only 2 of the
6 seeds with static-accuracy checkpoints have been carried through the full
corrected-methodology rollout pipeline), but on the best-controlled evidence
collected so far --- two independent seeds, two independent starting-geometry
axes, both known confounds removed --- momentum-conservation training is
associated with meaningfully more stable NVE dynamics, unanimously and by a
non-trivial margin. This reverses, rather than merely qualifies, the
project's earlier working conclusion (Section~\ref{sec:q3}) that
seed-to-seed variance alone explained the original hotter-with-momentum
result.

\appendix

\section{Implementation notes (debugging log, condensed)}
\label{app:impl}

Kept here for reproducibility, not because it's scientifically interesting
on its own. The water-box pipeline had never been run against a real
\texttt{torchmdnet} install before this round of work; four real bugs
surfaced in sequence, each one only visible after fixing the last:

\begin{enumerate}
  \item \textbf{Dead download URL}: \texttt{torchmdnet}'s shipped
  \texttt{WaterBox} downloader points at a Materials Cloud API endpoint
  (\texttt{record\_id=} query param) retired when Materials Cloud migrated
  to short-ID URLs. Fixed by overriding the URL to the same record's current
  address rather than patching the installed package.
  \item \textbf{Wrong physical units}: the raw \texttt{dataset\_1593.xyz} is
  in Bohr/Hartree atomic units despite a quippy-style header that normally
  implies \AA/eV --- confirmed via physical priors (box density, O--H bond
  length, force magnitude), not any stated documentation. Fixed via a
  per-sample unit-conversion transform.
  \item \textbf{Non-periodic bond inference}: the bond-inference utility
  (shared with the aspirin study) used plain Euclidean distance, which
  misgroups molecules whose atoms straddle a periodic box face. Fixed by
  adding minimum-image distance handling, gated behind an optional
  \texttt{box} argument (default \texttt{None} preserves exact old behavior
  for MD17).
  \item \textbf{GPU memory}: TensorNet's per-edge tensor representation,
  combined with this system's higher local connectivity (192-atom periodic
  liquid vs.\ a 21-atom sparse organic molecule), made the aspirin study's
  batch size far too large. Resolved via a smaller real batch size plus
  gradient accumulation to preserve the same effective optimizer batch size.
  \item \textbf{Lightning's \texttt{save\_last=True} isn't what it sounds
  like}: intended as a free, zero-training-cost way to get a
  non-cherry-picked final-epoch checkpoint for comparison against the
  min-validation-loss ``best'' one. Actually only updates ``whenever a
  checkpoint file gets saved,'' which under \texttt{save\_top\_k=1} means
  only on improvement --- so it silently duplicates ``best'' whenever
  nothing later ever beats it, discovered only after re-evaluating all six
  ``last'' checkpoints and finding they matched ``best'' almost exactly (a
  result that initially looked like a script bug). See
  Section~\ref{sec:checkpoint-fix} for the fix.
  \item \textbf{The weight-annealing schedule (Section~\ref{sec:anneal})
  broke checkpoint selection}: monitoring the literal training loss
  (\texttt{val\_total\_mse\_loss}) is unsafe once its own weights change
  mid-run, since the metric's scale shifts at exactly the epoch being
  tested, structurally favoring pre-anneal epochs regardless of real model
  quality. See Section~\ref{sec:checkpoint-fix} for the fix and
  Section~\ref{sec:q1} for confirmation it worked (5/6 cells now select
  post-anneal).
\end{enumerate}

\bibliographystyle{plainnat}
\begin{thebibliography}{9}

\bibitem{md17}
S.~Chmiela, A.~Tkatchenko, H.~E.~Sauceda, I.~Poltavsky, K.~T.~Sch{\"u}tt,
and K.-R.~M{\"u}ller.
\newblock Machine learning of accurate energy-conserving molecular force
fields.
\newblock \emph{Science Advances}, 3(5):e1603015, 2017.

\bibitem{cheng2019water}
B.~Cheng, E.~A.~Engel, J.~Behler, C.~Dellago, and M.~Ceriotti.
\newblock Ab initio thermodynamics of liquid and solid water.
\newblock \emph{Proceedings of the National Academy of Sciences},
116(4):1110--1115, 2019.
\newblock Dataset: Materials Cloud Archive,
DOI: 10.24435/materialscloud:2018.0020/v1.

\bibitem{tensornet}
G.~Simeon and G.~De Fabritiis.
\newblock TensorNet: Cartesian tensor representations for efficient learning
of molecular potentials.
\newblock In \emph{Advances in Neural Information Processing Systems
(NeurIPS)}, 2023.
\newblock arXiv:2306.06482.

\bibitem{torchmdnet2}
R.~P.~Pel{\'a}ez, G.~Simeon, R.~Galvelis, A.~Mirarchi, P.~Eastman,
S.~Doerr, P.~Th{\"o}lke, T.~E.~Markland, and G.~De Fabritiis.
\newblock TorchMD-Net 2.0: Fast neural network potentials for molecular
simulations.
\newblock \emph{Journal of Chemical Theory and Computation},
20(10):4076--4087, 2024.

\bibitem{mace}
I.~Batatia, D.~P.~Kov{\'a}cs, G.~N.~C.~Simm, C.~Ortner, and G.~Cs{\'a}nyi.
\newblock MACE: Higher order equivariant message passing neural networks for
fast and accurate force fields.
\newblock In \emph{Advances in Neural Information Processing Systems
(NeurIPS)}, 2022.
\newblock arXiv:2206.07697.

\bibitem{mlip-validation}
J.~D.~Morrow, J.~L.~A.~Gardner, and V.~L.~Deringer.
\newblock How to validate machine-learned interatomic potentials.
\newblock \emph{Journal of Chemical Physics}, 158:121501, 2023.
\newblock arXiv:2211.12484.

\bibitem{mlip-error-metrics}
Y.~Liu, X.~He, and Y.~Mo.
\newblock Discrepancies and error evaluation metrics for machine learning
interatomic potentials.
\newblock \emph{npj Computational Materials}, 9:174, 2023.

\bibitem{nequip}
S.~Batzner, A.~Musaelian, L.~Sun, M.~Geiger, J.~P.~Mailoa, M.~Kornbluth,
N.~Molinari, T.~E.~Smidt, and B.~Kozinsky.
\newblock E(3)-equivariant graph neural networks for data-efficient and
accurate interatomic potentials.
\newblock \emph{Nature Communications}, 13:2453, 2022.

\bibitem{adaptive-weighting}
D.~Ocampo, D.~Posso, R.~Namakian, and W.~Gao.
\newblock Adaptive loss weighting for machine learning interatomic
potentials.
\newblock \emph{Computational Materials Science}, 2024.
\newblock arXiv:2403.18122.

\bibitem{ase}
A.~H.~Larsen, J.~J.~Mortensen, J.~Blomqvist, I.~E.~Castelli, R.~Christensen,
M.~Du{\l}ak, J.~Friis, M.~N.~Groves, B.~Hammer, C.~Hargus, E.~D.~Hermes,
P.~C.~Jennings, P.~B.~Jensen, J.~Kermode, J.~R.~Kitchin, E.~L.~Kolsbjerg,
J.~Kubal, K.~Kaasbjerg, S.~Lysgaard, J.~B.~Maronsson, T.~Maxson,
T.~Olsen, L.~Pastewka, A.~Peterson, C.~Rostgaard, J.~Schi{\o}tz, O.~Sch{\"u}tt,
M.~Strange, K.~S.~Thygesen, T.~Vegge, L.~Vilhelmsen, M.~Walter, Z.~Zeng, and
K.~W.~Jacobsen.
\newblock The atomic simulation environment --- a Python library for working
with atoms.
\newblock \emph{Journal of Physics: Condensed Matter}, 29(27):273002, 2017.

\bibitem{fu2022forces}
X.~Fu, Z.~Wu, W.~Wang, T.~Xie, S.~Keten, R.~Gomez-Bombarelli, and T.~Jaakkola.
\newblock Forces are not enough: Benchmark and critical evaluation for
machine learning force fields with molecular simulations.
\newblock arXiv:2210.07237, 2022.

\bibitem{stocker2022robust}
S.~Stocker, J.~Gasteiger, F.~Becker, S.~G{\"u}nnemann, and J.~T.~Margraf.
\newblock How robust are modern graph neural network potentials in long and
hot molecular dynamics simulations?
\newblock \emph{Machine Learning: Science and Technology}, 3:045010, 2022.
\newblock \url{https://doi.org/10.1088/2632-2153/ac9955}.

\bibitem{raja2025stable}
S.~Raja, I.~Amin, F.~Pedregosa, and A.~S.~Krishnapriyan.
\newblock Stability-aware training of machine learning force fields with
differentiable Boltzmann estimators.
\newblock \emph{Transactions on Machine Learning Research}, 2025.
\newblock arXiv:2402.13984.

\bibitem{thaler2021difftre}
S.~Thaler and J.~Zavadlav.
\newblock Learning neural network potentials from experimental data via
differentiable trajectory reweighting.
\newblock \emph{Nature Communications}, 12:6884, 2021.
\newblock \url{https://doi.org/10.1038/s41467-021-27241-4}.

\bibitem{zwanzig1954}
R.~W.~Zwanzig.
\newblock High-temperature equation of state by a perturbation method. I.
Nonpolar gases.
\newblock \emph{Journal of Chemical Physics}, 22(8):1420--1426, 1954.

\bibitem{zhang2020dpgen}
Y.~Zhang, H.~Wang, W.~Chen, J.~Zeng, L.~Zhang, H.~Wang, and W.~E.
\newblock DP-GEN: A concurrent learning platform for the generation of
reliable deep learning based potential energy models.
\newblock \emph{Computer Physics Communications}, 253:107206, 2020.

\bibitem{zbl1985}
J.~F.~Ziegler, J.~P.~Biersack, and U.~Littmark.
\newblock \emph{The Stopping and Range of Ions in Solids}.
\newblock Pergamon Press, 1985.

\bibitem{shortrangeprior2025}
Z.~Yan, Z.~Fan, and Y.~Zhu.
\newblock Improving robustness and training efficiency of machine-learned
potentials by incorporating short-range empirical potentials.
\newblock \emph{Journal of Chemical Information and Modeling}, 2025.
\newblock arXiv:2504.15925.

\bibitem{nequip-code}
mir-group.
\newblock NequIP: E(3)-equivariant interatomic potentials (software
repository).
\newblock \url{https://github.com/mir-group/nequip}, \texttt{nequip/nn/
pair\_potential.py}, accessed 2026-08-18.

\bibitem{kingmaba2015adam}
D.~P.~Kingma and J.~Ba.
\newblock Adam: A method for stochastic optimization.
\newblock In \emph{International Conference on Learning Representations
(ICLR)}, 2015.
\newblock arXiv:1412.6980.

\bibitem{mace-finetuning-docs}
I.~Batatia, D.~Kovacs, G.~Simm, and contributors.
\newblock Fine-tuning foundation models; multihead replay fine-tuning
(MACE documentation).
\newblock \url{https://mace-docs.readthedocs.io/en/latest/guide/finetuning.html},
\url{https://mace-docs.readthedocs.io/en/latest/guide/multihead_finetuning.html},
accessed 2026-08-28.

\bibitem{torchmdnet-code}
G.~Simeon, R.~Pinsler, and contributors.
\newblock TorchMD-Net (software repository), \texttt{torchmdnet/module.py}
\texttt{LNNP.configure\_optimizers}.
\newblock \url{https://github.com/torchmd/torchmd-net}, accessed 2026-08-28.

\end{thebibliography}

\end{document}
